\section{辐射}

\subsection{局域辐射源的频谱分析}

前文已从推迟势出发，得到局域电荷体系在长波近似下的远区多极场，以及能量和角动量的瞬时辐射率。这些时域公式给出的是某一时刻源的辐射强度；频谱分析并不引入新的动力学假设，而是把同一组多极矩按时间尺度分解，从而回答辐射能量分布在哪些频率、是否形成连续谱或离散线谱，以及不同频段由轨道的哪些部分控制。因此，具体问题中的势能只负责确定轨道和多极矩，频域中的辐射学部分可以先在一般形式下完成。

以下对非周期过程采用 Fourier 变换约定
\begin{equation*}
	\widetilde f(\omega)=\int_{-\infty}^{\infty}f(t)e^{i\omega t}\d t,
	\qquad
	f(t)=\frac{1}{2\pi}\int_{-\infty}^{\infty}\widetilde f(\omega)e^{-i\omega t}\d\omega.
\end{equation*}
对实函数 $f(t)$ 取共轭即得 $\widetilde f(-\omega)=\widetilde f(\omega)^*$。把 $A(t)$、$B(t)$ 的反变换代入积分并先对时间积分，利用 $\int e^{-i(\omega+\omega')t}\d t=2\pi\delta(\omega+\omega')$，再把积分区间拆成正负频率并两次使用实函数的共轭性质，便得到正频率形式的 Parseval 恒等式
\begin{equation*}
	&\int_{-\infty}^{\infty} A(t):B(t)\d t
	=\frac{1}{2\pi}\int_{-\infty}^{\infty}\widetilde A(\omega):\widetilde B(-\omega)\d\omega\\
	&\quad=\frac{1}{\pi}\int_0^\infty
	\operatorname{Re}\!\left[\widetilde A(\omega):\widetilde B(\omega)^*\right]\d\omega,
\end{equation*}
其中对矢量时冒号理解为普通内积。特别地，分部积分 $n$ 次（非周期过程的相应边界项消失，或可由绝热收敛因子定义）得频率微分规则
\begin{equation*}
	\widetilde{f^{(n)}}(\omega)=(-i\omega)^n\widetilde f(\omega).
\end{equation*}

记观测方向为 $\vb n$，只保留远场中按 $R^{-1}$ 衰减的部分。沿用前文的电偶极矩 $\vb p$、磁偶极矩 $\vb m$ 和无迹电四极矩 $\vb D$，远区电场可统一写成
\begin{equation*}
	\vb E_{\mathrm{rad}}(R\vb n,t)
	=\frac{1}{4\pi\varepsilon_0c^2R}
	\left\{
	\vb n\times\left(\vb n\times\ddot{\vb p}\right)
	+\frac{1}{c}\vb n\times\ddot{\vb m}
	+\frac{1}{6c}\vb n\times
	\left[\vb n\times\left(\dddot{\vb D}\cdot\vb n\right)\right]
	\right\}_{t-R/c}.
\end{equation*}
相应的辐射磁场为 $\vb B_{\mathrm{rad}}=\vb n\times\vb E_{\mathrm{rad}}/c$。对远区电场做 Fourier 变换，迟滞时间换元 $t=t'+R/c$ 只产生相位因子 $e^{i\omega R/c}$，频域远区场为
\begin{equation*}
	\widetilde{\vb E}_{\mathrm{rad}}(R\vb n,\omega)
	=\frac{e^{i\omega R/c}}{4\pi\varepsilon_0c^2R}
	\boldsymbol{\mathcal A}(\omega,\vb n),
\end{equation*}
式中频域辐射振幅定义为
\begin{equation*}
	&\boldsymbol{\mathcal A}(\omega,\vb n)
	=\vb n\times\left(\vb n\times\widetilde{\ddot{\vb p}}\right)
	+\frac{1}{c}\vb n\times\widetilde{\ddot{\vb m}}\\
	&\quad +\frac{1}{6c}\vb n\times
	\left[\vb n\times\left(\widetilde{\dddot{\vb D}}\cdot\vb n\right)\right].
\end{equation*}
辐射磁场使能量以速率 $|\vb E|^2/\mu_0c$ 沿 $\uv n$ 流出，故单位立体角、单位频率的辐射能量为
\begin{equation*}
	\frac{\d^2E}{\d\omega\d\Omega}
	=\frac{R^2}{\pi\mu_0c}
	\left|\widetilde{\vb E}_{\mathrm{rad}}(R\vb n,\omega)\right|^2
	=\frac{1}{16\pi^3\varepsilon_0c^3}
	\left|\boldsymbol{\mathcal A}(\omega,\vb n)\right|^2,
\end{equation*}
其中第二等号是通量对时间的积分与正频率 Parseval 恒等式的直接结果：迟滞相位在模方中抵消，系数合并时用了 $\mu_0\varepsilon_0c^2=1$。上式中的平方必须在各多极振幅相加之后进行，因此在固定观测方向上一般包含 $E1$--$M1$、$E1$--$E2$ 等干涉项。单独取各阶振幅时，有
\begin{equation*}
	\frac{\d^2E_{E1}}{\d\omega\d\Omega}
	=\frac{1}{16\pi^3\varepsilon_0c^3}
	\left|\widetilde{\ddot{\vb p}}\times\vb n\right|^2,
\end{equation*}
\begin{equation*}
	\frac{\d^2E_{M1}}{\d\omega\d\Omega}
	=\frac{1}{16\pi^3\varepsilon_0c^5}
	\left|\widetilde{\ddot{\vb m}}\times\vb n\right|^2,
\end{equation*}
以及
\begin{equation*}
	\frac{\d^2E_{E2}}{\d\omega\d\Omega}
	=\frac{1}{576\pi^3\varepsilon_0c^5}
	\left|\left(\widetilde{\dddot{\vb D}}\cdot\vb n\right)\times\vb n\right|^2.
\end{equation*}
对观测方向积分时使用
\begin{equation*}
	\int n_i n_j\d\Omega=\frac{4\pi}{3}\delta_{ij},
	\qquad
	\int n_i n_jn_kn_l\d\Omega
	=\frac{4\pi}{15}
	\left(\delta_{ij}\delta_{kl}+\delta_{ik}\delta_{jl}+\delta_{il}\delta_{jk}\right).
\end{equation*}
这两个恒等式由各向同性可直接验证：$\int n_in_j\d\Omega$ 只能正比于 $\delta_{ij}$，与 $\delta_{ij}$ 收缩得 $\int\d\Omega=4\pi=3A$；四指标情形按对称性取 $\delta$ 组合的相同系数，收缩两次即定出 $\frac{4\pi}{15}$。以 $E1$ 为例做收缩，$|\widetilde{\ddot{\vb p}}\times\vb n|^2=\widetilde{\ddot p}_i\widetilde{\ddot p}_j^*(\delta_{ij}-n_in_j)$，得 $\int|\widetilde{\ddot{\vb p}}\times\vb n|^2\d\Omega=\frac{8\pi}{3}\widetilde{\ddot p}_i\widetilde{\ddot p}_i^*$；$M1$ 相同。$E2$ 的收缩为 $|(\vb D''\cdot\vb n)\times\vb n|^2=D''_{ij}D''^*_{kl}n_jn_l(\delta_{ik}-n_in_k)$，式中 $\vb D''=\widetilde{\dddot{\vb D}}$，第二个恒等式给出 $\frac{4\pi}{3}-\frac{8\pi}{15}=\frac{4\pi}{5}$ 倍 $D''_{ij}D''^*_{ij}$——无迹条件使 $\delta_{ij}\delta_{kl}$ 项消失。代入各角分布谱并完成系数，任意局域非相对论源的角积分频谱为
\begin{equation*}
	\frac{\d E_{E1}}{\d\omega}
	=\frac{1}{6\pi^2\varepsilon_0c^3}
	\widetilde{\ddot p}_i\widetilde{\ddot p}_i^*
	=\frac{\omega^4}{6\pi^2\varepsilon_0c^3}
	\widetilde p_i\widetilde p_i^*,
\end{equation*}
\begin{equation*}
	\frac{\d E_{M1}}{\d\omega}
	=\frac{1}{6\pi^2\varepsilon_0c^5}
	\widetilde{\ddot m}_i\widetilde{\ddot m}_i^*
	=\frac{\omega^4}{6\pi^2\varepsilon_0c^5}
	\widetilde m_i\widetilde m_i^*,
\end{equation*}
以及
\begin{equation*}
	\frac{\d E_{E2}}{\d\omega}
	=\frac{1}{720\pi^2\varepsilon_0c^5}
	\widetilde{\dddot D}_{ij}\widetilde{\dddot D}_{ij}^*
	=\frac{\omega^6}{720\pi^2\varepsilon_0c^5}
	\widetilde D_{ij}\widetilde D_{ij}^*.
\end{equation*}
第二个等号是频率微分规则的直接结果。不难看到，$E1$、$M1$ 谱按 $\omega^4$ 增长而 $E2$ 谱按 $\omega^6$ 增长，四极项对高频的权重更高。频谱对正频率积分后恢复前文的时域总辐射能量，亦即
\begin{equation*}
	E_X=\int_0^\infty\frac{\d E_X}{\d\omega}\d\omega
	=\int_{-\infty}^{\infty}P_X(t)\d t,
	\qquad X=E1,M1,E2.
\end{equation*}
这个等式既是频谱归一化的依据，也是具体闭式计算最重要的校验之一。

非束缚运动的多极矩在 $t\to\pm\infty$ 具有不同的渐近状态，故上述角积分频谱给出连续谱。在非相对论长波理论内部，最低频项可在求完整 Fourier 积分之前由分部积分的边界项直接确定：例如 $\widetilde{\ddot{\vb p}}(0)=\int\ddot{\vb p}\d t=[\dot{\vb p}]_{-\infty}^{+\infty}$，即
\begin{equation*}
	\widetilde{\ddot{\vb p}}(0)=\dot{\vb p}(+\infty)-\dot{\vb p}(-\infty),
	\qquad
	\widetilde{\ddot{\vb m}}(0)=\dot{\vb m}(+\infty)-\dot{\vb m}(-\infty),
\end{equation*}
\begin{equation*}
	\widetilde{\dddot{\vb D}}(0)=\ddot{\vb D}(+\infty)-\ddot{\vb D}(-\infty).
\end{equation*}
代入前面的角积分频谱公式取零频极限，得到
\begin{equation*}
	&\left.\frac{\d E_{E1}}{\d\omega}\right|_{\omega\to0}
	=\frac{|\Delta\dot{\vb p}|^2}{6\pi^2\varepsilon_0c^3},\\
	&\left.\frac{\d E_{M1}}{\d\omega}\right|_{\omega\to0}
	=\frac{|\Delta\dot{\vb m}|^2}{6\pi^2\varepsilon_0c^5},\\
	&\left.\frac{\d E_{E2}}{\d\omega}\right|_{\omega\to0}
	=\frac{\Delta\ddot D_{ij}\Delta\ddot D_{ij}}{720\pi^2\varepsilon_0c^5}.
\end{equation*}
式中 $\Delta\dot{\vb p}=\dot{\vb p}(+\infty)-\dot{\vb p}(-\infty)$ 为初末渐近值之差，$\Delta\dot{\vb m}$、$\Delta\ddot{\vb D}$ 含义类似。这些式子表明，非相对论多极谱的零频极限只依赖初末渐近运动。然而软辐射的这一性质并不依赖非相对论或长波展开；它实际上是完整 Lienard--Wiechert 理论中的相对论普适结论。下面从点电荷的精确频谱直接推导这一结果，并由此说明其适用范围和更深的红外结构。

考虑若干条点电荷世界线 $\vb r_a(t)$，记
\begin{equation*}
	\boldsymbol\beta_a=\frac{\dot{\vb r}_a}{c},
	\qquad
	\Phi_a(t,\vb n)=t-\frac{\vb n\cdot\vb r_a(t)}{c}.
\end{equation*}
由前文 Lienard--Wiechert 远区场 $\vb E_{\mathrm{rad},a}(R\vb n,t)=\frac{q_a}{4\pi\varepsilon_0cR}\left[\frac{\vb n\times[(\vb n-\boldsymbol\beta_a)\times\dot{\boldsymbol\beta}_a]}{(1-\vb n\cdot\boldsymbol\beta_a)^3}\right]_{t_s}$ 做 Fourier 变换：把观测时间换为粒子固有坐标时，$\d t=(1-\vb n\cdot\boldsymbol\beta_a)\d t_s$，相位为 $\omega\Phi_a$，迟滞因子 $e^{i\omega R/c}$ 在模方中抵消；再按上面的通量--Parseval 程序组装，完整相对论谱角分布为
\begin{equation*}
	\frac{\d^2E}{\d\omega\d\Omega}
	=\frac{1}{16\pi^3\varepsilon_0c}
	\left|
	\sum_a q_a\int_{-\infty}^{\infty}\d t
	\frac{\vb n\times\left[(\vb n-\boldsymbol\beta_a)\times\dot{\boldsymbol\beta}_a\right]}
	{(1-\vb n\cdot\boldsymbol\beta_a)^2}
	e^{i\omega\Phi_a(t,\vb n)}
	\right|^2.
\end{equation*}
这里的时间是粒子世界线的坐标时，式中已经保留了完整的迟滞相位 $\omega\vb n\cdot\vb r/c$ 和全部相对论束射因子。关键恒等式为
\begin{equation*}
	\frac{\d}{\d t}
	\left[
	\frac{\vb n\times(\vb n\times\boldsymbol\beta_a)}
	{1-\vb n\cdot\boldsymbol\beta_a}
	\right]
	=
	\frac{\vb n\times\left[(\vb n-\boldsymbol\beta_a)\times\dot{\boldsymbol\beta}_a\right]}
	{(1-\vb n\cdot\boldsymbol\beta_a)^2}.
\end{equation*}
它可由三重矢量积展开 $\vb n\times(\vb n\times\boldsymbol\beta)=\vb n(\vb n\cdot\boldsymbol\beta)-\boldsymbol\beta$ 与商的求导直接验证：分子展开为 $(\vb n\cdot\dot{\boldsymbol\beta})\vb n-(\vb n\cdot\dot{\boldsymbol\beta})\boldsymbol\beta-(1-\vb n\cdot\boldsymbol\beta)\dot{\boldsymbol\beta}$，恰为 $\vb n\times[(\vb n-\boldsymbol\beta)\times\dot{\boldsymbol\beta}]$ 的三重矢量积展开。若各粒子在 $t\to\pm\infty$ 趋于匀速，便可在精确谱角分布中先取 $\omega\to0$：此时相位 $e^{i\omega\Phi_a}\to1$，积分即全导数的边界值。将所有渐近外线统一编号为 $A$，约定出射线 $\eta_A=+1$、入射线 $\eta_A=-1$，得到领先软振幅
\begin{equation*}
	\boldsymbol{\mathcal S}^{(0)}(\vb n)
	=\sum_A\eta_Aq_A
	\frac{\vb n\times(\vb n\times\boldsymbol\beta_A)}
	{1-\vb n\cdot\boldsymbol\beta_A}.
\end{equation*}
于是任意相对论散射的领先软谱为
\begin{equation*}
	\left.\frac{\d^2E}{\d\omega\d\Omega}\right|_{\omega\to0}
	=\frac{1}{16\pi^3\varepsilon_0c}
	\left|\boldsymbol{\mathcal S}^{(0)}(\vb n)\right|^2.
\end{equation*}
这个结果只含各带电粒子的初末速度，不含相互作用区内的加速度历史。换言之，求领先软辐射时不需要先求完整的时域辐射波形；动力学只需提供渐近散射映射。

领先软谱还可以完成全立体角积分。物理散射过程满足电荷守恒
\begin{equation*}
	\sum_A\eta_Aq_A=0.
\end{equation*}
展开模方得 $\sum_{A,B}\eta_A\eta_Bq_Aq_B\frac{g(\boldsymbol\beta_A)\cdot g(\boldsymbol\beta_B)}{(1-\vb n\cdot\boldsymbol\beta_A)(1-\vb n\cdot\boldsymbol\beta_B)}$，其中 $g(\boldsymbol\beta)=\vb n\times(\vb n\times\boldsymbol\beta)$，且
\begin{equation*}
	g(\boldsymbol\beta_A)\cdot g(\boldsymbol\beta_B)
	=\boldsymbol\beta_A\cdot\boldsymbol\beta_B-(\vb n\cdot\boldsymbol\beta_A)(\vb n\cdot\boldsymbol\beta_B).
\end{equation*}
把 $(\vb n\cdot\boldsymbol\beta_A)(\vb n\cdot\boldsymbol\beta_B)$ 拆成 $1-(1-\vb n\cdot\boldsymbol\beta_A)-(1-\vb n\cdot\boldsymbol\beta_B)+(1-\vb n\cdot\boldsymbol\beta_A)(1-\vb n\cdot\boldsymbol\beta_B)$，除以分母后，常数项 $-(\sum_A\eta_Aq_A)^2$ 与含 $(1-\vb n\cdot\boldsymbol\beta_A)^{-1}$ 的单项（因子 $\sum_B\eta_Bq_B=0$）均因电荷守恒为零，只余
\begin{equation*}
	\left|\boldsymbol{\mathcal S}^{(0)}(\vb n)\right|^2
	=-
	\sum_{A,B}\eta_A\eta_Bq_Aq_B
	\frac{1-\boldsymbol\beta_A\cdot\boldsymbol\beta_B}
	{(1-\vb n\cdot\boldsymbol\beta_A)(1-\vb n\cdot\boldsymbol\beta_B)}.
\end{equation*}
所需角积分可以从两个初等公式推出：
\begin{equation*}
	\frac{1}{AB}=\int_0^1\frac{\d x}{[xA+(1-x)B]^2},
	\qquad
	\int\frac{\d\Omega}{(1-\vb n\cdot\vb w)^2}
	=\frac{4\pi}{1-w^2},
	\quad w<1.
\end{equation*}
第二个等式取极轴沿 $\vb w$ 后 $2\pi\int_{-1}^{1}\frac{\d c}{(1-wc)^2}$ 可直接积出。令
\begin{equation*}
	\Gamma_{AB}
	=\gamma_A\gamma_B
	\left(1-\boldsymbol\beta_A\cdot\boldsymbol\beta_B\right),
	\qquad
	\beta_{AB}=\sqrt{1-\Gamma_{AB}^{-2}},
\end{equation*}
其中 $\Gamma_{AB}$ 是两条渐近外线之间的相对 Lorentz 因子，$\beta_{AB}$ 是相应的相对速度。对成对求和形式使用 Feynman 参数，把两个分母合并在 $\vb w(x)=x\boldsymbol\beta_A+(1-x)\boldsymbol\beta_B$ 中并交换积分次序，得
\begin{equation*}
	\int\d\Omega\,
	\frac{1-\boldsymbol\beta_A\cdot\boldsymbol\beta_B}
	{(1-\vb n\cdot\boldsymbol\beta_A)(1-\vb n\cdot\boldsymbol\beta_B)}
	=4\pi
	(1-\boldsymbol\beta_A\cdot\boldsymbol\beta_B)
	\int_0^1\frac{\d x}
	{1-|x\boldsymbol\beta_A+(1-x)\boldsymbol\beta_B|^2}.
\end{equation*}
余下的 $x$ 积分是配平方后的初等积分：分母 $1-|\vb w(x)|^2$ 为 $x$ 的二次式，平移区间后化为 $\int_{-1}^{1}\frac{\d u}{\alpha^2-\gamma^2u^2}$ 型，由 $\Gamma_{AB}$ 的定义化简为
\begin{equation*}
	\int_0^1\frac{\d x}{1-|x\boldsymbol\beta_A+(1-x)\boldsymbol\beta_B|^2}
	=\frac{\operatorname{artanh}\beta_{AB}}
	{\beta_{AB}(1-\boldsymbol\beta_A\cdot\boldsymbol\beta_B)}.
\end{equation*}
（等速率、夹角 $\vartheta$ 时可直接复核：$1-|\vb w|^2=1-\beta^2\cos^2\frac{\vartheta}{2}-\beta^2\sin^2\frac{\vartheta}{2}(2x-1)^2$，平移区间后即化为 $\frac{1}{ab}\operatorname{artanh}(b/a)$ 型初等积分。）因此角积分后的相对论软谱具有完全闭合的形式
\begin{equation*}
	\left.\frac{\d E}{\d\omega}\right|_{\omega\to0}
	=-\frac{1}{4\pi^2\varepsilon_0c}
	\sum_{A,B}\eta_A\eta_Bq_Aq_B
	\left[
	\frac{\operatorname{artanh}\beta_{AB}}{\beta_{AB}}-1
	\right].
\end{equation*}
方括号中的 $-1$ 可以加入，是因为电荷守恒条件使相应的双重求和为零。当 $A=B$ 时，按连续极限 $\operatorname{artanh}\beta_{AA}/\beta_{AA}\to1$，自项自动消失。

最简单而重要的情形是一枚电荷 $q$ 由初速度 $\boldsymbol\beta_i$ 偏转到末速度 $\boldsymbol\beta_f$。此时只有一对入射、出射外线，闭合软谱中非零的只有 $A\neq B$ 两项，化为
\begin{equation*}
	\left.\frac{\d E}{\d\omega}\right|_{\omega\to0}
	=\frac{q^2}{2\pi^2\varepsilon_0c}
	\left[
	\frac{\operatorname{artanh}\beta_{if}}{\beta_{if}}-1
	\right]
	=\frac{q^2}{2\pi^2\varepsilon_0c}
	\left[
	\frac{\Gamma_{if}\operatorname{arccosh}\Gamma_{if}}
	{\sqrt{\Gamma_{if}^2-1}}-1
	\right].
\end{equation*}
若初末速率相同为 $\beta c$，夹角为 $\Theta$，则
\begin{equation*}
	\Gamma_{if}
	=\gamma^2(1-\beta^2\cos\Theta)
	=1+2\gamma^2\beta^2\sin^2\frac{\Theta}{2},
\end{equation*}
第二个等号用了 $\gamma^{-2}=1-\beta^2$ 与倍角公式。在非相对论极限，$\beta_{if}=|\boldsymbol\beta_f-\boldsymbol\beta_i|+O(\beta^3)$——把 $\beta_{if}^2=1-\Gamma_{if}^{-2}$ 按 $\beta^2$ 展开得 $\frac{(1-\beta^2\cos\Theta)^2-(1-\beta^2)^2}{(1-\beta^2\cos\Theta)^2}\simeq4\beta^2\sin^2\frac{\Theta}{2}$——并且
\begin{equation*}
	\frac{\operatorname{artanh}x}{x}-1
	=\frac{x^2}{3}+\frac{x^4}{5}+O(x^6),
\end{equation*}
取自 $\operatorname{artanh}$ 的幂级数。故单电荷软谱恢复
\begin{equation*}
	\left.\frac{\d E}{\d\omega}\right|_{\omega\to0}
	=\frac{q^2|\Delta\vb v|^2}{6\pi^2\varepsilon_0c^3}
	+O(c^{-5}),
\end{equation*}
这正是前面非相对论边界形式中的电偶极部分。相反，当 $\Gamma_{if}\gg1$ 时，
\begin{equation*}
	\frac{\Gamma_{if}\operatorname{arccosh}\Gamma_{if}}
	{\sqrt{\Gamma_{if}^2-1}}-1
	=\ln(2\Gamma_{if})-1+O(\Gamma_{if}^{-2}\ln\Gamma_{if}),
\end{equation*}
由 $\operatorname{arccosh}\Gamma\simeq\ln(2\Gamma)$ 展开得到。软能谱因相对论前向束射而出现对数增强。若散射角很小，则
\begin{equation*}
	\Gamma_{if}-1\simeq\frac{1}{2}\gamma^2\Theta^2.
\end{equation*}
于是 $\gamma\Theta\ll1$ 时软谱正比于 $\gamma^2\Theta^2$，而 $\gamma\Theta\gg1$ 时则按 $2\ln(\gamma\Theta)$ 增长。

零频振幅还等于远区辐射场对迟滞时间的积分。频域远区场在零频处的值正是该积分：$\widetilde{\vb E}_{\mathrm{rad}}(R\vb n,0)=\int\vb E_{\mathrm{rad}}(R\vb n,u)\d u$；沿用上面的边界值推导，若 $u=t_{\mathrm{obs}}-R/c$，则
\begin{equation*}
	\int_{-\infty}^{\infty}\vb E_{\mathrm{rad}}(R\vb n,u)\d u
	=\frac{1}{4\pi\varepsilon_0cR}
	\boldsymbol{\mathcal S}^{(0)}(\vb n).
\end{equation*}
因此领先软定理与电磁记忆效应是同一个边界量的频域和时域表述。另一方面，由于
\begin{equation*}
	\frac{\d N_\gamma}{\d\omega}
	=\frac{1}{\hbar\omega}\frac{\d E}{\d\omega},
\end{equation*}
频率为 $\omega$ 的光子携带能量 $\hbar\omega$。当 $\d E/\d\omega$ 在 $\omega\to0$ 趋于非零常数时，软光子数按 $1/\omega$ 发散，而总软辐射能量 $\int_0^{\omega_{\max}}(\d E/\d\omega)\d\omega$ 仍然有限。实验中有限观测时间、能量分辨率以及虚光子修正共同决定可分辨的红外下限。

还需注意，长程相互作用的次领先软展开一般不是普通的 $\omega$ 幂级数。短程力下，轨道在足够大的 $|t|$ 上严格趋于直线，领先软项之后可按 $\omega$ 展开；库仑力则给出
\begin{equation*}
	\vb r(t)=c\boldsymbol\beta_\pm t+\vb c_\pm\ln|t|+\vb b_\pm+O(t^{-1}),
	\qquad
	\boldsymbol\beta(t)=\boldsymbol\beta_\pm+\frac{\vb c_\pm}{ct}+O(t^{-2}).
\end{equation*}
其本质可从尾积分看出。对 $a>0$，由分部积分或指数积分定义，
\begin{equation*}
	\int_T^\infty\frac{e^{ia\omega t}}{t^2}\d t
	=\frac{e^{ia\omega T}}{T}+ia\omega E_1(-ia\omega T)
	=\frac{1}{T}-ia\omega\ln(\omega T)+O(\omega),
\end{equation*}
其中 $E_1$ 为指数积分，最后一个等号只显示了非解析部分。故四维库仑散射的软振幅一般具有
\begin{equation*}
	\boldsymbol{\mathcal A}(\omega,\vb n)
	=\boldsymbol{\mathcal S}^{(0)}(\vb n)
	+\omega\ln\omega\,\boldsymbol{\mathcal S}^{(\ln)}(\vb n)
	+O(\omega),
\end{equation*}
而不是简单的 $\boldsymbol{\mathcal S}^{(0)}+\omega\boldsymbol{\mathcal S}^{(1)}+\cdots$。领先常数项仍由领先软因子完全确定，但若要求次领先精度，就必须保留库仑轨道的长程对数尾。

硬频渐近则由时间复平面中离实轴最近的奇点或转折点控制。若某一实际进入辐射振幅的函数在上半平面最近奇点 $t_\star=t_0+i\tau_\star$ 附近具有
\begin{equation*}
	F(t)\sim A(t-t_\star)^{-\lambda},
	\qquad \tau_\star>0,
\end{equation*}
则由鞍点法（陡降路径）的初等形式，指数由奇点到实轴的距离决定、幂由奇点的指数决定，
\begin{equation*}
	\widetilde F(\omega)
	\sim C\,\omega^{\lambda-1}e^{i\omega t_\star},
	\qquad
	|\widetilde F(\omega)|^2
	\sim |C|^2\omega^{2\lambda-2}e^{-2\omega\tau_\star}.
\end{equation*}
因此硬频指数截止给出最近复时间奇点到实轴的距离；当奇点随轨道参数趋近实轴，或积分端点本身不光滑时，指数尾会退化为幂律或转折点一致渐近。

若运动是周期为 $T=2\pi/\Omega$ 的束缚运动，应使用 Fourier 级数
\begin{equation*}
	f_N=\frac{1}{T}\int_0^T f(t)e^{iN\Omega t}\d t,
	\qquad
	f(t)=\sum_{N=-\infty}^{\infty}f_Ne^{-iN\Omega t}.
\end{equation*}
实函数满足 $f_{-N}=f_N^*$。把连续谱公式中的频域多极矩换成周期情形的 $\widetilde{\vb p}(\omega)=2\pi\sum_N\vb p_N\delta(\omega-N\Omega)$，对正频率积分时只有 $N\geq1$ 的谱线贡献，例如
\begin{equation*}
	E_{E1}=\frac{1}{6\pi^2\varepsilon_0c^3}
	\int_0^\infty\omega^4|\widetilde{\vb p}(\omega)|^2\d\omega
	=\frac{4\pi^2}{6\pi^2\varepsilon_0c^3}
	\sum_{N=1}^{\infty}(N\Omega)^4p_{i,N}p_{i,N}^*,
\end{equation*}
其中利用了 $\delta$ 函数在正半轴的分离性 $\int_0^\infty\delta(\omega-N\Omega)\delta(\omega-M\Omega)\d\omega=\delta_{NM}$。每个周期的平均功率为 $\overline P=E/T=E\Omega/2\pi$，故第 $N\geq1$ 条谱线的周期平均功率为
\begin{equation*}
	&P_{E1,N}=\frac{(N\Omega)^4}{3\pi\varepsilon_0c^3}p_{i,N}p_{i,N}^*,\\
	&P_{M1,N}=\frac{(N\Omega)^4}{3\pi\varepsilon_0c^5}m_{i,N}m_{i,N}^*,\\
	&P_{E2,N}=\frac{(N\Omega)^6}{360\pi\varepsilon_0c^5}D_{ij,N}D_{ij,N}^*.
\end{equation*}
$M1$、$E2$ 由同样的代换得到。所有正整数谐波求和后恢复周期平均的时域功率：
\begin{equation*}
	\overline P_X=\sum_{N=1}^{\infty}P_{X,N}.
\end{equation*}
这可由周期 Parseval 恒等式 $\int_0^T|f(t)|^2\d t=T\sum_N|f_N|^2$ 直接验证。由此可见，连续谱和离散线谱只是 Fourier 变换与 Fourier 级数在非周期、周期运动中的两种表现：周期运动的频域场是 $\delta$ 函数谱，$\delta$ 平方在连续谱归一化中没有定义，必须改用周期 Parseval——这正是两者不能混用归一化的原因。下面以圆周运动为例，说明周期谱线公式在精确相对论运动中的用法。

\begin{example}
圆周运动的同步辐射谱.半径为 $a$、速率为 $\beta c$ 的圆周运动可以完全积分，是检验周期谱线公式并展示相对论效应如何改造谱形的最佳范例。取轨道位于 $xy$ 平面内
\begin{equation*}
	\vb r_q(t)=a(\sin\varphi,-\cos\varphi,0),
	\qquad
	\boldsymbol\beta=\beta(\cos\varphi,\sin\varphi,0),
	\qquad
	\varphi=\Omega t,
	\qquad
	\Omega=\frac{\beta c}{a}.
\end{equation*}
取场点方向与偏振基矢为
\begin{equation*}
	\uv R=(\cos\theta,0,\sin\theta),
	\qquad
	\uv e_a=(-\sin\theta,0,\cos\theta),
	\qquad
	\uv e_b=(0,1,0).
\end{equation*}

谐波谱的由来。由于轨道以 $T=2\pi/\Omega$ 为周期，精确谱积分中的被积函数（含完整迟滞相位 $e^{-i\omega\uv R\cdot\vb r_q/c}$）是固有时间 $t_s$ 的周期函数，其 Fourier 变换化为一列 $\delta$ 谱线：对任意周期函数 $g$，
\begin{equation*}
	\int_{-\infty}^{\infty}g(t_s)e^{i\omega t_s}\d t_s
	=2\pi\sum_{m=-\infty}^{\infty}g_m\delta(\omega-m\Omega),
	\qquad
	g_m=\frac{1}{T}\int_0^Tg(t_s)e^{im\Omega t_s}\d t_s.
\end{equation*}
这就是频率 $m\Omega=m\beta c/a$ 处谱线的来源。先用前面的全导数恒等式对精确谱积分分部积分（周期运动的边界项消失），横向矢量化为 $\uv R\times(\uv R\times\boldsymbol\beta)$；在 $(\uv e_a,\uv e_b,\uv R)$ 基中
\begin{equation*}
	\uv R\times(\uv R\times\boldsymbol\beta)
	=\beta\sin\theta\cos\varphi\,\uv e_a-\beta\sin\varphi\,\uv e_b,
	\qquad
	e^{-i\omega\uv R\cdot\vb r_q/c}=e^{-iu\sin\varphi},
	\quad u=\frac{\omega a\cos\theta}{c}.
\end{equation*}
对 $e^{-iu\sin\varphi}$ 用 Jacobi--Anger 展开 $e^{-iu\sin\varphi}=\sum_m J_m(u)e^{-im\varphi}$：$\cos\varphi\,e^{-iu\sin\varphi}$ 的 $m$ 阶 Fourier 系数为 $\frac12[J_{m-1}(u)+J_{m+1}(u)]=\frac{m}{u}J_m(u)$，$\sin\varphi\,e^{-iu\sin\varphi}$ 的为 $\frac{1}{2i}[J_{m+1}(u)-J_{m-1}(u)]=iJ_m'(u)$，其中用了 Bessel 递推。完成系数后，远区电场为
\begin{equation*}
	&\vb E_{\mathrm{rad}}(\vb r,\omega)
	=-\frac{iq\omega\beta e^{i\omega r/c}}{2\varepsilon_0cr}
	\sum_{m=-\infty}^{\infty}
	\delta\left(\omega-m\frac{\beta c}{a}\right)\\
	&\quad\times
	\left[
	-i\uv e_b J_m'\left(\frac{\omega a\cos\theta}{c}\right)
	+\uv e_a m\sin\theta
	\frac{J_m\left(\omega a\cos\theta/c\right)}{\omega a\cos\theta/c}
	\right].
\end{equation*}
故频谱由频率 $\omega=m\beta c/a$ 处的 $\delta$ 函数谱线构成：$m$ 次谐波中垂直于轨道平面的偏振分量由 $J_m'$ 控制，面内偏振分量由 $m\sin\theta\,J_m/(m\beta\cos\theta)$ 控制。对 $m\geq1$ 谐波应用周期 Parseval 恒等式，每周期能量在单位立体角中的分布为
\begin{equation*}
	&\frac{\d W_m}{\d\Omega}
	=\frac{2r^2T}{\mu_0c}|\vb E_m|^2
	=\frac{Tq^2\omega^2\beta^2}{8\pi^2\varepsilon_0c}
	\left[J_m'^2(m\beta\cos\theta)\right.\\
	&\qquad
	\left.+\frac{\sin^2\theta}{\beta^2\cos^2\theta}J_m^2(m\beta\cos\theta)\right],
\end{equation*}
其中 $\vb E_m=\widetilde{\vb E}_m/2\pi$，正负谐波贡献相同，最后在 $\omega=m\Omega$、$u=m\beta\cos\theta$ 处取值并完成系数。对 $m\geq1$ 的谐波求和可封闭完成，得到每周期能量的角分布
\begin{equation*}
	\frac{\d W}{\d\Omega}
	=\frac{q^2\beta^3}{64\pi\varepsilon_0a}
	\frac{4(1+\sin^2\theta)-\beta^2\cos^4\theta(1+3\beta^2)}
	{(1-\beta^2\cos^2\theta)^{7/2}},
\end{equation*}
再对立体角积分得每周期总能量
\begin{equation*}
	W=2\pi\int_{-\pi/2}^{\pi/2}\frac{\d W}{\d\Omega}\cos\theta\d\theta
	=\frac{q^2\gamma^4\beta^3}{3\varepsilon_0a},
\end{equation*}
折回功率 $P=W/T=q^2\gamma^4\beta^4c/(6\pi\varepsilon_0a^2)$，与相对论 Larmor 公式一致，验证了谱线求和校验在精确相对论情形下仍然成立。

连续化极限。当速度接近光速 ($\gamma\gg1$)、观测角很小 ($\theta\ll1$) 时，谱的特征宽度 $\sim\gamma^3\Omega$ 远大于谐波间距 $\Omega$，每个特征宽度内包含大量谐波，求和可化为积分：对缓慢变化的 $f$，
\begin{equation*}
	\sum_{m=-\infty}^{\infty}f(m)\delta(\omega-m\Omega)
	=\frac{1}{\Omega}f\left(\frac{\omega}{\Omega}\right),
	\qquad
	\Omega=\frac{\beta c}{a}.
\end{equation*}
此时谐波参数 $u=m\beta\cos\theta\simeq m\left(1-\frac{\theta^2+\gamma^{-2}}{2}\right)$ 与阶数 $m$ 之差为 $\frac{m}{2}(\theta^2+\gamma^{-2})$ 量级，适用 Bessel 函数的 Debye 一致渐近（可由 $K_\nu$ 的积分表示 $K_\nu(z)=\int_0^\infty e^{-z\cosh t}\cosh(\nu t)\d t$ 复核），上述线性组合化为修正贝塞尔函数
\begin{equation*}
	&\sum_{m=-\infty}^{\infty}
	\delta\left(\omega-m\frac{\beta c}{a}\right)
	\left(
	-i\uv e_b J_m'\left(\frac{\omega a\cos\theta}{c}\right)
	+\uv e_a\sin\theta
	\frac{mJ_m\left(\omega a\cos\theta/c\right)}{\omega a\cos\theta/c}
	\right)\\
	&\quad=-\frac{ia}{\sqrt{3}\pi\beta c}(\theta^2+\gamma^{-2})
	\left(
	\uv e_b K_{2/3}(\xi)
	+i\uv e_a\frac{\theta}{\sqrt{\theta^2+\gamma^{-2}}}K_{1/3}(\xi)
	\right),
	\qquad
	\xi=\frac{\omega a}{3c}(\theta^2+\gamma^{-2})^{3/2},
\end{equation*}
其中 $\xi=\frac{m}{3}(\theta^2+\gamma^{-2})^{3/2}$ 在 $\omega=m\Omega$ 处的值，且 $\beta\simeq1$。代入每周期能量谱即得连续化的角分布频谱
\begin{equation*}
	&\frac{\d^2W}{\d\Omega\d\omega}
	=\frac{q^2}{12\pi^3\varepsilon_0c}
	\left(\frac{\omega a}{c}\right)^2(\theta^2+\gamma^{-2})^2
	\left(
	(K_{2/3}(\xi))^2
	+\frac{\theta^2}{\theta^2+\gamma^{-2}}(K_{1/3}(\xi))^2
	\right),\\
	&\frac{\d W}{\d\Omega}
	=\frac{7q^2}{64\pi\varepsilon_0a(\theta^2+\gamma^{-2})^{5/2}}
	\left(
	1+\frac{5}{7}\frac{\theta^2}{\theta^2+\gamma^{-2}}
	\right),\\
	&W=2\pi\int_{-\infty}^{\infty}\frac{\d W}{\d\Omega}\d\theta
	=\frac{q^2\gamma^4}{3\varepsilon_0a}.
\end{equation*}
修正贝塞尔函数在 $\xi\gg1$ 时按 $K_\nu(\xi)\sim\sqrt{\pi/2\xi}\,e^{-\xi}$ 指数衰减，故每个观测方向上辐射显著的频率满足 $\xi\lesssim1$，即 $\theta^2+\gamma^{-2}\lesssim(3c/\omega a)^{2/3}$；在谱峰频率 $\omega\sim\gamma^3\beta c/a$ 处——$\theta=0$ 时谱 $\propto z^2K_{2/3}^2(z)$，$z=\omega a\gamma^{-3}/3c$，而 $z^2K_{2/3}^2(z)$ 在 $z\sim1$ 处最大——这一判据给出 $\theta\lesssim1/\gamma$。谱集中在前向 $\theta\sim1/\gamma$ 的窄锥内并展宽到约 $\gamma^3$ 条谐波：低速时圆周运动只辐射 $m=1$ 基波，$\gamma\gg1$ 时谱被相对论束射聚束展宽——相对论效应正是以这种锥形结构改造了谱形。
\end{example}

在经典散射问题中，固定碰撞参数 $\vb b$ 的 $\d E/\d\omega$ 描述一条给定轨道的辐射。若入射束在碰撞参数平面上均匀（单位面积、单位频率的注入能量相同），则定义能量加权的谱辐射截面
\begin{equation*}
	\frac{\d\chi_X}{\d\omega}
	=\int\d^2b\,\frac{\d E_X(\vb b)}{\d\omega}.
\end{equation*}
轴对称散射中 $\d^2b=b\d b\d\phi$ 的角向积分给 $2\pi$，即 $\d^2b=2\pi b\d b$。频率为 $\omega$ 的光子携带能量 $\hbar\omega$，相应的光子数截面为
\begin{equation*}
	\frac{\d\sigma_{\gamma,X}}{\d\omega}
	=\frac{1}{\hbar\omega}\frac{\d\chi_X}{\d\omega}.
\end{equation*}
谱辐射截面的大碰撞参数端通常控制红外行为，小碰撞参数端通常控制硬频行为；屏蔽、有限尺寸、量子反冲和相对论效应分别为这些端点提供物理截断。

前面的角积分频谱及周期谱线公式属于非相对论长波多极理论，其共同条件是
\begin{equation*}
	\frac{v_{\mathrm{src}}}{c}\ll1,
	\qquad
	\frac{\omega a_{\mathrm{eff}}(\omega)}{c}\ll1.
\end{equation*}
后一条件必须对实际贡献于该频率的轨道区段检查，而不能只以无穷远速度或某个固定几何尺度代替。通常 $E2$ 振幅相对 $E1$ 振幅低一阶 $\omega a/c$；若电偶极矩因对称性或荷质比关系抵消，则必须把四极项作为领先项保留。与此不同，领先软谱角分布和角积分软谱是完整相对论电动力学在 $\omega\to0$ 的领先结论，不要求 $v/c\ll1$ 或长波多极展开；它们的限制是渐近入射、出射状态必须存在，并且所研究频率满足软条件 $\omega\tau_{\mathrm{sc}}\ll1$。

因此，一个具体辐射问题的频谱分析应先判断所需的是完整频谱还是仅需软区。完整的非相对论频谱应先由动力学求出轨道并构造 $\vb p(t)$、$\vb m(t)$、$\vb D(t)$，再根据运动是否周期选择连续 Fourier 变换或 Fourier 级数约定；所得频域多极矩代入完整角分布或相应的角积分频谱公式。若只研究任意速度散射的领先软辐射，则不必求完整轨道 Fourier 变换，只需由散射动力学确定渐近四动量，再代入领先软谱角分布或角积分软谱。散射系综最后按谱辐射截面定义对碰撞参数积分；频谱归一化恒等式、非相对论软极限公式、相对论软极限公式和硬频渐近公式分别检查归一化与三个渐近区的行为。这些频域结果把轨道当作已知源；要让轨道本身由电磁场决定，还必须先解决点电荷自作用问题，这正是下面两节的任务。

\subsection{点粒子的辐射反作用}

上一节把电荷轨道当作已知源,完成了固定轨道上的频谱分析。现在要让轨道本身由电磁场决定,就必须回答一个更基本的问题:点电荷的推迟场在世界线上发散,哪一部分应当解释为粒子惯性的重定义,哪一部分才是有限的辐射反作用?这个问题不能从非相对论的“附加质量”概念猜测,而应先在四维时空中完成 Dirac 的场分解与守恒量计算,求得世界线上的有限正则场后,再在有限尺寸粒子的极限下读出熟悉的电磁质量与延迟自作用。

\subsubsection{四维 Green 函数与推迟、超前场}

解决上述问题的第一步,是把推迟场、超前场及其分解在四维时空中写清楚。为此统一使用国际单位制和号差 $+2$,固定记号
\begin{equation*}
 g_{\mu\nu}=\diag(1,-1,-1,-1),
 \qquad
 x^\mu=(ct,\vb x),
 \qquad
 A^\mu=(\Phi/c,\vb A),
 \qquad
 J^\mu=(c\rho,\vb J).
\end{equation*}
粒子世界线记为 $z^\mu(\tau)$,其四速度、四加速度和四维 jerk 分别为
\begin{equation*}
 u^\mu={\d z^\mu\over\d\tau},
 \qquad
 a^\mu={\d u^\mu\over\d\tau},
 \qquad
 \dot a^\mu={\d a^\mu\over\d\tau},
 \qquad
 u^2=c^2,
 \quad
 u\cdot a=0.
\end{equation*}
大写 $J^\mu$ 始终表示电流,点表示对固有时求导;世界管半径用 $\ell$,近世界线的类空位移用 $\xi^\mu$。记号约定小结:$\gamma$ 只保留给三维相对论中的 Lorentz 因子,不再兼任空间位移。

在 Lorenz 规范下,
\begin{equation*}
 F^{\mu\nu}=\partial^\mu A^\nu-\partial^\nu A^\mu,
 \qquad
 \partial_\mu A^\mu=0,
 \qquad
 \Box A^\mu=\mu_0J^\mu,
\end{equation*}
其中
\begin{equation*}
 \Box=\partial_\mu\partial^\mu
 ={\partial^2\over\partial(x^0)^2}-\nabla^2
 ={1\over c^2}{\partial^2\over\partial t^2}-\nabla^2.
\end{equation*}
为避免在一开始把因果边界条件隐藏在三维推迟势中,先求四维 Green 函数。令
\begin{equation*}
 \Box G(x)=\delta^{(4)}(x),
 \qquad
 G(x)=\int{\d^4k\over(2\pi)^4}\,\widetilde G(k)e^{-ik\cdot x}.
\end{equation*}
由于
\begin{equation*}
 \Box e^{-ik\cdot x}=-k^2e^{-ik\cdot x},
 \qquad
 k^2=(k^0)^2-\vb k^2,
\end{equation*}
动量空间 Green 函数必须满足
\begin{equation*}
 -k^2\widetilde G(k)=1.
\end{equation*}
分母在 $k^0=\pm|\vb k|$ 处有极点;不同的极点绕行规定给出不同的边界条件。推迟与超前 Green 函数分别写成
\begin{align*}
 G_{\mathrm{ret}}(x)
 &=-\int{\d^4k\over(2\pi)^4}
 {e^{-ik\cdot x}\over k^2+i0\,k^0},\\
 G_{\mathrm{adv}}(x)
 &=-\int{\d^4k\over(2\pi)^4}
 {e^{-ik\cdot x}\over k^2-i0\,k^0}.
\end{align*}
对 $k^0$ 作围道积分。对推迟解,当 $x^0>0$ 时闭合下半平面,两个极点都被压到下方;当 $x^0<0$ 时闭合上半平面而没有极点。因此
\begin{align*}
 G_{\mathrm{ret}}(x)
 &=\theta(x^0)
 \int{\d^3k\over(2\pi)^3}
 e^{i\vb k\cdot\vb x}
 {\sin(|\vb k|x^0)\over|\vb k|},\\
 G_{\mathrm{adv}}(x)
 &=-\theta(-x^0)
 \int{\d^3k\over(2\pi)^3}
 e^{i\vb k\cdot\vb x}
 {\sin(|\vb k|x^0)\over|\vb k|}.
\end{align*}
令 $r=|\vb x|$,利用
\begin{equation*}
 \int\d\Omega_{\vb k}\,e^{i\vb k\cdot\vb x}
 =4\pi{\sin(kr)\over kr},
\end{equation*}
推迟解的径向积分为
\begin{align*}
 G_{\mathrm{ret}}(x)
 &= {\theta(x^0)\over2\pi^2r}
 \int_0^\infty\d k\,\sin(kr)\sin(kx^0)\\
 &= {\theta(x^0)\over4\pi r}
 \left[\delta(x^0-r)-\delta(x^0+r)\right]\\
 &= {\delta(x^0-r)\over4\pi r}.
\end{align*}
这里第二步用到了 $\int_0^\infty\d k\,\sin(ka)\sin(kb)=(\pi/2)[\delta(a-b)-\delta(a+b)]$,最后一步利用了 $x^0>0$ 与 $r\geq0$。同理
\begin{equation*}
 G_{\mathrm{adv}}(x)={\delta(x^0+r)\over4\pi r}.
\end{equation*}
由于
\begin{equation*}
 \delta(x^2)
 =\delta((x^0)^2-r^2)
 ={\delta(x^0-r)+\delta(x^0+r)\over2r},
\end{equation*}
上面两式又可写成显式协变的光锥形式
\begin{equation*}
 G_{\mathrm{ret/adv}}(x)
 ={1\over2\pi}\theta(\pm x^0)\delta(x^2).
\end{equation*}
因此推迟与超前 Green 函数是同一个波动算符的两个基本解;它们只在对光锥两支的支集选择上不同。

给定任意守恒电流,两个相应的特解为
\begin{equation*}
 A_{\mathrm{ret/adv}}^\mu(x)
 =\mu_0\int G_{\mathrm{ret/adv}}(x-x')J^\mu(x')\d^4x'.
\end{equation*}
它们满足同一个非齐次方程,故半和仍含源而半差无源:
\begin{equation*}
 \Box{A_{\mathrm{ret}}^\mu+A_{\mathrm{adv}}^\mu\over2}
 =\mu_0J^\mu,
 \qquad
 \Box{A_{\mathrm{ret}}^\mu-A_{\mathrm{adv}}^\mu\over2}=0.
\end{equation*}

点电荷沿世界线 $z^\mu(\tau)$ 运动时,四维电流为
\begin{equation*}
 J^\mu(x)=qc\int_{-\infty}^{\infty}
 u^\mu(\tau')\delta^{(4)}(x-z(\tau'))\d\tau'.
\end{equation*}
把它代入上面的势公式,记
\begin{equation*}
 R^\mu(\tau')=x^\mu-z^\mu(\tau'),
 \qquad
 \Phi(\tau')=R^2(\tau'),
\end{equation*}
则
\begin{equation*}
 {\d\Phi\over\d\tau'}
 =2R\cdot{\d R\over\d\tau'}
 =-2R\cdot u.
\end{equation*}
利用单变量 delta 函数公式
\begin{equation*}
 \delta(\Phi(\tau'))
 =\sum_i{\delta(\tau'-\tau_i)\over|\Phi'(\tau_i)|}
 =\sum_i{\delta(\tau'-\tau_i)\over2|R\cdot u|_{\tau_i}},
\end{equation*}
其中 $\tau_i$ 是 $R^2=0$ 的根,便得到
\begin{align*}
 A_{\mathrm{ret}}^\mu(x)
 &= {\mu_0qc\over4\pi}
 \left.{u^\mu\over R\cdot u}\right|_{\tau_{\mathrm{ret}}},\\
 A_{\mathrm{adv}}^\mu(x)
 &=-{\mu_0qc\over4\pi}
 \left.{u^\mu\over R\cdot u}\right|_{\tau_{\mathrm{adv}}}.
\end{align*}
这里推迟根位于过去光锥,满足 $R\cdot u>0$;超前根位于未来光锥,满足 $R\cdot u<0$,故第二式的负号来自绝对值本身,并不是额外的约定。

实际场是同一源方程的某个物理解。任取常数 $\alpha$,它都可写成
\begin{equation*}
 A_{\mathrm{act}}^\mu
 =\alpha A_{\mathrm{ret}}^\mu
 +(1-\alpha)A_{\mathrm{adv}}^\mu
 +A_{\mathrm{hom}}^\mu,
\end{equation*}
其中 $A_{\mathrm{hom}}^\mu$ 是齐次解。对同一个 $A_{\mathrm{act}}^\mu$ 重新分组,定义
\begin{align*}
 A_{\mathrm{in}}^\mu
 &=(1-\alpha)(A_{\mathrm{adv}}^\mu-A_{\mathrm{ret}}^\mu)
 +A_{\mathrm{hom}}^\mu,\\
 A_{\mathrm{out}}^\mu
 &=\alpha(A_{\mathrm{ret}}^\mu-A_{\mathrm{adv}}^\mu)
 +A_{\mathrm{hom}}^\mu,
\end{align*}
则
\begin{equation*}
 A_{\mathrm{act}}^\mu
 =A_{\mathrm{ret}}^\mu+A_{\mathrm{in}}^\mu
 =A_{\mathrm{adv}}^\mu+A_{\mathrm{out}}^\mu,
 \qquad
 A_{\mathrm{out}}^\mu-A_{\mathrm{in}}^\mu
 =A_{\mathrm{ret}}^\mu-A_{\mathrm{adv}}^\mu.
\end{equation*}
因此 $A_{\mathrm{in}}^\mu$ 与 $A_{\mathrm{out}}^\mu$ 都是无源场。对场强作同样的分解,定义 Dirac 奇异场和辐射场
\begin{equation*}
 F_{\mathrm S}^{\mu\nu}
 ={F_{\mathrm{ret}}^{\mu\nu}+F_{\mathrm{adv}}^{\mu\nu}\over2},
 \qquad
 F_{\mathrm R}^{\mu\nu}
 ={F_{\mathrm{ret}}^{\mu\nu}-F_{\mathrm{adv}}^{\mu\nu}\over2},
\end{equation*}
并把在世界线上将保持有限的总正则场记为
\begin{equation*}
 F_{\mathrm{reg}}^{\mu\nu}
 \equiv F_{\mathrm{in}}^{\mu\nu}+F_{\mathrm R}^{\mu\nu}.
\end{equation*}
于是同一个实际场具有连续的等价改写链
\begin{equation*}
 F_{\mathrm{act}}^{\mu\nu}
 =F_{\mathrm{ret}}^{\mu\nu}+F_{\mathrm{in}}^{\mu\nu}
 =F_{\mathrm{adv}}^{\mu\nu}+F_{\mathrm{out}}^{\mu\nu}
 =F_{\mathrm S}^{\mu\nu}+F_{\mathrm{reg}}^{\mu\nu}.
\end{equation*}
后面的近世界线计算将证明:$F_{\mathrm S}^{\mu\nu}$ 精确收集推迟场和超前场共有的局部发散项,而 $F_{\mathrm R}^{\mu\nu}$ 在世界线上有有限极限。超前场在这里只是局部奇异性分解的数学工具;实际散射问题的因果边界条件仍由 $F_{\mathrm{in}}^{\mu\nu}$ 决定。

\subsubsection{近世界线展开：推迟与超前场的共同奇异部分}

上面的等价改写链把问题归结为:读出 $F_{\mathrm S}^{\mu\nu}$ 的局部奇性与 $F_{\mathrm R}^{\mu\nu}$ 的世界线极限,这要求在近世界线区域把推迟场与超前场统一展开。以下全部使用带量纲的四速度、四加速度和 jerk,并以固有时为参数;近线小量是管半径 $\ell$。由 $u^2=c^2$ 连续微分得到
\begin{equation*}
 u^2=c^2,
 \qquad
 u\cdot a=0,
 \qquad
 u\cdot\dot a=-a^2.
\end{equation*}
前两条说明四加速度恒与四速度正交;第三条由 $u\cdot a=0$ 再求导得 $u\cdot\dot a+a^2=0$,是 jerk 与加速度的内积关系,在下面的分母展开中会反复用到。

先把上面的 Liénard--Wiechert 势公式对观测点求导。由光锥条件 $R^2=0$,
\begin{equation*}
 0=\partial_\alpha(R^2)
 =2R_\alpha-2(R\cdot u){\partial\tau'\over\partial x^\alpha},
\end{equation*}
故
\begin{equation*}
 {\partial\tau'\over\partial x^\alpha}
 ={R_\alpha\over R\cdot u}.
\end{equation*}
将上式用于势 $A^\beta=\mu_0qc\,u^\beta/(4\pi\,R\cdot u)$ 的导数,链式法则给出
\begin{equation*}
 \partial_\alpha A^\beta
 ={\mu_0qc\over4\pi}{1\over R\cdot u}
 R_\alpha{\d\over\d\tau'}
 \left({u^\beta\over R\cdot u}\right),
\end{equation*}
再反对称化。利用 $\d R^\alpha/\d\tau'=-u^\alpha$ 把对 $u^\beta/(R\cdot u)$ 的求导转移到 $R^\alpha u^\beta$ 上:
\begin{equation*}
 R^\alpha{\d\over\d\tau'}\left({u^\beta\over R\cdot u}\right)
 ={\d\over\d\tau'}\left({R^\alpha u^\beta\over R\cdot u}\right)
 +{u^\alpha u^\beta\over R\cdot u},
\end{equation*}
组合成反对称后末项的贡献为 $u^{[\alpha}u^{\beta]}=0$,并代回 $\mu_0qc/(4\pi)=q/(4\pi\varepsilon_0c)$,推迟与超前场可统一写成
\begin{equation*}
 F_\eta^{\mu\nu}(x)
 =\left.
 \eta{q\over4\pi\varepsilon_0c}
 {1\over R\cdot u}{\d\over\d\tau'}
 \left({R^{[\mu}u^{\nu]}\over R\cdot u}\right)
 \right|_{\tau'=\tau_\eta},
 \qquad
 \eta=\begin{cases}+1,&\mathrm{ret},\\-1,&\mathrm{adv}.\end{cases}
\end{equation*}
其中 $X^{[\mu}Y^{\nu]}=X^\mu Y^\nu-X^\nu Y^\mu$。

在世界线上固定一点 $\tau_0$。取观测点位于与 $u^\mu(\tau_0)$ 正交的瞬时静止超平面:
\begin{equation*}
 x^\mu=z^\mu(\tau_0)+\xi^\mu,
 \qquad
 \xi\cdot u=0,
 \qquad
 \xi^2=-\ell^2,
 \qquad
 \xi^\mu=\ell n^\mu,
 \quad
 n^2=-1.
\end{equation*}
令源点参数为
\begin{equation*}
 \tau'=\tau_0-\sigma.
\end{equation*}
在 $\tau_0$ 处 Taylor 展开世界线和四速度:
\begin{align*}
 z^\mu(\tau_0-\sigma)
 &=z^\mu(\tau_0)-u^\mu\sigma
 +{1\over2}a^\mu\sigma^2
 -{1\over6}\dot a^\mu\sigma^3
 +O(\sigma^4),\\
 u^\mu(\tau_0-\sigma)
 &=u^\mu-a^\mu\sigma
 +{1\over2}\dot a^\mu\sigma^2
 +O(\sigma^3),
\end{align*}
更高一阶的世界线导数只出现在被略去的项中。于是
\begin{equation*}
 R^\mu(\tau')
 =\xi^\mu+u^\mu\sigma
 -{1\over2}a^\mu\sigma^2
 +{1\over6}\dot a^\mu\sigma^3
 +O(\sigma^4).
\end{equation*}
近世界线时 $\xi=O(\ell)$,而 $\sigma$ 为同阶的固有时差($c$ 的幂次只补量纲,不影响计阶)。为了把场算到有限阶,需要把统一场式中的分子、分母和光锥根一致展开。

先计算反对称分子。把世界线展开与 $R^\mu$ 的展开逐项相乘,按 $\sigma$ 幂次收集
\begin{align*}
 R^{[\mu}u^{\nu]}(\tau')={}&\xi^{[\mu}u^{\nu]}
 +\left(u^{[\mu}u^{\nu]}-\xi^{[\mu}a^{\nu]}\right)\sigma\\
 &+\left({1\over2}\xi^{[\mu}\dot a^{\nu]}
 -u^{[\mu}a^{\nu]}
 -{1\over2}a^{[\mu}u^{\nu]}\right)\sigma^2\\
 &+\left({1\over2}u^{[\mu}\dot a^{\nu]}
 +{1\over2}a^{[\mu}a^{\nu]}
 +{1\over6}\dot a^{[\mu}u^{\nu]}\right)\sigma^3
 +O(\sigma^4).
\end{align*}
利用
\begin{equation*}
 u^{[\mu}u^{\nu]}=0,
 \qquad
 a^{[\mu}a^{\nu]}=0,
 \qquad
 a^{[\mu}u^{\nu]}=-u^{[\mu}a^{\nu]},
 \qquad
 u^{[\mu}\dot a^{\nu]}=-\dot a^{[\mu}u^{\nu]},
\end{equation*}
并注意更高一阶的世界线导数项最终只进入比所需更高的阶,上式化成
\begin{equation*}
 R^{[\mu}u^{\nu]}(\tau')
 =\xi^{[\mu}u^{\nu]}
 -\xi^{[\mu}a^{\nu]}\sigma
 +\left({1\over2}\xi^{[\mu}\dot a^{\nu]}
 -{1\over2}u^{[\mu}a^{\nu]}\right)\sigma^2
 -{1\over3}\dot a^{[\mu}u^{\nu]}\sigma^3
 +O(\ell^4).
\end{equation*}
这里 $\sigma^2$ 系数中的 $-u^{[\mu}a^{\nu]}-a^{[\mu}u^{\nu]}/2$ 合并为 $-u^{[\mu}a^{\nu]}/2$,而 $\sigma^3$ 系数中的 $a^{[\mu}a^{\nu]}$ 为零、$u^{[\mu}\dot a^{\nu]}$ 与 $\dot a^{[\mu}u^{\nu]}$ 合并为 $-(2/3)\dot a^{[\mu}u^{\nu]}$。

再计算分母。把 $R^\mu(\tau')$ 与 $u^\mu(\tau')$ 逐项收缩:
\begin{align*}
 R\cdot u(\tau')={}&\xi\cdot u
 +\left(u^2-\xi\cdot a\right)\sigma
 +\left({1\over2}\xi\cdot\dot a
 -u\cdot a
 -{1\over2}a\cdot u\right)\sigma^2\\
 &+\left({1\over2}u\cdot\dot a
 +{1\over2}a^2
 +{1\over6}\dot a\cdot u\right)\sigma^3
 +O(\sigma^4).
\end{align*}
利用近点条件与运动学恒等式,并略去只影响更高阶的世界线导数项,得到
\begin{equation*}
 R\cdot u(\tau')
 =c^2\Lambda\sigma
 +{1\over2}(\xi\cdot\dot a)\sigma^2
 -{1\over6}a^2\sigma^3
 +O(\sigma^4),
 \qquad
 \Lambda\equiv1-{\xi\cdot a\over c^2}.
\end{equation*}
逐项核对:$\sigma$ 项由 $\xi\cdot u=0$ 与 $u^2=c^2$ 给出 $c^2-\xi\cdot a$;$\sigma^2$ 项中 $u\cdot a$ 的两次出现相互抵消,只剩 $\xi\cdot\dot a/2$;$\sigma^3$ 项的系数 $u\cdot\dot a/2+a^2/2+\dot a\cdot u/6$ 在 $u\cdot\dot a=-a^2$ 代入后为 $-a^2/6$,并不为零,这一项正是下面分母倒数展开中 $a^2\sigma^2/(6c^2\Lambda)$ 的来源。由于 $\xi\cdot a=O(\ell)$,有 $\Lambda=1+O(\ell)$。

为了随后对统一场式求导,需要先展开分母倒数。把上面的分母写成
\begin{equation*}
 R\cdot u=c^2\Lambda\sigma
 \left[1+{\xi\cdot\dot a\over2c^2\Lambda}\sigma
 -{a^2\over6c^2\Lambda}\sigma^2
 +O(\sigma^3)\right],
\end{equation*}
对括号作 $(1+X)^{-1}=1-X+X^2+\cdots$ 展开,
\begin{align*}
 {1\over R\cdot u}
 &={1\over c^2\Lambda\sigma}
 \left[
 1-{\xi\cdot\dot a\over2c^2\Lambda}\sigma
 +{a^2\over6c^2\Lambda}\sigma^2
 +{(\xi\cdot\dot a)^2\over4c^4\Lambda^2}\sigma^2
 +O(\ell^3)
 \right].
\end{align*}
这里 $(\xi\cdot\dot a)^2\sigma^2/c^4=O(\ell^4)$,在最终场的有限阶以前不产生贡献,故以后可省略。将化简后的分子与分母倒数相乘,先不取光锥根,得到
\begin{align*}
 {R^{[\mu}u^{\nu]}\over R\cdot u}
 ={1\over c^2\Lambda\sigma}\bigg[{}&
 \xi^{[\mu}u^{\nu]}
 -\xi^{[\mu}a^{\nu]}\sigma
 -{(\xi\cdot\dot a)\xi^{[\mu}u^{\nu]}\over2c^2\Lambda}\sigma\\
 &+{a^2\xi^{[\mu}u^{\nu]}\over6c^2\Lambda}\sigma^2
 +{1\over2}\xi^{[\mu}\dot a^{\nu]}\sigma^2
 -{1\over2}u^{[\mu}a^{\nu]}\sigma^2\\
 &+{(\xi\cdot\dot a)\xi^{[\mu}a^{\nu]}\over2c^2\Lambda}\sigma^2
 -{1\over3}\dot a^{[\mu}u^{\nu]}\sigma^3
 +O(\ell^4)
 \bigg].
\end{align*}
式中含 $\xi^{[\mu}a^{\nu]}\sigma^2$ 的交叉项在再求导并乘一个 $(R\cdot u)^{-1}$ 后只影响 $O(\ell)$,故可在所需精度下略去;同样,在方括号内部可把分母中的 $\Lambda$ 替换为 $1$——这只会丢弃 $O(\ell)$ 的修正——只有稍后与 $u^{[\mu}a^{\nu]}$ 相乘的线性 $\xi\cdot a$ 修正必须保留。于是
\begin{align*}
 {R^{[\mu}u^{\nu]}\over R\cdot u}
 ={1\over c^2\Lambda\sigma}\bigg[{}&
 \xi^{[\mu}u^{\nu]}
 -\xi^{[\mu}a^{\nu]}\sigma
 -{(\xi\cdot\dot a)\xi^{[\mu}u^{\nu]}\over2c^2}\sigma\\
 &+{a^2\xi^{[\mu}u^{\nu]}\over6c^2}\sigma^2
 +{1\over2}\xi^{[\mu}\dot a^{\nu]}\sigma^2
 -{1\over2}u^{[\mu}a^{\nu]}\sigma^2
 -{1\over3}\dot a^{[\mu}u^{\nu]}\sigma^3
 \bigg]+O(\ell^3).
\end{align*}
由于 $\tau'=\tau_0-\sigma$,有 $\d/\d\tau'=-\d/\d\sigma$。对上面的比值逐项求导,含 $\sigma^{-1}$ 的首项把幂次降低一级:
\begin{align*}
 {\d\over\d\tau'}
 \left({R^{[\mu}u^{\nu]}\over R\cdot u}\right)
 ={1\over c^2\Lambda}\bigg[{}&
 {\xi^{[\mu}u^{\nu]}\over\sigma^2}
 -{a^2\xi^{[\mu}u^{\nu]}\over6c^2}
 -{1\over2}\xi^{[\mu}\dot a^{\nu]}
 +{1\over2}u^{[\mu}a^{\nu]}
 +{2\over3}\dot a^{[\mu}u^{\nu]}\sigma
 \bigg]+O(\ell^2).
\end{align*}
现在再乘分母倒数。其中由外层展开的 $a^2\sigma^2/(6c^2\Lambda)$ 项与导数首项 $\xi^{[\mu}u^{\nu]}/\sigma^2$ 相乘产生的 $+a^2\xi^{[\mu}u^{\nu]}/(6c^2\Lambda)$,恰好抵消导数式中显式的 $-a^2\xi^{[\mu}u^{\nu]}/(6c^2)$;两者在乘以外层因子后均只差 $O(\ell)$,落在已声明的精度之内。因此在尚未代入光锥根时,场化成
\begin{align*}
 F_\eta^{\mu\nu}
 =\eta{q\over4\pi\varepsilon_0c^5\Lambda^2}
 \bigg[{}&
 {\xi^{[\mu}u^{\nu]}\over\sigma^3}
 -{(\xi\cdot\dot a)\xi^{[\mu}u^{\nu]}\over2c^2\sigma^2}
 -{\xi^{[\mu}\dot a^{\nu]}\over2\sigma}
 +{u^{[\mu}a^{\nu]}\over2\sigma}
 +{2\over3}\dot a^{[\mu}u^{\nu]}
 \bigg]_{\sigma=\sigma_\eta}
 +O(\ell).
\end{align*}

接下来必须求真正的推迟与超前根;如果简单令 $\sigma=\pm\ell/c$,会漏掉与发散幂相乘后留下的有限项。由 $R^\mu$ 的展开,
\begin{equation*}
 R^2
 =\xi^2+2\xi\cdot u\,\sigma
 +\left(u^2-\xi\cdot a\right)\sigma^2
 +\left({1\over3}\xi\cdot\dot a-u\cdot a\right)\sigma^3
 +\left({1\over4}a^2+{1\over3}u\cdot\dot a\right)\sigma^4
 +O(\ell^5).
\end{equation*}
代入 $\xi^2=-\ell^2$、$\xi\cdot u=0$、$u^2=c^2$、$u\cdot a=0$ 和 $u\cdot\dot a=-a^2$:一次项消失,$\sigma^3$ 系数中的 $u\cdot a$ 消失,而 $\sigma^4$ 系数化为 $a^2/4-a^2/3=-a^2/12$。光锥条件化成
\begin{equation*}
 0=R^2
 =-\ell^2+c^2\Lambda\sigma^2
 +{1\over3}(\xi\cdot\dot a)\sigma^3
 -{1\over12}a^2\sigma^4
 +O(\ell^5).
\end{equation*}
令 $\eta=+1$ 表示推迟根、$\eta=-1$ 表示超前根,并作 ansatz
\begin{equation*}
 \sigma_\eta
 =\eta{\ell\over c\sqrt{\Lambda}}
 \left(1+A_\eta\ell+B\ell^2+O(\ell^3)\right).
\end{equation*}
由上式,
\begin{align*}
 \sigma_\eta^2
 &={\ell^2\over c^2\Lambda}
 \left[1+2A_\eta\ell+(2B+A_\eta^2)\ell^2+O(\ell^3)\right],\\
 \sigma_\eta^3
 &=\eta{\ell^3\over c^3\Lambda^{3/2}}
 \left[1+3A_\eta\ell+O(\ell^2)\right],\\
 \sigma_\eta^4
 &={\ell^4\over c^4\Lambda^2}
 \left[1+O(\ell)\right].
\end{align*}
代入光锥条件:首项 $-\ell^2+c^2\Lambda\sigma_\eta^2$ 中的 $\ell^2$ 部分相消,除以 $\ell^2$ 后得到
\begin{align*}
 0={}&2A_\eta\ell
 +(2B+A_\eta^2)\ell^2
 +\eta{\xi\cdot\dot a\over3c^3\Lambda^{3/2}}\ell
 \left(1+3A_\eta\ell\right)
 -{a^2\over12c^4\Lambda^2}\ell^2
 +O(\ell^3).
\end{align*}
注意 $\xi\cdot\dot a=O(\ell)$,它保证组合 $(\xi\cdot\dot a)/c^3$ 的量纲与量级正确,所以 $A_\eta=O(\ell)$;因而 $A_\eta^2\ell^2$ 和 $(\xi\cdot\dot a)A_\eta\ell^2/c^3$ 都是 $O(\ell^4)$,超出当前精度。逐阶比较:由 $O(\ell)$ 阶,
\begin{equation*}
 2A_\eta+\eta{\xi\cdot\dot a\over3c^3\Lambda^{3/2}}=0,
\end{equation*}
由 $O(\ell^2)$ 阶(略去 $A_\eta^2$ 与 $(\xi\cdot\dot a)A_\eta/c^3$),
\begin{equation*}
 2B-{a^2\over12c^4\Lambda^2}=0,
\end{equation*}
所以
\begin{equation*}
 A_\eta=-\eta{\xi\cdot\dot a\over6c^3\Lambda^{3/2}},
 \qquad
 B={a^2\over24c^4\Lambda^2},
\end{equation*}
\begin{equation*}
 \sigma_\eta
 =\eta{\ell\over c\sqrt{\Lambda}}
 \left[
 1-\eta{\xi\cdot\dot a\over6c^3\Lambda^{3/2}}\ell
 +{a^2\over24c^4\Lambda^2}\ell^2
 +O(\ell^3)
 \right].
\end{equation*}
推迟根 $\sigma_{+}>0$,超前根 $\sigma_{-}<0$,与 $\tau'=\tau_0-\sigma$ 的定义一致。

场式含 $\sigma^{-3}$、$\sigma^{-2}$ 和 $\sigma^{-1}$,故把 $(1+A_\eta\ell+B\ell^2)^{-k}$ 按二项式展开到 $\ell^2$ 阶($A_\eta^2$ 部分为 $O(\ell^4)$ 可略),分别得到
\begin{align*}
 \sigma_\eta^{-1}
 &=\eta{c\sqrt{\Lambda}\over\ell}
 \left[
 1+\eta{\xi\cdot\dot a\over6c^3\Lambda^{3/2}}\ell
 -{a^2\over24c^4\Lambda^2}\ell^2
 +O(\ell^3)
 \right],\\
 \sigma_\eta^{-2}
 &={c^2\Lambda\over\ell^2}
 \left[
 1+\eta{\xi\cdot\dot a\over3c^3\Lambda^{3/2}}\ell
 -{a^2\over12c^4\Lambda^2}\ell^2
 +O(\ell^3)
 \right],\\
 \sigma_\eta^{-3}
 &=\eta{c^3\Lambda^{3/2}\over\ell^3}
 \left[
 1+\eta{\xi\cdot\dot a\over2c^3\Lambda^{3/2}}\ell
 -{a^2\over8c^4\Lambda^2}\ell^2
 +O(\ell^3)
 \right].
\end{align*}
现在逐项代入场式。首先,$\sigma^{-3}$ 项给出
\begin{align*}
 \eta{q\over4\pi\varepsilon_0c^5\Lambda^2}
 {\xi^{[\mu}u^{\nu]}\over\sigma_\eta^3}
 ={q\over4\pi\varepsilon_0c}
 \Lambda^{-1/2}\bigg[
 {\xi^{[\mu}u^{\nu]}\over c\ell^3}
 +\eta{(\xi\cdot\dot a)\xi^{[\mu}u^{\nu]}\over2c^4\Lambda^{3/2}\ell^2}
 -{a^2\xi^{[\mu}u^{\nu]}\over8c^5\Lambda^2\ell}
 \bigg]+O(\ell).
\end{align*}
其次,$\sigma^{-2}$ 项给出
\begin{equation*}
 -\eta{q\over4\pi\varepsilon_0c^5\Lambda^2}
 {(\xi\cdot\dot a)\xi^{[\mu}u^{\nu]}\over2c^2\sigma_\eta^2}
 =-{q\over4\pi\varepsilon_0c}
 \Lambda^{-1/2}
 {\eta(\xi\cdot\dot a)\xi^{[\mu}u^{\nu]}\over2c^4\Lambda^{1/2}\ell^2}
 +O(1).
\end{equation*}
两式的 $\ell^{-2}$ 项分别为 $+\eta q(\xi\cdot\dot a)\xi^{[\mu}u^{\nu]}/(8\pi\varepsilon_0c^5\Lambda^2\ell^2)$ 与 $-\eta q(\xi\cdot\dot a)\xi^{[\mu}u^{\nu]}/(8\pi\varepsilon_0c^5\Lambda\ell^2)$;将 $\Lambda^{-1}=1+\xi\cdot a/c^2+O(\ell^2)$ 一致展开后,两者在 $O(1)$ 阶恰好相消,残差 $O(\ell)$ 落在已声明的精度之内。这个抵消必须同时使用光锥根的加速度修正(它产生 $\sigma^{-3}$ 展开中的 $\ell$ 修正项)和场的显式 $\sigma^{-2}$ 项;若直接令 $\sigma_\eta=\eta\ell/c$,便看不到这一结构。

对两个 $\sigma^{-1}$ 项,使用逆幂展开并保留 $\Lambda=1-\xi\cdot a/c^2$ 的线性修正,得到
\begin{align*}
 -\eta{q\over4\pi\varepsilon_0c^5\Lambda^2}
 {\xi^{[\mu}\dot a^{\nu]}\over2\sigma_\eta}
 &=-{q\over4\pi\varepsilon_0c}
 \Lambda^{-1/2}
 {\xi^{[\mu}\dot a^{\nu]}\over2c^3\ell}
 +O(1),\\
 \eta{q\over4\pi\varepsilon_0c^5\Lambda^2}
 {u^{[\mu}a^{\nu]}\over2\sigma_\eta}
 &={q\over4\pi\varepsilon_0c}
 \Lambda^{-1/2}
 {(1+\xi\cdot a/c^2)u^{[\mu}a^{\nu]}\over2c^3\ell}
 +O(1).
\end{align*}
最后,有限项保留分支符号:
\begin{equation*}
 \eta{q\over4\pi\varepsilon_0c^5\Lambda^2}
 {2\over3}\dot a^{[\mu}u^{\nu]}
 =\eta{q\over4\pi\varepsilon_0c}
 \Lambda^{-1/2}
 {2\over3c^4}\dot a^{[\mu}u^{\nu]}
 +O(\ell).
\end{equation*}
合并以上各式,推迟场和超前场的近线展开化成
\begin{align*}
 F_\eta^{\mu\nu}
 ={q\over4\pi\varepsilon_0c}\Lambda^{-1/2}
 \bigg[{}&
 {\xi^{[\mu}u^{\nu]}\over c\ell^3}
 -{a^2\xi^{[\mu}u^{\nu]}\over8c^5\ell}
 -{\xi^{[\mu}\dot a^{\nu]}\over2c^3\ell}
 +{(1+\xi\cdot a/c^2)u^{[\mu}a^{\nu]}\over2c^3\ell}
 +\eta{2\over3c^4}\dot a^{[\mu}u^{\nu]}
 \bigg]+O(\ell).
\end{align*}
上面的展开式是前面冗长计算的简洁形式,并且仍描述同一个 $F_\eta^{\mu\nu}$。现在可以直接读出:所有发散项都与 $\eta$ 无关,只有有限的最后一项在推迟和超前两支间变号。因此
\begin{align*}
 F_{\mathrm S}^{\mu\nu}
 &={F_{+}^{\mu\nu}+F_{-}^{\mu\nu}\over2}\\
 &={q\over4\pi\varepsilon_0c}\Lambda^{-1/2}
 \left[
 {\xi^{[\mu}u^{\nu]}\over c\ell^3}
 -{a^2\xi^{[\mu}u^{\nu]}\over8c^5\ell}
 -{\xi^{[\mu}\dot a^{\nu]}\over2c^3\ell}
 +{(1+\xi\cdot a/c^2)u^{[\mu}a^{\nu]}\over2c^3\ell}
 \right]+O(\ell),
\end{align*}
而两支之差为
\begin{equation*}
 F_{+}^{\mu\nu}-F_{-}^{\mu\nu}
 ={q\over4\pi\varepsilon_0c}\Lambda^{-1/2}
 {4\over3c^4}\dot a^{[\mu}u^{\nu]}+O(\ell),
\end{equation*}
取 $\ell\to0$ 时 $\Lambda\to1$,故
\begin{equation*}
 F_{\mathrm R}^{\mu\nu}(z)
 ={1\over2}\lim_{\ell\to0}
 \left(F_{+}^{\mu\nu}-F_{-}^{\mu\nu}\right)
 ={q\over6\pi\varepsilon_0c}
 {\dot a^{[\mu}u^{\nu]}\over c^4}
 ={q\over6\pi\varepsilon_0c^5}
 \left(\dot a^\mu u^\nu-\dot a^\nu u^\mu\right).
\end{equation*}
于是同一个实际场的等价改写链现在具有明确的局部内容:$F_{\mathrm S}$ 是时间对称且携带全部 Coulomb 型奇性的部分,$F_{\mathrm{reg}}=F_{\mathrm{in}}+F_{\mathrm R}$ 在世界线上有限。奇异场携带的四动量将交给世界管上的守恒量计算。

\subsubsection{世界管几何与管壁四动量通量}

上一小节只证明了正则场 $F_{\mathrm{reg}}$ 在世界线上有有限极限;这还不足以确定运动方程,因为质量重整化来自奇异场携带的四动量,而不是世界线上的场值。为此,在世界线周围取固有半径(管壁到世界线的固有距离)为 $\ell$ 的世界管,计算奇异场穿过管壁的四动量通量。管壁参数化为

\begin{equation*}
 y^\mu(\tau,\theta,\phi)
 =z^\mu(\tau)+\ell n^\mu(\tau,\theta,\phi),
 \qquad
 n^2=-1,
 \qquad
 n\cdot u=0.
\end{equation*}

在每个 $\tau$ 的瞬时静止系中,

\begin{equation*}
 u^\mu=(c,\vb{0}),
 \qquad
 n^\mu=(0,\sin\theta\cos\phi,\sin\theta\sin\phi,\cos\theta).
\end{equation*}

三个切向量为

\begin{equation*}
 {\partial y^\mu\over\partial\tau}=u^\mu+\ell\dot n^\mu,
 \qquad
 {\partial y^\mu\over\partial\theta}=\ell\partial_\theta n^\mu,
 \qquad
 {\partial y^\mu\over\partial\phi}=\ell\partial_\phi n^\mu.
\end{equation*}

管壁有向面积元定义为(设 $\epsilon_{0123}=+1$)

\begin{equation*}
 \d\Sigma_\nu
 =\epsilon_{\nu\alpha\beta\xi}
 {\partial y^\alpha\over\partial\tau}
 {\partial y^\beta\over\partial\theta}
 {\partial y^\xi\over\partial\phi}
 \d\tau\d\theta\d\phi.
\end{equation*}

由 $n\cdot u=0$ 对固有时微分得

\begin{equation*}
 \dot n\cdot u=-n\cdot a.
\end{equation*}

在瞬时静止系中,切向量矩阵的 $3\times3$ 非零子式为 $\ell^2\sin\theta\,(c+\ell\dot n^0)$:空间部分 $\partial_\theta n^i$、$\partial_\phi n^i$ 张成切平面,其混合积 $|\partial_\theta n\times\partial_\phi n|=\sin\theta$;由上式得 $\dot n^0=-n\cdot a/c$,故行列式为

\begin{equation*}
 \ell^2\sin\theta\left(c+\ell\dot n^0\right)
 =c\ell^2\sin\theta\left(1-{\ell n\cdot a\over c^2}\right).
\end{equation*}

外向空间法矢为 $n^\mu$,但降指标后 $n_i=-n^i$,故

\begin{equation*}
 \d\Sigma_\nu
 =-c\ell^2\left(1-{\ell n\cdot a\over c^2}\right)n_\nu\d\tau\d\Omega
 =-c\ell\left(1-{\xi\cdot a\over c^2}\right)\xi_\nu\d\tau\d\Omega,
\end{equation*}

最后一步用了 $\xi^\mu=\ell n^\mu$。这里出现的几何因子正是近线展开中的 $\Lambda=1-\xi\cdot a/c^2$。

在号差 $+2$ 下,SI 制电磁能动量张量为

\begin{equation*}
 T_{\mathrm{em}}^{\mu\nu}
 ={1\over\mu_0}
 \left(
 -F^{\mu\alpha}F^\nu{}_{\alpha}
 +{1\over4}g^{\mu\nu}F^{\alpha\beta}F_{\alpha\beta}
 \right).
\end{equation*}

管壁上的四动量通量为

\begin{equation*}
 \d P_{\mathrm{wall}}^\mu
 ={1\over c}T_{\mathrm{em}}^{\mu\nu}\d\Sigma_\nu.
\end{equation*}

系数 $1/c$ 来自 SI 制中 $x^0=ct$:四动量的空间分量为 $\int T^{i0}\d^3x/c$,时间分量为 $\int T^{00}\d^3x/c=E/c$。由于面积元为 $O(\ell^2)$,若希望保留积分后发散的 $O(\ell^{-1})$ 项和有限的 $O(1)$ 项,能动量张量中必须保留奇异场平方到 $O(\ell^{-3})$,以及奇异场与正则场的交叉项到 $O(\ell^{-2})$。

为缩短中间代数,定义三个组合

\begin{equation*}
 \mathcal G^{\mu\nu}={\xi^{[\mu}u^{\nu]}\over c},
 \qquad
 \mathcal H^{\mu\nu}={\xi^{[\mu}\dot a^{\nu]}\over c^3},
 \qquad
 \mathcal J^{\mu\nu}={u^{[\mu}a^{\nu]}\over c^3},
\end{equation*}

它们分别承载奇异场展开中的 $\xi u$、$\xi\dot a$ 与 $u a$ 三类反对称结构。上一小节的奇异场展开式使实际场在管壁上写成

\begin{equation*}
 F_{\mathrm{act}}^{\mu\nu}
 =F_{\mathrm{reg}}^{\mu\nu}
 +{q\over4\pi\varepsilon_0c}\Lambda^{-1/2}
 \left[
 A\mathcal G^{\mu\nu}
 +B\mathcal J^{\mu\nu}
 +C\mathcal H^{\mu\nu}
 \right]+O(\ell),
\end{equation*}

其中

\begin{equation*}
 A=\ell^{-3}-{a^2\over8c^4}\ell^{-1},
 \qquad
 B={1+\xi\cdot a/c^2\over2}\ell^{-1},
 \qquad
 C=-{1\over2}\ell^{-1}.
\end{equation*}

计算张量平方所需的乘积先保持未收缩形式。由反对称化定义逐项相乘,并用 $\xi\cdot u=0$、$u^2=c^2$、$u\cdot a=0$、$u\cdot\dot a=-a^2$、$\xi^2=-\ell^2$,得到

\begin{align*}
 \mathcal G^\mu{}_\alpha\mathcal G^{\alpha\nu}
 &=-\xi^\mu\xi^\nu+\ell^2{u^\mu u^\nu\over c^2},\\
 \mathcal G^\mu{}_\alpha\mathcal J^{\alpha\nu}
 &={\xi^\mu a^\nu\over c^2}+{(\xi\cdot a)u^\mu u^\nu\over c^4},\\
 \mathcal J^\mu{}_\alpha\mathcal G^{\alpha\nu}
 &={a^\mu\xi^\nu\over c^2}+{(\xi\cdot a)u^\mu u^\nu\over c^4},\\
 \mathcal G^\mu{}_\alpha\mathcal H^{\alpha\nu}
 &={a^2\xi^\mu\xi^\nu\over c^4}
 +{\ell^2u^\mu\dot a^\nu\over c^4}
 +{(\xi\cdot\dot a)u^\mu\xi^\nu\over c^4},\\
 \mathcal H^\mu{}_\alpha\mathcal G^{\alpha\nu}
 &={a^2\xi^\mu\xi^\nu\over c^4}
 +{\ell^2\dot a^\mu u^\nu\over c^4}
 +{(\xi\cdot\dot a)\xi^\mu u^\nu\over c^4},\\
 \mathcal J^\mu{}_\alpha\mathcal J^{\alpha\nu}
 &=-{a^2u^\mu u^\nu\over c^6}-{a^\mu a^\nu\over c^4}.
\end{align*}

例如第一式逐项展开为

\begin{align*}
 {1\over c^2}(\xi^\mu u_\alpha-u^\mu\xi_\alpha)
 (\xi^\alpha u^\nu-u^\alpha\xi^\nu)
 &={(\xi\cdot u)\xi^\mu u^\nu\over c^2}
 -{u^2\xi^\mu\xi^\nu\over c^2}
 -{\xi^2u^\mu u^\nu\over c^2}
 +{(\xi\cdot u)u^\mu\xi^\nu\over c^2}\\
 &=-\xi^\mu\xi^\nu+{\ell^2u^\mu u^\nu\over c^2};
\end{align*}

其余各式完全同理,例如 $\mathcal G^\mu{}_\alpha\mathcal J^{\alpha\nu}$ 中 $(\xi\cdot u)u^\mu a^\nu/c^4$、$\xi^\mu(u\cdot a)u^\nu/c^4$、$u^\mu(\xi\cdot u)a^\nu/c^4$ 三项分别因 $\xi\cdot u=0$ 与 $u\cdot a=0$ 消失,只留下 $\xi^\mu a^\nu/c^2$ 与 $(\xi\cdot a)u^\mu u^\nu/c^4$。

先求能动量张量第一项中出现的

\begin{equation*}
 \mathcal X^\mu
 \equiv F_{\mathrm{act}}^\mu{}_\alpha
 F_{\mathrm{act}}^{\alpha\nu}\xi_\nu.
\end{equation*}

把管壁上的场展开为奇异场平方与奇异--正则交叉项:

\begin{align*}
 \mathcal X^\mu
 ={}&{q^2\over16\pi^2\varepsilon_0^2c^2\Lambda}
 \bigg[
 A^2\mathcal G^\mu{}_\alpha\mathcal G^{\alpha\nu}
 +AB\left(
 \mathcal G^\mu{}_\alpha\mathcal J^{\alpha\nu}
 +\mathcal J^\mu{}_\alpha\mathcal G^{\alpha\nu}
 \right)\\
 &+AC\left(
 \mathcal G^\mu{}_\alpha\mathcal H^{\alpha\nu}
 +\mathcal H^\mu{}_\alpha\mathcal G^{\alpha\nu}
 \right)
 +B^2\mathcal J^\mu{}_\alpha\mathcal J^{\alpha\nu}
 \bigg]\xi_\nu\\
 &+{q\over4\pi\varepsilon_0c}\ell^{-3}
 \left[
 \mathcal G^\mu{}_\alpha F_{\mathrm{reg}}^{\alpha\nu}
 +F_{\mathrm{reg}}^\mu{}_\alpha\mathcal G^{\alpha\nu}
 \right]\xi_\nu
 +O(1).
\end{align*}

这里 $BH$、$CH$ 和 $H^2$ 至多在面积积分后产生 $O(\ell)$,所以没有写出;交叉项中也只需保留最强的 $A\mathcal G\sim\ell^{-2}$,其 $\Lambda^{-1/2}$ 因子已并入 $\ell^{-3}$ 的 $O(\ell^2)$ 修正中。

把乘积公式与 $\xi_\nu$ 收缩。例如

\begin{align*}
 (\mathcal G\mathcal G)^{\mu\nu}\xi_\nu
 &=\left(-\xi^\mu\xi^\nu+{\ell^2u^\mu u^\nu\over c^2}\right)\xi_\nu
 =-\xi^\mu(\xi\cdot\xi)+{\ell^2u^\mu(\xi\cdot u)\over c^2}
 =\ell^2\xi^\mu,\\
 (\mathcal G\mathcal J)^{\mu\nu}\xi_\nu
 &=\xi^\mu{(a\cdot\xi)\over c^2}+{(\xi\cdot a)u^\mu(u\cdot\xi)\over c^4}
 ={(\xi\cdot a)\xi^\mu\over c^2},
\end{align*}

其余各式用同一套恒等式,完整收缩表为

\begin{align*}
 (\mathcal G\mathcal G)^{\mu\nu}\xi_\nu
 &=\ell^2\xi^\mu,\\
 (\mathcal G\mathcal J)^{\mu\nu}\xi_\nu
 &={(\xi\cdot a)\xi^\mu\over c^2},\\
 (\mathcal J\mathcal G)^{\mu\nu}\xi_\nu
 &=-{\ell^2a^\mu\over c^2},\\
 (\mathcal G\mathcal H)^{\mu\nu}\xi_\nu
 &=-{\ell^2a^2\xi^\mu\over c^4},\\
 (\mathcal H\mathcal G)^{\mu\nu}\xi_\nu
 &=-{\ell^2a^2\xi^\mu\over c^4},\\
 (\mathcal J\mathcal J)^{\mu\nu}\xi_\nu
 &=-{(\xi\cdot a)a^\mu\over c^4}.
\end{align*}

正则场交叉项则为

\begin{align*}
 &\left[
 \mathcal G^\mu{}_\alpha F_{\mathrm{reg}}^{\alpha\nu}
 +F_{\mathrm{reg}}^\mu{}_\alpha\mathcal G^{\alpha\nu}
 \right]\xi_\nu\\
 &\qquad
 ={\ell^2\over c}F_{\mathrm{reg}}^\mu{}_\nu u^\nu
 +{\xi^\mu\over c}F_{\mathrm{reg}\,\alpha\beta}u^\alpha\xi^\beta,
\end{align*}

其中第一项来自 $\mathcal G^{\alpha\nu}\xi_\nu=\ell^2u^\alpha/c$,第二项来自 $\mathcal G^\mu{}_\alpha F_{\mathrm{reg}}^{\alpha\nu}\xi_\nu$,而 $F_{\mathrm{reg}}^{\alpha\nu}\xi_\nu\xi_\alpha=0$(反对称张量收缩于对称张量)。

将 $A,B,C$ 代入,并用到

\begin{align*}
 A^2&=\ell^{-6}-{a^2\over4c^4}\ell^{-4}+O(\ell^{-2}),\\
 AB&={1+\xi\cdot a/c^2\over2}\ell^{-4}+O(\ell^{-2}),\\
 AC&=-{1\over2}\ell^{-4}+O(\ell^{-2}),\\
 B^2&={1\over4}\ell^{-2}+O(\ell^{-1}),
\end{align*}

$\mathcal X^\mu$ 逐项化成

\begin{align*}
 \mathcal X^\mu
 ={q^2\over16\pi^2\varepsilon_0^2c^2\Lambda}
 \bigg\{{}
 &\left(\ell^{-6}-{a^2\over4c^4}\ell^{-4}\right)
 \ell^2\xi^\mu\\
 &+{1+\xi\cdot a/c^2\over2}\ell^{-4}
 \left[{(\xi\cdot a)\xi^\mu\over c^2}
 -{\ell^2a^\mu\over c^2}\right]\\
 &+{a^2\ell^{-2}\xi^\mu\over c^4}
 -{(\xi\cdot a)a^\mu\ell^{-2}\over4c^4}
 \bigg\}\\
 &+{q\over4\pi\varepsilon_0c}\ell^{-3}
 \left[
 {\ell^2\over c}F_{\mathrm{reg}}^\mu{}_\nu u^\nu
 +{\xi^\mu\over c}F_{\mathrm{reg}\,\alpha\beta}u^\alpha\xi^\beta
 \right]+O(1).
\end{align*}

这就是 $F^\mu{}_\alpha F^{\alpha\nu}\xi_\nu$ 的完整所需展开。

再求能动量张量的迹项。所需的反对称张量全收缩为

\begin{align*}
 \mathcal G_{\alpha\beta}\mathcal G^{\alpha\beta}
 &=-2\ell^2,\\
 \mathcal G_{\alpha\beta}\mathcal J^{\alpha\beta}
 &=-{2\xi\cdot a\over c^2},\\
 \mathcal G_{\alpha\beta}\mathcal H^{\alpha\beta}
 &={2\ell^2a^2\over c^4},\\
 \mathcal J_{\alpha\beta}\mathcal J^{\alpha\beta}
 &={2a^2\over c^4}.
\end{align*}

例如

\begin{align*}
 \mathcal G_{\alpha\beta}\mathcal J^{\alpha\beta}
 &={1\over c^4}(\xi_\alpha u_\beta-\xi_\beta u_\alpha)
 (u^\alpha a^\beta-u^\beta a^\alpha)\\
 &={1\over c^4}\left[-(\xi\cdot a)u^2-(u\cdot u)(\xi\cdot a)\right]
 =-{2\xi\cdot a\over c^2},
\end{align*}

其中利用了 $u\cdot a=0$、$u^2=c^2$。因此

\begin{align*}
 F_{\mathrm{act}\,\alpha\beta}F_{\mathrm{act}}^{\alpha\beta}
 ={}&{q^2\over16\pi^2\varepsilon_0^2c^2\Lambda}
 \bigg[
 -2A^2\ell^2
 -4AB{\xi\cdot a\over c^2}
 +4AC{\ell^2a^2\over c^4}
 +2B^2{a^2\over c^4}
 \bigg]\\
 &+{q\over\pi\varepsilon_0c}\ell^{-3}
 {F_{\mathrm{reg}\,\alpha\beta}\xi^\alpha u^\beta\over c}
 +O(1),
\end{align*}

其中 $-4AB(\xi\cdot a)$ 来自交叉项 $2AB\bigl[\mathcal G_{\alpha\beta}\mathcal J^{\alpha\beta}+\mathcal J_{\alpha\beta}\mathcal G^{\alpha\beta}\bigr]=2AB\cdot2(-2\xi\cdot a/c^2)$ 型收缩的两重组合。把 $A,B,C$ 代入并保留所需阶数,得到

\begin{align*}
 F_{\mathrm{act}\,\alpha\beta}F_{\mathrm{act}}^{\alpha\beta}
 ={}&{q^2\over16\pi^2\varepsilon_0^2c^2\Lambda}
 \left[
 -2\left(1+{\xi\cdot a\over c^2}
 +\left({\xi\cdot a\over c^2}\right)^2\right)\ell^{-4}
 -{a^2\ell^{-2}\over c^4}
 \right]\\
 &+{q\over\pi\varepsilon_0c}\ell^{-3}
 {F_{\mathrm{reg}\,\alpha\beta}\xi^\alpha u^\beta\over c}
 +O(1).
\end{align*}

于是迹项为

\begin{align*}
 {1\over4}\xi^\mu
 F_{\mathrm{act}\,\alpha\beta}F_{\mathrm{act}}^{\alpha\beta}
 ={}&{q^2\over16\pi^2\varepsilon_0^2c^2\Lambda}
 \left[
 -{1\over2}
 \left(1+{\xi\cdot a\over c^2}
 +\left({\xi\cdot a\over c^2}\right)^2\right)\ell^{-4}
 -{a^2\ell^{-2}\over4c^4}
 \right]\xi^\mu\\
 &+{q\over4\pi\varepsilon_0c}\ell^{-3}
 {\xi^\mu F_{\mathrm{reg}\,\alpha\beta}\xi^\alpha u^\beta\over c}
 +O(1).
\end{align*}

现在把 $\mathcal X^\mu$ 的展开与迹项相加。由于 $F_{\mathrm{reg}\,\alpha\beta}$ 反对称,

\begin{equation*}
 F_{\mathrm{reg}\,\alpha\beta}u^\alpha\xi^\beta
 =-F_{\mathrm{reg}\,\alpha\beta}\xi^\alpha u^\beta,
\end{equation*}

两个正则场交叉项中与 $\xi^\mu$ 成正比的部分相消,只留下 $F_{\mathrm{reg}}^\mu{}_\nu u^\nu$;$\xi^\mu$ 的系数按 $\ell$ 幂次收集:$\ell^{-4}$ 部分为 $1+\frac12(\xi\cdot a/c^2)+\frac12(\xi\cdot a/c^2)^2-\frac12\bigl(1+\xi\cdot a/c^2+(\xi\cdot a/c^2)^2\bigr)=\frac12$,平方项相消,$\ell^{-2}$ 部分为 $-a^2/(4c^4)+a^2/c^4-a^2/(4c^4)=a^2/(2c^4)$;$a^\mu$ 的系数为 $-(1+\xi\cdot a/c^2)/(2c^2)-(\xi\cdot a/c^2)/(4c^2)=-(1+3\xi\cdot a/2c^2)/(2c^2)$。同一组合化成简洁形式

\begin{align*}
 \mathcal Q^\mu
 &\equiv
 F_{\mathrm{act}}^\mu{}_\alpha
 F_{\mathrm{act}}^{\alpha\nu}\xi_\nu
 +{1\over4}\xi^\mu
 F_{\mathrm{act}\,\alpha\beta}F_{\mathrm{act}}^{\alpha\beta}\\
 &={q^2\over16\pi^2\varepsilon_0^2c^2\Lambda}
 \left[
 \left({1\over2}\ell^{-4}
 +{a^2\ell^{-2}\over2c^4}\right)\xi^\mu
 -{\ell^{-2}\over2c^2}
 \left(1+{3\xi\cdot a\over2c^2}\right)a^\mu
 \right]\\
 &\quad
 +{q\over4\pi\varepsilon_0c^2}\ell^{-1}
 F_{\mathrm{reg}}^\mu{}_\nu u^\nu
 +O(1).
\end{align*}

这一步把看似众多的张量平方化成了三个结构:沿径向的奇函数项、沿加速度的发散惯性项,以及有限的正则 Lorentz 力。

把能动量张量收缩到面积元上:$\d P^\mu=(1/c)T_{\mathrm{em}}^{\mu\nu}\d\Sigma_\nu$ 中,奇异场平方部分恰好按 $\mathcal Q^\mu$ 组装,其中面积元因子 $(1-\xi\cdot a/c^2)$ 与 $\mathcal Q^\mu$ 携带的 $\Lambda^{-1}$ 精确合并为 $1$(因为 $\Lambda=1-\xi\cdot a/c^2$);正则交叉部分合并面积元因子后给出与 $qF_{\mathrm{reg}}^{\mu\nu}u_\nu$ 成正比的项。使用 $\mu_0\varepsilon_0c^2=1$,单位固有时和单位立体角的管壁四动量通量为

\begin{align*}
 {\d^2P_{\mathrm{wall}}^\mu\over\d\tau\d\Omega}
 =-{1\over4\pi}\bigg\{ {q^2\over4\pi\varepsilon_0}
 \bigg[{}
 &\left({1\over2\ell^3}
 +{a^2\over2c^4\ell}\right)\xi^\mu
 -{a^\mu\over2c^2\ell}
 -{3(\xi\cdot a)a^\mu\over4c^4\ell}
 \bigg]\\
 &+\left(1+{\xi\cdot a\over c^2}\right)
 qF_{\mathrm{reg}}^{\mu\nu}u_\nu
 \bigg\}+O(\ell).
\end{align*}

这是同一个通量,但每项的角奇偶性已经显式:方括号内第一、二项为径向奇函数,第三、四项沿加速度方向。

瞬时静止系中的球面对称给出

\begin{equation*}
 \int n^\mu\d\Omega=0,
 \qquad
 \int n^\mu n^\nu\d\Omega
 =-{4\pi\over3}
 \left(g^{\mu\nu}-{u^\mu u^\nu\over c^2}\right).
\end{equation*}

因为 $\xi^\mu=\ell n^\mu$,通量中所有含单个 $\xi^\mu$ 的项积分为零;含 $(\xi\cdot a)a^\mu$ 的项也因 $\int\xi\cdot a\,\d\Omega=0$ 而消失。正则场在 $\ell\to0$ 时与角方向无关。角积分后只剩两项:

\begin{align*}
 \int\left[-{1\over4\pi}\cdot{q^2\over4\pi\varepsilon_0}
 \left(-{a^\mu\over2c^2\ell}\right)\right]\d\Omega
 &={q^2a^\mu\over8\pi\varepsilon_0c^2\ell},\\
 \int\left[-{1\over4\pi}
 \left(1+{\xi\cdot a\over c^2}\right)
 qF_{\mathrm{reg}}^{\mu\nu}u_\nu\right]\d\Omega
 &=-qF_{\mathrm{reg}}^{\mu\nu}u_\nu,
\end{align*}

故同一个管壁通量化成

\begin{equation*}
 {\d P_{\mathrm{wall}}^\mu\over\d\tau}
 ={q^2\over8\pi\varepsilon_0c^2\ell}a^\mu
 -qF_{\mathrm{reg}}^{\mu\nu}u_\nu
 +O(\ell).
\end{equation*}

这就是整个近线展开和世界管积分后的简洁结果:第一项与加速度平行,表现为随管半径发散的局部惯性,其系数 $q^2/(8\pi\varepsilon_0c^2\ell)$ 正是电磁质量;第二项是在世界线上有限的正则 Lorentz 四力。下一步把世界管看成封闭四维区域的边界,把管壁通量与端帽上的动量变化联系起来。

\subsubsection{端帽动量、质量重整化与有限自力}

世界管本身并不是封闭曲面:它有两个端帽。取由管壁以及 $\tau_1$、$\tau_2$ 两个端帽围成的四维区域 $\mathcal V$。真空区域中场能动量的 Gauss 定理为

\begin{equation*}
 \int_{\mathcal V}\partial_\nu T_{\mathrm{em}}^{\mu\nu}\d^4x
 =\int_{\partial\mathcal V}T_{\mathrm{em}}^{\mu\nu}\d\Sigma_\nu
 =\Delta P_{\mathrm{cap}}^\mu+P_{\mathrm{wall}}^\mu.
\end{equation*}

如果只讨论管外的真空区域,$\partial_\nu T_{\mathrm{em}}^{\mu\nu}=0$,于是管壁流出的四动量必须由两个端帽所截获的管内四动量变化补偿。不过,点粒子世界线本身已从积分区域中挖去;Maxwell 方程并不单独规定被挖去部分含有哪些非电磁自由度。因此还必须加入一个点粒子局域性假设:管内未解析自由度可由只依赖当前世界线状态的端帽四动量 $p_0^\mu$ 描述。

若粒子没有额外内部方向,并且在最低导数阶只允许使用 $u^\mu$,则

\begin{equation*}
 p_0^\mu(\ell)=m_0(\ell)u^\mu.
\end{equation*}

这正是 Dirac 把管壁通量视为端点量之差、并选择 $B^\mu\propto u^\mu$ 的物理内容。若允许 $a^\mu$ 或更高导数进入 $p_0^\mu$,便等价于给点粒子增加新的内部结构常数;那将是另一个有效理论,而不是最小的点电子模型。

在无外场时,端帽内粒子动量与管壁流量满足

\begin{equation*}
 {\d p_0^\mu\over\d\tau}
 +{\d P_{\mathrm{wall}}^\mu\over\d\tau}=0.
\end{equation*}

有外部无源场时,正则场已经包含 $F_{\mathrm{in}}^{\mu\nu}$,故把管壁通量公式代入,并使用 ${\d p_0^\mu/\d\tau}=m_0(\ell)a^\mu$,得到预重整化方程

\begin{equation*}
 qF_{\mathrm{reg}}^{\mu\nu}u_\nu
 =\left[
 m_0(\ell)
 +{q^2\over8\pi\varepsilon_0c^2\ell}
 \right]a^\mu+O(\ell).
\end{equation*}

括号内的两个量分别发散,但它们只以同一个加速度系数出现。实验所测得的惯性质量因而定义为这一组合的有限极限:

\begin{equation*}
 m
 =\lim_{\ell\to0}
 \left[
 m_0(\ell)
 +{q^2\over8\pi\varepsilon_0c^2\ell}
 \right].
\end{equation*}

这不是把无限量"丢掉",而是说明裸质量 $m_0(\ell)$ 本身不是可测量参数;只有这一组合进入低能运动方程。

此处的发散系数

\begin{equation*}
 {q^2\over8\pi\varepsilon_0c^2\ell}
 ={1\over c^2}{q^2\over8\pi\varepsilon_0\ell}
 ={U_{\mathrm C}(\ell)\over c^2}
\end{equation*}

等于半径 $\ell$ 的 Coulomb 外场能 $U_{\mathrm C}(\ell)=q^2/(8\pi\varepsilon_0\ell)$ 除以 $c^2$,即球对称近场的 $P_{\mathrm{bound}}^0/c$ 型质量。它来自与世界线正交的协变端帽极限;前面实验室同时面上扩展球电荷的 $4U_{\mathrm{es}}/(3c^2)$ 来自非封闭电磁子系统的动量积分。两者对应不同的积分超曲面和不同的系统划分,不能仅凭"电磁质量"这一名称互相替换。

取预重整化方程的有限极限,得到最紧凑的重整化方程

\begin{equation*}
 ma^\mu=qF_{\mathrm{reg}}^{\mu\nu}u_\nu.
\end{equation*}

这里的 $F_{\mathrm{reg}}$ 仍是 Green 函数一节定义的同一个物理量。把它重新展开为

\begin{equation*}
 F_{\mathrm{reg}}^{\mu\nu}
 =F_{\mathrm{in}}^{\mu\nu}+F_{\mathrm R}^{\mu\nu},
\end{equation*}

并代入前一节的正则场在世界线上的有限极限式 $F_{\mathrm R}^{\mu\nu}(z)=(q/6\pi\varepsilon_0c^5)(\dot a^\mu u^\nu-\dot a^\nu u^\mu)$,有限自场四力为

\begin{align*}
 f_{\mathrm{self}}^\mu
 &\equiv qF_{\mathrm R}^{\mu\nu}u_\nu\\
 &={q^2\over6\pi\varepsilon_0c^5}
 \left(
 \dot a^\mu u^\nu u_\nu
 -u^\mu\dot a^\nu u_\nu
 \right).
\end{align*}

由 $u^2=c^2$ 一次微分得 $u\cdot a=0$,再微分一次得

\begin{equation*}
 a^2+u\cdot\dot a=0,
 \qquad
 u\cdot\dot a=-a^2.
\end{equation*}

因此有限自场四力化成

\begin{equation*}
 f_{\mathrm{self}}^\mu
 ={q^2\over6\pi\varepsilon_0c^3}
 \left(
 \dot a^\mu+{a^2\over c^2}u^\mu
 \right).
\end{equation*}

在号差 $+2$ 下,类时 $u^2=c^2$ 而类空 $a^2<0$。式中第二项必须写成加号,因为

\begin{align*}
 u_\mu f_{\mathrm{self}}^\mu
 &={q^2\over6\pi\varepsilon_0c^3}
 \left(u\cdot\dot a+{a^2\over c^2}u^2\right)\\
 &={q^2\over6\pi\varepsilon_0c^3}(-a^2+a^2)=0.
\end{align*}

这保证自力与四速度正交,不改变 $u^2=c^2$ 的归一化。

把正则场的分解与有限自力代回重整化方程,得到 Abraham--Lorentz--Dirac 方程

\begin{equation*}
 ma^\mu
 =f_{\mathrm{ext}}^\mu
 +{q^2\over6\pi\varepsilon_0c^3}
 \left(
 \dot a^\mu+{a^2\over c^2}u^\mu
 \right),
 \qquad
 f_{\mathrm{ext}}^\mu=qF_{\mathrm{in}}^{\mu\nu}u_\nu.
\end{equation*}

所以正则场在世界线上的有限极限、这里的有限自力和 ALD 方程不是三个独立结论。它们依次是同一个有限对象的场强形式、四力形式和运动方程形式。最后,世界管方法还可以给出同一自作用的第三种等价组织方式:把有限自力改写为束缚近场动量与远区辐射动量的守恒关系。

\subsubsection{束缚动量、Schott 项与远区辐射}

世界管方法还可以把有限自力改写成近场束缚动量和远区辐射动量。端帽内的电磁四动量 $P_{\mathrm{bound}}^\mu(\ell)$ 由奇异场在管内的积分 $(1/c)\int_{r\le\ell}T_{\mathrm{em}}^{\mu0}\d^3x$ 直接展开:Coulomb 部分给出 $q^2u^\mu/(8\pi\varepsilon_0c^2\ell)$(共动系中只有时间分量 $U_{\mathrm C}(\ell)/c$,提升后即与 $u^\mu$ 平行),世界线加速度的线性修正给出 $-q^2a^\mu/(6\pi\varepsilon_0c^3)$,故在小 $\ell$ 下

\begin{equation*}
 P_{\mathrm{bound}}^\mu(\ell)
 ={q^2\over8\pi\varepsilon_0c^2\ell}u^\mu
 -{q^2\over6\pi\varepsilon_0c^3}a^\mu
 +O(\ell).
\end{equation*}

下面两个性质将验证这一展开:第一项的导数复现管壁通量的发散部分,而总守恒式回到 ALD 方程。第一项是 Coulomb 近场携带的发散四动量,它的导数为

\begin{equation*}
 {\d\over\d\tau}
 \left({q^2\over8\pi\varepsilon_0c^2\ell}u^\mu\right)
 ={q^2\over8\pi\varepsilon_0c^2\ell}a^\mu,
\end{equation*}

正好复现管壁通量中的发散惯性。有限部分定义为 Schott 四动量

\begin{equation*}
 P_{\mathrm{Schott}}^\mu
 =-{q^2\over6\pi\varepsilon_0c^3}a^\mu.
\end{equation*}

它依赖瞬时加速度,是附着在粒子附近的可逆束缚场动量,而不是已经传播到无穷远的辐射动量。

远区 Lienard 辐射的四动量率必须与 $u^\mu$ 平行,因为在瞬时静止系中辐射的净三动量为零,而能量为正。协变形式为

\begin{equation*}
 {\d P_{\mathrm{rad}}^\mu\over\d\tau}
 =-{q^2\over6\pi\varepsilon_0c^5}a^2u^\mu.
\end{equation*}

由于 $a^2<0$,右侧时间分量为正。

把裸粒子动量、束缚场动量和已辐射动量放在同一个守恒式中:

\begin{equation*}
 {\d\over\d\tau}
 \left[m_0(\ell)u^\mu
 +P_{\mathrm{bound}}^\mu(\ell)\right]
 +{\d P_{\mathrm{rad}}^\mu\over\d\tau}
 =f_{\mathrm{ext}}^\mu.
\end{equation*}

逐项代入束缚动量和辐射四动量率的展开:

\begin{align*}
 &\left[m_0(\ell)
 +{q^2\over8\pi\varepsilon_0c^2\ell}\right]a^\mu
 -{q^2\over6\pi\varepsilon_0c^3}\dot a^\mu
 -{q^2\over6\pi\varepsilon_0c^5}a^2u^\mu
 +O(\ell)
 =f_{\mathrm{ext}}^\mu.
\end{align*}

使用质量重整化定义并把后两项移到右边,立即回到 ALD 方程。因此 Green 函数半差法、世界管壁通量法以及"裸粒子加束缚近场加远区辐射"的守恒法,是同一个自作用问题的三种等价组织方式。

具体地,由 Schott 四动量与远区辐射四动量率,

\begin{align*}
 -{\d P_{\mathrm{Schott}}^\mu\over\d\tau}
 -{\d P_{\mathrm{rad}}^\mu\over\d\tau}
 &={q^2\over6\pi\varepsilon_0c^3}\dot a^\mu
 +{q^2\over6\pi\varepsilon_0c^5}a^2u^\mu\\
 &={q^2\over6\pi\varepsilon_0c^3}
 \left(\dot a^\mu+{a^2\over c^2}u^\mu\right)
 =f_{\mathrm{self}}^\mu.
\end{align*}

第一行把自力看成两个四动量流的负和,第二行把同一个量化成局域世界线导数。

辐射四动量率的时间分量给出 Lienard 功率。四加速度不变量为

\begin{equation*}
 -a^2
 =\gamma^6
 \left[
 \vb a^2-{(\vb v\times\vb a)^2\over c^2}
 \right].
\end{equation*}

这一恒等式来自 $a^\mu$ 的三维分解:$a^0=\gamma^4\vb v\cdot\vb a/c$、$\vb a_{(4)}=\gamma^2\vb a+\gamma^4(\vb v\cdot\vb a)\vb v/c^2$,代入 $-a^2=(a^0)^2-\vb a_{(4)}^2$ 并合并同类项即得。故

\begin{equation*}
 P_{\mathrm{rad}}
 ={q^2\gamma^6\over6\pi\varepsilon_0c^3}
 \left[
 \vb a^2-{(\vb v\times\vb a)^2\over c^2}
 \right].
\end{equation*}

确实,由 $P^0=E/c$ 与 $u^0=\gamma c$,功率 $P=\d E/\d t=(1/\gamma)\d E/\d\tau=(c/\gamma)\d P_{\mathrm{rad}}^0/\d\tau=-q^2a^2/(6\pi\varepsilon_0c^3)$,代入上式一致。

在非相对论极限,定义辐射反作用时间

\begin{equation*}
 \tau_{\mathrm{rr}}
 ={q^2\over6\pi\varepsilon_0mc^3},
 \qquad
 m\tau_{\mathrm{rr}}={q^2\over6\pi\varepsilon_0c^3}.
\end{equation*}

则

\begin{equation*}
 \vb F_{\mathrm{AL}}=m\tau_{\mathrm{rr}}\dot{\vb a},
 \qquad
 P_{\mathrm{rad}}=m\tau_{\mathrm{rr}}\vb a^2.
\end{equation*}

机械功率由同一个自力给出

\begin{align*}
 \vb F_{\mathrm{AL}}\cdot\vb v
 &=m\tau_{\mathrm{rr}}\dot{\vb a}\cdot\vb v\\
 &=m\tau_{\mathrm{rr}}{\d\over\d t}(\vb a\cdot\vb v)
 -m\tau_{\mathrm{rr}}\vb a^2\\
 &=-{\d E_{\mathrm{Schott}}\over\d t}-P_{\mathrm{rad}},
 \qquad
 E_{\mathrm{Schott}}=-m\tau_{\mathrm{rr}}\vb a\cdot\vb v.
\end{align*}

所以瞬时机械功率并不总等于负的远区辐射功率;二者之差是束缚近场能量的时间导数。只有当散射过程初末 $\vb a\to0$ 时,Schott 边界项才消失,机械总能量损失才等于远区全频谱积分。

作为检验,考虑一维恒定固有加速度 $\alpha$ 的双曲运动。此时

\begin{equation*}
 a^2=-\alpha^2,
 \qquad
 \dot a^\mu={\alpha^2\over c^2}u^\mu
 =-{a^2\over c^2}u^\mu.
\end{equation*}

代入有限自力得

\begin{equation*}
 f_{\mathrm{self}}^\mu=0,
\end{equation*}

而 Lienard 功率仍给出非零辐射。由前面的链式关系,此时

\begin{equation*}
 {\d P_{\mathrm{Schott}}^\mu\over\d\tau}
 =-{\d P_{\mathrm{rad}}^\mu\over\d\tau},
\end{equation*}

即远区辐射四动量完全由束缚近场四动量的减少供给;这不是能量守恒的矛盾,而是瞬时场能分组的结果。

把 ALD 方程写成三维形式还可检验低速极限。四力与三力满足

\begin{equation*}
 f^\mu=\gamma\left({\vb F\cdot\vb v\over c},\vb F\right),
\end{equation*}

故三力是四力空间分量除以 $\gamma$。由 $u^\mu=\gamma(c,\vb v)$ 连续求导,

\begin{align*}
 a^0&=\gamma^4{\vb v\cdot\vb a\over c},\\
 \vb a_{(4)}&=\gamma^2\vb a
 +\gamma^4{\vb v\cdot\vb a\over c^2}\vb v,
\end{align*}

再对固有时求导得四 jerk 的各个分量,连同 $a^2u^\mu/c^2$ 项代入有限自力的空间分量 $\vb f_{\mathrm{self}}=(q^2/6\pi\varepsilon_0c^3)(\dot{\vb a}_{(4)}+a^2\gamma\vb v/c^2)$,除以 $\gamma$ 并合并同类项,得到相对论三维 ALD 力

\begin{align*}
 \vb F_{\mathrm{ALD}}
 ={q^2\over6\pi\varepsilon_0c^3}
 \bigg[{}
 &\gamma^2\dot{\vb a}
 +{3\gamma^4\over c^2}(\vb v\cdot\vb a)\vb a
 +{\gamma^4\over c^2}(\vb v\cdot\dot{\vb a})\vb v\\
 &+{3\gamma^6\over c^4}(\vb v\cdot\vb a)^2\vb v
 \bigg].
\end{align*}

当 $v/c\ll1$ 时,$\gamma=1+O(v^2/c^2)$,后三项均比第一项高两个 $1/c$ 阶,故

\begin{equation*}
 \vb F_{\mathrm{AL}}
 ={q^2\over6\pi\varepsilon_0c^3}\dot{\vb a}
 +O(c^{-5}).
\end{equation*}

因此非相对论 Abraham--Lorentz 力不是另一个独立模型,而是 Dirac 正则场自力在低速极限下的同一表达。

\subsubsection{Abraham--Lorentz--Dirac 方程的可解应用、无 runaway 条件与降阶}

前面已经给出点电荷的完整协变方程。为看清这个三阶世界线方程的解空间，先不作低速近似，而考虑严格的一维运动：用快度 $\psi(\tau)$ 参数化四速度，
\begin{equation*}
 u^\mu=c\bigl(\cosh\psi,\sinh\psi,0,0\bigr),
 \qquad
 v=c\tanh\psi,
 \qquad
 p=mc\sinh\psi.
\end{equation*}
对固有时逐次求导，得到
\begin{align*}
 a^\mu
 &=c\dot\psi\bigl(\sinh\psi,\cosh\psi,0,0\bigr),\\
 \dot a^\mu
 &=c\ddot\psi\bigl(\sinh\psi,\cosh\psi,0,0\bigr)
 +c\dot\psi^2\bigl(\cosh\psi,\sinh\psi,0,0\bigr),\\
 a^2&=-c^2\dot\psi^2.
\end{align*}
其中第三式是 $a^\mu$ 自缩并的结果，用到 $\sinh^2\psi-\cosh^2\psi=-1$。把最后一式代入 ALD 组合，含 $\dot\psi^2$ 的两项逐分量抵消：
\begin{align*}
 \dot a^\mu+{a^2\over c^2}u^\mu
 &={}c\ddot\psi(\sinh\psi,\cosh\psi,0,0)
 +c\dot\psi^2(\cosh\psi,\sinh\psi,0,0)
 -c\dot\psi^2(\cosh\psi,\sinh\psi,0,0)\\
 &=c\ddot\psi(\sinh\psi,\cosh\psi,0,0).
\end{align*}
令 $F_\parallel(\tau)$ 表示实验室中的一维普通力，即 $F_\parallel=\d p/\d t$。由 $v=c\tanh\psi$ 解出 $\gamma=(1-\tanh^2\psi)^{-1/2}=\cosh\psi$、$\gamma v/c=\sinh\psi$，相应四力为
\begin{align*}
 f_{\mathrm{ext}}^\mu
 &=\gamma\left({F_\parallel v\over c},F_\parallel,0,0\right)\\
 &=F_\parallel(\sinh\psi,\cosh\psi,0,0).
\end{align*}
以上三个式子中，惯性项、外力项和自力项都平行于同一个单位类空方向 $(\sinh\psi,\cosh\psi,0,0)$。快度参数化的威力正在于此：三条平行方向使四维向量方程退化为单个标量方程，因此四维 ALD 方程严格化为
\begin{equation*}
 mc\left(\dot\psi-\tau_{\mathrm{rr}}\ddot\psi\right)
 =F_\parallel(\tau),
 \qquad
 \tau_{\mathrm{rr}}={q^2\over6\pi\varepsilon_0mc^3}.
\end{equation*}
这一化简没有使用 $v/c\ll1$：把 $a^\mu=c\dot\psi(\sinh\psi,\cosh\psi,0,0)$、$\dot a^\mu+a^2u^\mu/c^2=c\ddot\psi(\sinh\psi,\cosh\psi,0,0)$ 与四力代入 $ma^\mu=f_{\mathrm{ext}}^\mu+m\tau_{\mathrm{rr}}(\dot a^\mu+a^2u^\mu/c^2)$，三个平行方向上的系数 $mc\dot\psi$、$F_\parallel$、$mc\tau_{\mathrm{rr}}\ddot\psi$ 相等即得上式。非线性的相对论运动学全部被快度吸收，而辐射反作用只作用于 $\dot\psi$。

令 $w(\tau)=\dot\psi(\tau)$，标量方程改写为
\begin{equation*}
 \dot w-{1\over\tau_{\mathrm{rr}}}w
 =-{F_\parallel(\tau)\over mc\tau_{\mathrm{rr}}}.
\end{equation*}
乘以积分因子 $e^{-\tau/\tau_{\mathrm{rr}}}$，得到
\begin{equation*}
 {\d\over\d\tau}
 \left(e^{-\tau/\tau_{\mathrm{rr}}}w\right)
 =-{e^{-\tau/\tau_{\mathrm{rr}}}\over mc\tau_{\mathrm{rr}}}
 F_\parallel(\tau).
\end{equation*}
从任意 $\tau_0$ 积分到 $\tau$，一般解为
\begin{equation*}
 w(\tau)
 =e^{(\tau-\tau_0)/\tau_{\mathrm{rr}}}
 \bigg[
 w(\tau_0)
 -{1\over mc\tau_{\mathrm{rr}}}
 \int_{\tau_0}^{\tau}
 e^{-(s-\tau_0)/\tau_{\mathrm{rr}}}F_\parallel(s)\d s
 \bigg].
\end{equation*}
方括号中的常数不受外力历史约束；若它非零，$w$ 含有 $e^{\tau/\tau_{\mathrm{rr}}}$ 的 runaway 模态。要求外力在未来关闭以后加速度不指数增长，相当于施加未来边界条件
\begin{equation*}
 \lim_{\tau\to+\infty}
 e^{-\tau/\tau_{\mathrm{rr}}}w(\tau)=0.
\end{equation*}
将积分因子关系式从 $\tau$ 积分到 $+\infty$：$\left[e^{-s/\tau_{\mathrm{rr}}}w(s)\right]_{\tau}^{\infty}=-\frac{1}{mc\tau_{\mathrm{rr}}}\int_{\tau}^{\infty}e^{-s/\tau_{\mathrm{rr}}}F_\parallel(s)\d s$，未来边界条件使上端为零，再令 $r=(s-\tau)/\tau_{\mathrm{rr}}$、$\d s=\tau_{\mathrm{rr}}\d r$，得到唯一的无 runaway 解
\begin{align*}
 \dot\psi(\tau)
 &={1\over mc\tau_{\mathrm{rr}}}
 \int_{\tau}^{\infty}
 e^{-(s-\tau)/\tau_{\mathrm{rr}}}F_\parallel(s)\d s\\
 &={1\over mc}
 \int_0^\infty e^{-r}
 F_\parallel(\tau+\tau_{\mathrm{rr}}r)\d r.
\end{align*}
同一个解的快度和动量为
\begin{equation*}
 \psi(\tau)=\psi(\tau_i)+\int_{\tau_i}^{\tau}\dot\psi(s)\d s,
 \qquad
 p(\tau)=mc\sinh\psi(\tau).
\end{equation*}
无 runaway 解的积分形式清楚显示：消去 runaway 的代价是当前加速度含有未来一个宽度约为 $\tau_{\mathrm{rr}}$ 的外力平均，这就是预加速的数学来源。

取最简单的阶跃力
\begin{equation*}
 F_\parallel(\tau)=F_0\Theta(\tau).
\end{equation*}
直接计算无 runaway 积分：对 $\tau<0$，阶跃在 $r=-\tau/\tau_{\mathrm{rr}}$ 处开启，$\int_{-\tau/\tau_{\mathrm{rr}}}^{\infty}e^{-r}\d r=e^{\tau/\tau_{\mathrm{rr}}}$；对 $\tau\geq0$ 全区间积分得 $1$。于是
\begin{equation*}
 \dot\psi(\tau)
 ={F_0\over mc}
 \begin{cases}
 e^{\tau/\tau_{\mathrm{rr}}},&\tau<0,\\
 1,&\tau\geq0.
 \end{cases}
\end{equation*}
若 $\psi(-\infty)=\psi_-$，再积分一次得到闭式
\begin{equation*}
 \psi(\tau)=\psi_-+{F_0\over mc}
 \begin{cases}
 \tau_{\mathrm{rr}}e^{\tau/\tau_{\mathrm{rr}}},&\tau<0,\\
 \tau+\tau_{\mathrm{rr}},&\tau\geq0,
 \end{cases}
 \qquad
 p(\tau)=mc\sinh\psi(\tau).
\end{equation*}
速度和快度在开关时刻连续，但加速度在外力开启以前已经按 $e^{\tau/\tau_{\mathrm{rr}}}$ 建立，这就是预加速。

为把“开启”和“关闭”放在同一个有限过程内，考虑矩形脉冲
\begin{equation*}
 F_\parallel(\tau)=F_0\bigl[\Theta(\tau)-\Theta(\tau-T)\bigr],
 \qquad T>0.
\end{equation*}
同一个无 runaway 积分给出：对 $0\leq\tau\leq T$ 只有前一个阶跃开启，$\int_0^{(T-\tau)/\tau_{\mathrm{rr}}}e^{-r}\d r=1-e^{-(T-\tau)/\tau_{\mathrm{rr}}}$；对 $\tau<0$ 积分区间为 $(-\tau/\tau_{\mathrm{rr}},(T-\tau)/\tau_{\mathrm{rr}})$，$\int e^{-r}\d r=e^{\tau/\tau_{\mathrm{rr}}}\left(1-e^{-T/\tau_{\mathrm{rr}}}\right)$；$\tau>T$ 时两个阶跃同时开启，积分为零。于是
\begin{equation*}
 \dot\psi(\tau)={F_0\over mc}
 \begin{cases}
 e^{\tau/\tau_{\mathrm{rr}}}
 \left(1-e^{-T/\tau_{\mathrm{rr}}}\right),&\tau<0,\\[2mm]
 1-e^{-(T-\tau)/\tau_{\mathrm{rr}}},&0\leq\tau\leq T,\\[2mm]
 0,&\tau>T.
 \end{cases}
\end{equation*}
逐段积分并要求 $\psi$ 连续，得到
\begin{equation*}
 \psi(\tau)=\psi_-+{F_0\over mc}
 \begin{cases}
 \tau_{\mathrm{rr}}e^{\tau/\tau_{\mathrm{rr}}}
 (1-e^{-T/\tau_{\mathrm{rr}}}),&\tau<0,\\[1mm]
 \tau+\tau_{\mathrm{rr}}
 \left[1-e^{-(T-\tau)/\tau_{\mathrm{rr}}}\right],&0\leq\tau\leq T,\\[1mm]
 T,&\tau>T.
 \end{cases}
\end{equation*}
粒子在脉冲到来前预加速，并在脉冲关闭前提前减小加速度；脉冲结束后没有 runaway，最终快度增量恰为 $F_0T/(mc)$。

若 $F_\parallel$ 在尺度 $T$ 上光滑变化且 $T\gg\tau_{\mathrm{rr}}$，可在无 runaway 积分中作 Taylor 展开：
\begin{align*}
 F_\parallel(\tau+\tau_{\mathrm{rr}}r)
 &=\sum_{n=0}^\infty
 {\tau_{\mathrm{rr}}^nr^n\over n!}
 {\d^nF_\parallel\over\d\tau^n},\\
 \dot\psi(\tau)
 &= {1\over mc}
 \sum_{n=0}^\infty
 {\tau_{\mathrm{rr}}^n\over n!}
 {\d^nF_\parallel\over\d\tau^n}
 \int_0^\infty e^{-r}r^n\d r\\
 &= {1\over mc}
 \sum_{n=0}^\infty
 \tau_{\mathrm{rr}}^n
 {\d^nF_\parallel\over\d\tau^n}.
\end{align*}
这里使用了 $\int_0^\infty e^{-r}r^n\d r=n!$。保留到一阶，
\begin{equation*}
 mc\dot\psi
 =F_\parallel+\tau_{\mathrm{rr}}\dot F_\parallel
 +O\!\left({\tau_{\mathrm{rr}}^2\over T^2}\right).
\end{equation*}
降阶式是精确无 runaway 解的低频渐近，而不是与 ALD 方程并列的另一条基本定律。对阶跃力，$\dot F_\parallel=F_0\delta(\tau)$，局域截断把平滑预加速层替换成一个 $\delta$ 冲量；这说明导数展开不能在不光滑外力处逐项使用。

上述降阶可直接写成协变形式。定义与四速度正交的投影算符
\begin{equation*}
 \Delta^\mu{}_{\nu}
 =\delta^\mu{}_{\nu}-{u^\mu u_\nu\over c^2}.
\end{equation*}
由 $u^2=c^2$ 求导得 $u\cdot a=0$，再对 $u\cdot a=0$ 求导得 $u\cdot\dot a=-a^2$，
\begin{align*}
 \Delta^\mu{}_{\nu}\dot a^\nu
 &=\dot a^\mu-{u^\mu\over c^2}(u\cdot\dot a)\\
 &=\dot a^\mu+{a^2\over c^2}u^\mu.
\end{align*}
因此 ALD 方程等价地写成
\begin{equation*}
 ma^\mu=f_{\mathrm{ext}}^\mu
 +m\tau_{\mathrm{rr}}\Delta^\mu{}_{\nu}\dot a^\nu.
\end{equation*}
在自力项中使用零阶关系 $ma^\mu=f_{\mathrm{ext}}^\mu+O(\tau_{\mathrm{rr}})$：对该关系沿世界线求导得 $\dot a^\nu=\frac{1}{m}\d f_{\mathrm{ext}}^\nu/\d\tau+O(\tau_{\mathrm{rr}})$，代入 $m\tau_{\mathrm{rr}}\Delta^\mu{}_\nu\dot a^\nu$，得到一阶降阶式
\begin{equation*}
 ma^\mu=f_{\mathrm{ext}}^\mu
 +\tau_{\mathrm{rr}}\Delta^\mu{}_{\nu}
 {\d f_{\mathrm{ext}}^\nu\over\d\tau}
 +O(\tau_{\mathrm{rr}}^2).
\end{equation*}
对 Lorentz 四力 $f_{\mathrm{ext}}^\mu=qF_{\mathrm{ext}}^\mu{}_{\nu}u^\nu$，先求导
\begin{equation*}
 {\d f_{\mathrm{ext}}^\mu\over\d\tau}
 =q(\partial_\alpha F_{\mathrm{ext}}^\mu{}_{\nu})u^\alpha u^\nu
 +qF_{\mathrm{ext}}^\mu{}_{\nu}a^\nu,
\end{equation*}
把该导数代入降阶式中的 $\tau_{\mathrm{rr}}\Delta^\mu{}_\nu\dot f_{\mathrm{ext}}^\nu$，第一项（$\partial_\alpha F$ 部分）保留，第二项（$F a$ 部分）用零阶关系 $a^\nu=(q/m)F_{\mathrm{ext}}^\nu{}_{\beta}u^\beta+O(\tau_{\mathrm{rr}})$ 替换，一般降阶式化为 Landau--Lifshitz 方程
\begin{align*}
 ma^\mu={}&qF_{\mathrm{ext}}^\mu{}_{\nu}u^\nu
 +q\tau_{\mathrm{rr}}
 \Delta^\mu{}_{\rho}
 (\partial_\alpha F_{\mathrm{ext}}^\rho{}_{\nu})u^\alpha u^\nu\\
 &+{q^2\tau_{\mathrm{rr}}\over m}
 \Delta^\mu{}_{\rho}
 F_{\mathrm{ext}}^\rho{}_{\nu}
 F_{\mathrm{ext}}^\nu{}_{\beta}u^\beta
 +O(\tau_{\mathrm{rr}}^2).
\end{align*}
投影算符使右端自动与 $u^\mu$ 正交。若用普通三力 $\vb F=\d\vb p/\d t$ 表示，并把自力项中的加速度替换为零阶加速度 $\vb a_0$，一般降阶式的空间部分可化为三维形式。记 $f^\mu=\gamma(\vb F\cdot\vb v/c,\vb F)$。因 $u_\mu f^\mu=0$，投影可改写成 $\Delta^{\mu}{}_{\nu}\dot f^\nu=\dot f^\mu+(a\cdot f/c^2)u^\mu$；由 $u^\mu=\gamma(c,\vb v)$ 与 $\d\gamma/\d t=\gamma^3(\vb v\cdot\vb a_0)/c^2$ 直接计算得 $a\cdot f=-\gamma^3\vb a_0\cdot\vb F$、$\d f^i/\d\tau=\gamma\dot\gamma F^i+\gamma^2\d F^i/\d t$。把 $ma^i=\gamma\d p^i/\d t$ 的空间分量除以 $\gamma$ 并代入 $\dot\gamma$，
\begin{align*}
 {\d p^i\over\d t}
 &=F^i+\tau_{\mathrm{rr}}\left[\gamma{\d F^i\over\d t}
 +{\gamma^3\over c^2}(\vb v\cdot\vb a_0)F^i
 -{\gamma^3\over c^2}v^i(\vb a_0\cdot\vb F)\right]\\
 &=F^i+\tau_{\mathrm{rr}}\left[\gamma{\d F^i\over\d t}
 -{\gamma^3\over c^2}\left(\vb a_0\times(\vb v\times\vb F)\right)^i\right],
\end{align*}
末行用到恒等式 $(\vb v\cdot\vb a_0)\vb F-\vb v(\vb a_0\cdot\vb F)=-\vb a_0\times(\vb v\times\vb F)$。写成向量形式：
\begin{equation*}
 {\d\vb p\over\d t}
 =\vb F
 +\tau_{\mathrm{rr}}
 \left[
 \gamma{\d\vb F\over\d t}
 -{\gamma^3\over c^2}
 \vb a_0\times(\vb v\times\vb F)
 \right]
 +O(\tau_{\mathrm{rr}}^2),
\end{equation*}
其中 $\d/\d t$ 是沿零阶轨道的全导数。这三个公式是同一降阶动力学的抽象投影形式、外电磁场形式和三维实验室形式。

恒定一维普通力 $F_\parallel=F_0$ 是一个有用检验。标量方程的无 runaway 解为 $\dot\psi=F_0/(mc)$，故 $\ddot\psi=0$，ALD 自力严格为零；与此同时 Lienard 功率非零。前面关于 Schott 项与远区辐射的守恒链式关系已经说明，这一辐射由 Schott 束缚动量的变化补偿。快度形式把“恒定固有加速度下有辐射而局域自力为零”的关系化成一行可检验的等式。这一节始终在点模型内部求解；下面保留源的有限形状，看同一个自作用理论如何给出附加惯性与延迟自力。

\subsubsection{有限尺寸正则化:附加惯性与延迟自力}

Dirac 计算说明，点粒子推迟场的共同奇异部分只通过质量重整化进入可观测运动方程，而推迟场与超前场的半差留下有限的 ALD 自力。现在反过来把点流 $J^\mu_{\mathrm{point}}$ 平滑为有限尺寸守恒电流。这样做不是另起一种理论，而是在同一个四维 Green 函数和同一个能动量守恒律中保留源的空间形状。附加质量、频率依赖惯性和延迟自力将作为这一平滑理论的低频、低速表示出现。

对任意光滑守恒电流，电磁能动量张量仍满足
\begin{equation*}
 T_{\mathrm{em}}^{\mu\nu}
 ={1\over\mu_0}
 \left(
 -F^{\mu\alpha}F^\nu{}_{\alpha}
 +{1\over4}g^{\mu\nu}F^{\alpha\beta}F_{\alpha\beta}
 \right),
 \qquad
 \partial_\mu T_{\mathrm{em}}^{\mu\nu}
 =-F^{\nu\lambda}J_\lambda.
\end{equation*}
这一关系由麦克斯韦方程直接推出：对张量逐项微分，并用 $\partial_\mu F^{\mu\alpha}=\mu_0J^\alpha$ 与 Bianchi 恒等式，
\begin{align*}
 \partial_\mu T_{\mathrm{em}}^{\mu\nu}
 &={1\over\mu_0}\left[-\partial_\mu F^{\mu\alpha}F^\nu{}_\alpha
 -F^{\mu\alpha}\partial_\mu F^\nu{}_\alpha
 +{1\over4}\partial^\nu(F^{\alpha\beta}F_{\alpha\beta})\right]\\
 &=-J^\alpha F^\nu{}_\alpha
 -{1\over2}F^{\mu\alpha}\left(\partial_\mu F^\nu{}_\alpha
 -\partial_\alpha F^\nu{}_\mu\right)
 +{1\over2}F^{\mu\alpha}\partial^\nu F_{\mu\alpha}\\
 &=-F^{\nu\lambda}J_\lambda,
\end{align*}
其中第二步第二项把导数换成反对称组合，利用了 $F^{\mu\alpha}$ 的反对称性；Bianchi 恒等式给出 $\partial_\mu F^\nu{}_\alpha-\partial_\alpha F^\nu{}_\mu=\partial^\nu F_{\mu\alpha}$，使末两项相消。
令 $T_{\mathrm{mat}}^{\mu\nu}$ 包含带电物质的惯性、维持其结构的约束应力及其他非电磁相互作用，并把总场写成
\begin{equation*}
 F^{\mu\nu}=F_{\mathrm{ext}}^{\mu\nu}+F_{\mathrm{self}}^{\mu\nu}.
\end{equation*}
物质与自场之间交换的四动量在总和中相消，因此局域守恒律可连续写成
\begin{align*}
 \partial_\mu T_{\mathrm{mat}}^{\mu\nu}
 &=F_{\mathrm{ext}}^{\nu\lambda}J_\lambda
   +F_{\mathrm{self}}^{\nu\lambda}J_\lambda,\\
 \partial_\mu
 \left(T_{\mathrm{mat}}^{\mu\nu}+T_{\mathrm{self}}^{\mu\nu}\right)
 &=F_{\mathrm{ext}}^{\nu\lambda}J_\lambda.
\end{align*}
第二行不是新的假设，而是第一行与自场能动量张量方程相加后的同一个守恒关系。取包围带电体的世界管 $\mathcal W$，以两张类空超曲面 $\Sigma_1,\Sigma_2$ 截断，把守恒关系在 $\Omega$ 上积分。Gauss 定理把 $\partial_\mu(T_{\mathrm{mat}}+T_{\mathrm{self}})^{\mu\nu}$ 的体积分化为 $\Sigma_2$、$-\Sigma_1$ 与管壁三部分的边界积分；物质场不穿过管壁，管壁项只剩自场部分。按 $P^\nu[\Sigma]$ 的定义整理，得到
\begin{align*}
 P_{\mathrm{mat}}^\nu[\Sigma_2]-P_{\mathrm{mat}}^\nu[\Sigma_1]
 +P_{\mathrm{self}}^\nu[\Sigma_2]-P_{\mathrm{self}}^\nu[\Sigma_1]
 +\int_{\partial\mathcal W}T_{\mathrm{self}}^{\mu\nu}\d\Sigma_\mu
 &=\int_\Omega F_{\mathrm{ext}}^{\nu\lambda}J_\lambda\d^4x,\\
 P^\nu[\Sigma]&={1\over c}\int_\Sigma T^{\mu\nu}\d\Sigma_\mu.
\end{align*}
世界管取得很大时，侧面通量是远区辐射四动量；世界管缩到源附近时，侧面和端面还含有随源一起运动的束缚场四动量。后面出现的附加惯性、Schott 动量和辐射反作用，都是同一自场能动量在不同尺度下的表示。

自场由推迟 Green 函数确定。在 Lorenz 规范下，
\begin{equation*}
 \Box D_{\mathrm{ret}}(x-x')=\delta^{(4)}(x-x'),
 \qquad
 A_{\mathrm{self}}^\mu(x)
 =\mu_0\int D_{\mathrm{ret}}(x-x')J^\mu(x')\d^4x',
\end{equation*}
故
\begin{equation*}
 F_{\mathrm{self}}^{\mu\nu}(x)
 =\mu_0\int
 \left[
 \partial^\mu D_{\mathrm{ret}}(x-x')J^\nu(x')
 -\partial^\nu D_{\mathrm{ret}}(x-x')J^\mu(x')
 \right]\d^4x'.
\end{equation*}
将自场代回 Lorentz 力密度，同一个自作用首先写成最一般的电流泛函
\begin{equation*}
 f_{\mathrm{self}}^\mu(x)
 =\mu_0J_\nu(x)\int
 \left[
 \partial^\mu D_{\mathrm{ret}}(x-x')J^\nu(x')
 -\partial^\nu D_{\mathrm{ret}}(x-x')J^\mu(x')
 \right]\d^4x'.
\end{equation*}
对于有限尺寸的光滑电流，该泛函有限且因果，但它依赖源的整个历史。若源可由若干集体坐标 $Q^A(t)$ 描述，在线性近似下将电流展开为
\begin{equation*}
 J^\mu(x;Q)=J_0^\mu(x)+J_A^\mu(x)Q^A+O(Q^2),
\end{equation*}
并把自力泛函投影到 $Q^A$。线性化后的自作用泛函是 $Q$ 的线性、平移不变泛函，因果性（$t$ 时刻的力只依赖 $t'<t$ 的历史）使这种泛函必为卷积形式，同一个自作用便改写为推迟记忆核
\begin{equation*}
 \mathcal F_A^{\mathrm{self}}(t)
 =\int_{-\infty}^tK_{AB}(t-t')Q^B(t')\d t',
 \qquad
 \widetilde{\mathcal F}_A^{\mathrm{self}}(\omega)
 =K_{AB}(\omega)\widetilde Q^B(\omega).
\end{equation*}
因 $K_{AB}(t<0)=0$，$K_{AB}(\omega)$ 在上半复频率平面解析，并满足
\begin{equation*}
 K_{AB}(-\omega)=K_{AB}(\omega)^*.
\end{equation*}
这两条性质分别表达因果性和实值性。

现取最重要的集体坐标，即整个电荷分布的平移 $\vb X(t)$。均匀位移不能产生自恢复力，均匀速度也不能产生耗散，因此平移核在 $\omega=0$ 至少有二阶零点。于是同一个平移自力可以依次写成
\begin{equation*}
 \widetilde F_{\mathrm{self},i}
 =K_{ij}(\omega)\widetilde X_j
 =\omega^2\mathcal M_{ij}^{\mathrm{em}}(\omega)\widetilde X_j
 =-\mathcal M_{ij}^{\mathrm{em}}(\omega)\widetilde a_j.
\end{equation*}
其中第二种与第三种写法由 $\vb a=\ddot{\vb X}$ 的频域关系 $\widetilde a_j=-\omega^2\widetilde X_j$ 等价。三种写法描述的是同一个物理量：第一种强调它是位移的线性响应，第二种抽出平移零模，第三种把剩余响应解释为频率依赖的附加惯性张量。一般形状的带电体只有张量 $\mathcal M_{ij}^{\mathrm{em}}$；只有施加旋转对称性以后才出现单一的“附加质量”。

下面由 Maxwell 方程直接计算惯性张量 $\mathcal M_{ij}^{\mathrm{em}}(\omega)$。设静止电荷密度为 $\rho_0(\vb x)$，取电荷中心为坐标原点，
\begin{equation*}
 q=\int\rho_0(\vb x)\d^3x,
 \qquad
 \int\vb x\rho_0(\vb x)\d^3x=0.
\end{equation*}
作小振幅刚性平移并保留 $\vb X$ 的一次项，
\begin{equation*}
 \rho(\vb x,t)
 =\rho_0(\vb x)-\vb X(t)\cdot\nabla\rho_0(\vb x),
 \qquad
 \vb J(\vb x,t)=\rho_0(\vb x)\dot{\vb X}(t).
\end{equation*}
采用 Fourier 约定
\begin{equation*}
 g(\vb x,t)=\int{\d^3k\d\omega\over(2\pi)^4}
 \widetilde g(\vb k,\omega)e^{i\vb k\cdot\vb x-i\omega t},
 \qquad
 \widetilde\rho_0(\vb k)=\int\rho_0(\vb x)e^{-i\vb k\cdot\vb x}\d^3x,
\end{equation*}
在该约定下 $\nabla\mapsto i\vb k$、$\partial_t\mapsto -i\omega$；对 $\rho$ 的表达式作变换并对 $\rho_0$ 分部积分（$\widetilde{\nabla\rho_0}=i\vb k\widetilde\rho_0$），连续性方程化为
\begin{equation*}
 \widetilde{\delta\rho}
 =-i\widetilde\rho_0\,\vb k\cdot\widetilde{\vb X},
 \qquad
 \widetilde{\vb J}
 =-i\omega\widetilde\rho_0\widetilde{\vb X}.
\end{equation*}
波动方程 $\Box\Phi=-\rho/\varepsilon_0$、$\Box\vb A=-\mu_0\vb J$ 的 Fourier 形式为 $(-k^2+\omega^2/c^2)\widetilde\Phi=-\widetilde\rho/\varepsilon_0$ 等，解出推迟标势和矢势（分母中的 $-i0$ 是推迟 Green 函数 $1/(k^2-(\omega/c)^2-i0)$ 的因果处方）
\begin{equation*}
 \widetilde\Phi
 ={\widetilde{\delta\rho}\over
 \varepsilon_0[k^2-(\omega/c)^2-i0]},
 \qquad
 \widetilde{\vb A}
 ={\mu_0\widetilde{\vb J}\over
 k^2-(\omega/c)^2-i0},
\end{equation*}
由 $\vb E=-\nabla\Phi-\partial_t\vb A$，频域中 $\widetilde{\delta E}_i=-ik_i\widetilde{\delta\Phi}+i\omega\widetilde A_i$，代入上式并利用 $\mu_0=1/(\varepsilon_0c^2)$，运动引起的电场扰动为
\begin{equation*}
 \widetilde{\delta E}_i
 ={\widetilde\rho_0\over\varepsilon_0[k^2-(\omega/c)^2-i0]}
 \left({\omega^2\over c^2}\delta_{ij}-k_i k_j\right)
 \widetilde X_j.
\end{equation*}
总自力的一次项不是只积分 $\rho_0\delta\vb E$：位移密度 $\delta\rho$ 还作用于静态自场 $\vb E_0$，故
\begin{equation*}
 \widetilde F_{\mathrm{self},i}
 =\int\left(\rho_0\widetilde{\delta E}_i
 +\widetilde{\delta\rho}\,E_{0i}\right)\d^3x.
\end{equation*}
第二项与第一项的纵向部分精确相消：由 $\widetilde E_{0i}=-ik_i\widetilde\rho_0/(\varepsilon_0k^2)$，第二项按 Parseval 恒等式给出 $\int\frac{|\widetilde\rho_0|^2k_ik_j\widetilde X_j}{\varepsilon_0k^2}\frac{\d^3k}{(2\pi)^3}$；而 $\frac{\omega^2}{c^2}\delta_{ij}-k_ik_j=\frac{\omega^2}{c^2}\left(\delta_{ij}-\frac{k_ik_j}{k^2}\right)-k_ik_j\left(1-\frac{\omega^2}{c^2k^2}\right)$ 的纵向部分满足 $-\frac{k_ik_j(1-\omega^2/c^2k^2)}{k^2-\omega^2/c^2}=-\frac{k_ik_j}{k^2}$，恰好与之抵消。这正是整块电荷刚性平移不产生净自力的平移对称性。引入横向投影
\begin{equation*}
 P_{ij}^{\mathrm T}(\vb k)=\delta_{ij}-{k_i k_j\over k^2},
\end{equation*}
利用 Parseval 恒等式整理后，同一个自力依次化为
\begin{align*}
 \widetilde F_{\mathrm{self},i}
 &= {\omega^2\widetilde X_j\over\varepsilon_0c^2}
 \int{\d^3k\over(2\pi)^3}
 {\left|\widetilde\rho_0(\vb k)\right|^2P_{ij}^{\mathrm T}(\vb k)
 \over k^2-(\omega/c)^2-i0}\\
 &=-\mathcal M_{ij}^{\mathrm{em}}(\omega)\widetilde a_j,
\end{align*}
其中
\begin{equation*}
 \mathcal M_{ij}^{\mathrm{em}}(\omega)
 ={1\over\varepsilon_0c^2}
 \int{\d^3k\over(2\pi)^3}
 {\left|\widetilde\rho_0(\vb k)\right|^2P_{ij}^{\mathrm T}(\vb k)
 \over k^2-(\omega/c)^2-i0}.
\end{equation*}
这个惯性张量是线性平移问题中自力泛函的闭式化结果。它仍保留任意电荷形状、有限光行时和完整的线性辐射反作用，尚未使用长波展开。

把附加惯性分成实部和虚部，
\begin{equation*}
 \mathcal M_{ij}^{\mathrm{em}}
 =\mathcal M_{ij}'+i\mathcal M_{ij}'',
\end{equation*}
实部改变同相加速度响应，虚部产生不可逆功率。由
\begin{equation*}
 {1\over x-i0}=\mathcal P{1\over x}+i\pi\delta(x),
\end{equation*}
同一个 $\mathcal M_{ij}^{\mathrm{em}}$ 的虚部先写成三维壳层积分；$\delta(k^2-k_0^2)=\delta(k-k_0)/(2k_0)$ 使径向积分 $\int_0^\infty k^2\d k\,\delta(k^2-k_0^2)=k_0/2$（$k_0=\omega/c$），随后完成径向积分，得到更透明的角积分形式
\begin{align*}
 \mathcal M_{ij}''(\omega)
 &= {\pi\over\varepsilon_0c^2}
 \int{\d^3k\over(2\pi)^3}
 |\widetilde\rho_0(\vb k)|^2P_{ij}^{\mathrm T}(\vb k)
 \delta\!\left(k^2-{\omega^2\over c^2}\right)\\
 &= {\omega\over16\pi^2\varepsilon_0c^3}
 \int\d\Omega_{\vb n}
 \left|\widetilde\rho_0\!\left({\omega\over c}\vb n\right)\right|^2
 \left(\delta_{ij}-n_i n_j\right),
 \qquad \omega>0.
\end{align*}
对简谐加速度 $\vb a(t)=\Re(\vb a_0e^{-i\omega t})$，平均 $\langle F_iv_i\rangle=\frac12\Re(\widetilde F_i^*\widetilde v_i)$；由 $\widetilde v_i=i\widetilde a_i/\omega$、$\widetilde F_i=-\mathcal M_{ij}\widetilde a_j$ 与 $\mathcal M=\mathcal M'+i\mathcal M''$，实部 $\mathcal M'$ 的贡献为纯虚，自场所做平均功为
\begin{equation*}
 \left\langle\vb F_{\mathrm{self}}\cdot\vb v\right\rangle
 =-{1\over2\omega}a_{0i}^*\mathcal M_{ij}''(\omega)a_{0j}.
\end{equation*}
因此同一个虚部也可由辐射功率定义：
\begin{align*}
 \left\langle P_{\mathrm{rad}}\right\rangle
 &={1\over2\omega}a_{0i}^*\mathcal M_{ij}''(\omega)a_{0j}\\
 &={1\over32\pi^2\varepsilon_0c^3}
 \int\d\Omega_{\vb n}
 \left|\widetilde\rho_0\!\left({\omega\over c}\vb n\right)\right|^2
 \left|\vb n\times\vb a_0\right|^2.
\end{align*}
第一行把功率写成响应函数的耗散部分，第二行把同一功率写成前面频谱分析中的远区角积分，其中用了 $a_{0i}^*(\delta_{ij}-n_in_j)a_{0j}=|\vb n\times\vb a_0|^2$。两种计算相等，说明附加惯性和辐射谱不是两个独立问题。

因惯性张量是推迟响应，$K(\omega)$ 在上半复频率平面解析并在无穷远衰减。对 $K(\omega')/(\omega'-\omega)$ 沿实轴加无穷远半圆作 Cauchy 积分，得 $K(\omega)=\frac{1}{i\pi}\mathcal P\int_{-\infty}^{\infty}\frac{K(\omega')}{\omega'-\omega}\d\omega'$；用实值性 $K(-\omega)=K^*(\omega)$ 折到正频率，虚实部满足色散关系
\begin{equation*}
 \mathcal M_{ij}'(\omega)
 =\mathcal M_{ij}'(\infty)
 +{2\over\pi}\mathcal P\int_0^\infty
 {\omega'\mathcal M_{ij}''(\omega')\over\omega'^2-\omega^2}\d\omega'.
\end{equation*}
对有限尺寸源，$\mathcal M_{ij}'(\infty)=0$（谱表示中积分以 $k\sim|\omega|/c$ 为特征尺标，形状因子在高频衰减）；且在 $\omega=0$ 处 $\mathcal M_{ij}''(0)=0$，故 $\mathcal M_{ij}(0)=\mathcal M_{ij}'(0)$。令 $\omega=0$，同一个静态附加惯性又化为整个辐射谱密度的和规则
\begin{equation*}
 \mathcal M_{ij}^{\mathrm{em}}(0)
 ={2\over\pi}\int_0^\infty
 {\mathcal M_{ij}''(\omega')\over\omega'}\d\omega'.
\end{equation*}
因此，零频附加质量并不是只由近场“局域地”决定；它也等于全部频率辐射耗散的色散积分。

惯性张量在零频处为
\begin{equation*}
 \mathcal M_{ij}^{\mathrm{em}}(0)
 ={1\over\varepsilon_0c^2}
 \int{\d^3k\over(2\pi)^3}
 {|\widetilde\rho_0(\vb k)|^2\over k^2}
 \left(\delta_{ij}-{k_i k_j\over k^2}\right).
\end{equation*}
对指标取迹。横向投影的迹为 $\operatorname{tr}P^{\mathrm T}=3-1=2$；又由 $U_{\mathrm{es}}=\frac12\int\rho_0\Phi_0\d^3x$ 与 $\widetilde\Phi_0=\widetilde\rho_0/(\varepsilon_0k^2)$ 经 Parseval 恒等式，同一个积分可识别成静电能
\begin{equation*}
 U_{\mathrm{es}}
 ={1\over2\varepsilon_0}
 \int{\d^3k\over(2\pi)^3}{|\widetilde\rho_0(\vb k)|^2\over k^2},
\end{equation*}
便有连续的等式链
\begin{equation*}
 \operatorname{tr}\boldsymbol{\mathcal M}^{\mathrm{em}}(0)
 ={2\over\varepsilon_0c^2}
 \int{\d^3k\over(2\pi)^3}{|\widetilde\rho_0|^2\over k^2}
 ={4U_{\mathrm{es}}\over c^2}.
\end{equation*}
这是一条适用于任意形状电荷分布的迹定理。若再取球对称，旋转不变性使
\begin{equation*}
 \mathcal M_{ij}^{\mathrm{em}}(0)=\delta m_{\mathrm{dyn}}\delta_{ij},
\end{equation*}
于是迹定理的左端成为 $3\delta m_{\mathrm{dyn}}$，同一个结论才进一步简化为
\begin{equation*}
 \delta m_{\mathrm{dyn}}
 ={1\over3}\operatorname{tr}\boldsymbol{\mathcal M}^{\mathrm{em}}(0)
 ={4U_{\mathrm{es}}\over3c^2}.
\end{equation*}
这说明 $4/3$ 系数首先是一般附加惯性张量在球对称下的标量化结果，而不是可以在任意源模型中预先假定的常数。这是 $4/3$ 的第一次出现。

接下来在球对称条件下把同一个频率响应进一步化简，并展示 Jackson--Abraham--Lorentz 推导中实空间展开、闭式求和和形状因子表示之间的关系。设电荷分布作非相对论刚性平移，瞬时静止系中 $\vb J(\vb x,t)=\rho(\vb x)\vb v(t)$。定义电磁动量对机械方程的贡献为
\begin{equation*}
 {\d\vb p_{\mathrm{em}}\over\d t}
 \equiv-\vb F_{\mathrm{self}}.
\end{equation*}
将推迟势中的源量在 $t-R/c$ 处展开，
\begin{equation*}
 [g]_{\mathrm{ret}}
 =\sum_{n=0}^\infty{(-1)^n\over n!}
 \left({R\over c}\right)^n g^{(n)}(t),
 \qquad R=|\vb x-\vb x'|,
\end{equation*}
代入自场 Lorentz 力。推迟展开的 $n$ 阶项中，标势的 $\vb R/R^2$ 部分经连续性方程 $\dot\rho=-\nabla'\cdot\vb J$ 并入矢势项，$n$ 阶积分中出现
\begin{equation*}
 \vb J(\vb x')-{\vb R\over n+2}\nabla'\cdot\vb J(\vb x').
\end{equation*}
对第二项分部积分（$\partial'_jR_i=-\delta_{ij}$、$\partial'_jR=-R_j/R$），在双重积分内它等价于
\begin{equation*}
 {n+1\over n+2}\vb J
 -{n-1\over n+2}{(\vb J\cdot\vb R)\vb R\over R^2}.
\end{equation*}
再代入 $\vb J=\rho\vb v$。球对称使所有 $\vb R$ 方向等权：$\langle R_iR_j\rangle_\Omega=\alpha\delta_{ij}$，两侧缩并 $\langle R^2\rangle_\Omega=3\alpha$ 定出 $\alpha=R^2/3$，故
\begin{equation*}
 \left\langle{(\vb v\cdot\vb R)\vb R\over R^2}\right\rangle_{\Omega_R}
 ={1\over3}\vb v.
\end{equation*}
于是上面的冗长向量组合在同一个积分中化成
\begin{equation*}
 {n+1\over n+2}\vb J
 -{n-1\over n+2}{(\vb J\cdot\vb R)\vb R\over R^2}
 \longrightarrow {2\over3}\rho(\vb x')\vb v.
\end{equation*}
自场动量率因此得到只含标量距离矩的展开
\begin{equation*}
 {\d\vb p_{\mathrm{em}}\over\d t}
 ={1\over6\pi\varepsilon_0c^2}
 \sum_{n=0}^\infty{(-1)^n\over n!c^n}
 \vb v^{(n+1)}(t)
 \iint\rho(\vb x)\rho(\vb x')R^{n-1}\d^3x\d^3x'.
\end{equation*}
级数是同一个球对称自作用的时间域导数展开。其 $n=0$ 项是附加惯性，$n=1$ 项是形状无关的辐射反作用：
\begin{align*}
 {\d\vb p_{\mathrm{em}}\over\d t}
 &=\left[{1\over6\pi\varepsilon_0c^2}
 \iint{\rho(\vb x)\rho(\vb x')\over R}\d^3x\d^3x'\right]\vb a
 -{q^2\over6\pi\varepsilon_0c^3}\dot{\vb a}+\cdots\\
 &={4U_{\mathrm{es}}\over3c^2}\vb a
 -{q^2\over6\pi\varepsilon_0c^3}\dot{\vb a}+\cdots.
\end{align*}
其中 $n=0$ 项系数用 $U_{\mathrm{es}}=\frac12\int\rho\Phi\d^3x=\frac{1}{8\pi\varepsilon_0}\iint\frac{\rho(\vb x)\rho(\vb x')}{|\vb x-\vb x'|}\d^3x\d^3x'$ 化为 $4U_{\mathrm{es}}/(3c^2)$，$n=1$ 项系数用 $\iint\rho(\vb x)\rho(\vb x')\d^3x\d^3x'=q^2$ 给出 $q^2/(6\pi\varepsilon_0c^3)$。由于 $\vb F_{\mathrm{self}}=-\d\vb p_{\mathrm{em}}/\d t$，上式立即给出熟悉的惯性反力和 Abraham--Lorentz 项，但它们此时只是完整级数的前两项。

对级数作 Fourier 变换，$\vb v^{(n+1)}\mapsto(-i\omega)^{n+1}\widetilde{\vb v}$。系数部分求和（$(-1)^n(-i)^n=i^n$）：
\begin{equation*}
 \sum_{n=0}^\infty{(-1)^n\over n!c^n}(-i\omega)^{n+1}R^{n-1}
 ={-i\omega\over R}\sum_{n=0}^\infty{(i\omega R/c)^n\over n!}
 ={-i\omega\over R}e^{i\omega R/c},
\end{equation*}
幂级数中的距离矩重新求和为
\begin{equation*}
 \sum_{n=0}^\infty{1\over n!}
 \left({i\omega R\over c}\right)^n
 ={e^{i\omega R/c}},
\end{equation*}
因此同一个无穷导数级数精确重写成
\begin{align*}
 \widetilde{\dot{\vb p}}_{\mathrm{em}}
 &=-i\omega M_{\mathrm{em}}(\omega)\widetilde{\vb v},\\
 M_{\mathrm{em}}(\omega)
 &={1\over6\pi\varepsilon_0c^2}
 \iint\rho(\vb x)\rho(\vb x')
 {e^{i\omega R/c}\over R}\d^3x\d^3x'.
\end{align*}
级数形式和闭式形式是同一个物理量的两种写法；前者适合长波展开，后者保留全部光行时。

再定义球对称形状因子
\begin{equation*}
 \widetilde\rho(\vb k)=qf(k),
 \qquad f(0)=1.
\end{equation*}
用外向波 Green 函数的标准展开（$(\nabla^2+k_0^2)G=\delta^{(3)}$ 的 Fourier 解）
\begin{equation*}
 {e^{ik_0R}\over R}
 =4\pi\int{\d^3k\over(2\pi)^3}{e^{i\vb k\cdot\vb R}\over k^2-k_0^2-i0},
 \qquad k_0={\omega\over c},
\end{equation*}
代入闭式并对 $d^3x\,d^3x'$ 作卷积，$|\widetilde\rho(\vb k)|^2=q^2|f(k)|^2$，角向积分 $\int\d\Omega=4\pi$ 后，同一个 $M_{\mathrm{em}}(\omega)$ 化成谱表示
\begin{equation*}
 M_{\mathrm{em}}(\omega)
 ={q^2\over3\pi^2\varepsilon_0c^2}
 \int_0^\infty
 {k^2|f(k)|^2\over k^2-(\omega/c)^2-i0}\d k.
\end{equation*}
实空间表示和谱表示分别把同一附加惯性写成电荷对之间的推迟相互作用和空间频率的叠加：前者显示有限传播时延，后者显示源形状如何截断高频响应。

若裸机械质量为 $m_0$，机械动量和电磁动量相加后，外力方程为
\begin{equation*}
 -i\omega\mathsf M(\omega)\widetilde{\vb v}
 =\widetilde{\vb F}_{\mathrm{ext}},
 \qquad
 \mathsf M(\omega)=m_0+M_{\mathrm{em}}(\omega).
\end{equation*}
物理低频质量是同一个有效质量在零频处的值，
\begin{equation*}
 m\equiv\mathsf M(0)
 =m_0+{4U_{\mathrm{es}}\over3c^2}.
\end{equation*}
用低频质量消去裸质量后，外力方程可改写成更清楚、紫外收敛性更好的形式
\begin{align*}
 \mathsf M(\omega)
 &=m+{1\over6\pi\varepsilon_0c^2}
 \iint\rho\rho'
 {e^{i\omega R/c}-1\over R}\d^3x\d^3x'\\
 &=m+{q^2\omega^2\over3\pi^2\varepsilon_0c^4}
 \int_0^\infty
 {|f(k)|^2\over k^2-(\omega/c)^2-i0}\d k.
\end{align*}
第二行由第一行代入谱表示并相减：$\frac{e^{i\omega R/c}-1}{R}$ 对应频域差 $\frac{1}{k^2-k_0^2-i0}-\frac{1}{k^2}=\frac{k_0^2}{k^2(k^2-k_0^2-i0)}$。两行仍是同一个 $\mathsf M(\omega)$：减去零频值以后，原来线性发散的静态质量被隔离，剩余的频率依赖部分多了两个高频幂次的收敛性。

作为与 Dirac 点粒子结果的交叉检验，现在在已经减去零频质量的响应中取点形状 $f(k)=1$。利用
\begin{equation*}
 \int_0^\infty{\d k\over k^2-(\omega/c)^2-i0}
 ={i\pi c\over2\omega},
 \qquad \omega>0,
\end{equation*}
该值是主值 $\mathcal P\int_0^\infty\frac{\d k}{k^2-k_0^2}=0$（被积函数关于 $k_0$ 奇对称）与 $k=k_0$ 处留数 $\pi i/(2k_0)$ 之和。同一个有效质量化为
\begin{equation*}
 \mathsf M_{\mathrm{point}}(\omega)
 =m+i{q^2\omega\over6\pi\varepsilon_0c^3}
 =m(1+i\omega\tau_{\mathrm{rr}}),
 \qquad
 \tau_{\mathrm{rr}}={q^2\over6\pi\varepsilon_0mc^3}.
\end{equation*}
代回外力方程并逆变换，得到
\begin{equation*}
 m\vb a
 =\vb F_{\mathrm{ext}}+m\tau_{\mathrm{rr}}\dot{\vb a}.
\end{equation*}
因此这里重新得到的 Abraham--Lorentz 方程与 Dirac 正则场的低速极限完全相同；它也是同一个有限尺寸有效质量 $\mathsf M(\omega)$ 在先减去静态质量、再取点形状以后得到的特殊形式。

点形状抹掉了高频结构。为看清有限尺寸如何恢复时延，取半径为 $a$ 的均匀薄球壳，
\begin{equation*}
 f(k)={\sin ka\over ka},
 \qquad
 U_{\mathrm{es}}={q^2\over8\pi\varepsilon_0a},
 \qquad
 \delta m={4U_{\mathrm{es}}\over3c^2}
 ={q^2\over6\pi\varepsilon_0ac^2}.
\end{equation*}
将这个 $f(k)$ 代入形状因子公式。利用下列标准积分（第一式由 $\sin^2ka=\frac12(1-\cos2ka)$ 与 $\mathcal P\int_0^\infty\frac{\cos(bk)}{k^2-k_0^2}\d k=-\frac{\pi}{2k_0}\sin(bk_0)$ 得到；第二式取 $\Im\frac{1}{k^2-k_0^2-i0}=\pi\delta(k^2-k_0^2)$ 的留数）
\begin{equation*}
 \mathcal P\int_0^\infty{\sin^2(ka)\over k^2-k_0^2}\d k
 ={\pi\over4k_0}\sin(2k_0a),
 \qquad
 \Im\int_0^\infty{\sin^2(ka)\over k^2-k_0^2-i0}\d k
 ={\pi\over2k_0}\sin^2(k_0a),
\end{equation*}
得到
\begin{equation*}
 M_{\mathrm{em}}^{\mathrm{shell}}(\omega)
 =\delta m\,{e^{i\zeta_a}-1\over i\zeta_a},
 \qquad
 \zeta_a={2\omega a\over c}.
\end{equation*}
核对：$\frac{e^{i\zeta}-1}{i\zeta}=\frac{\sin\zeta}{\zeta}+i\frac{1-\cos\zeta}{\zeta}$，代入 $\delta m=q^2/(6\pi\varepsilon_0ac^2)$ 后实部、虚部与上面两个积分一致；零频极限 $M_{\mathrm{em}}^{\mathrm{shell}}(0)=\delta m$（$\int_0^\infty\frac{\sin^2x}{x^2}\d x=\pi/2$）。同一个球壳自力先在频域写成
\begin{align*}
 \widetilde{\vb F}_{\mathrm{self}}
 &=-M_{\mathrm{em}}^{\mathrm{shell}}(\omega)\widetilde{\vb a}\\
 &={q^2\over12\pi\varepsilon_0ca^2}
 \left(e^{2i\omega a/c}-1\right)\widetilde{\vb v},
\end{align*}
逆变换后仍是同一个自力，但其因果结构变得直接可见：
\begin{equation*}
 \vb F_{\mathrm{self}}(t)
 ={q^2\over12\pi\varepsilon_0ca^2}
 \left[\vb v\!\left(t-{2a\over c}\right)-\vb v(t)\right].
\end{equation*}
最后在 $|\omega|a/c\ll1$ 下展开这个精确延迟力。令 $\tau=2a/c$，Taylor 展开 $\vb v(t-\tau)-\vb v(t)=-\tau\vb a+\frac{\tau^2}{2}\dot{\vb a}-\frac{\tau^3}{6}\ddot{\vb a}+\cdots$，得
\begin{align*}
 \vb F_{\mathrm{self}}
 &={q^2\over12\pi\varepsilon_0ca^2}
 \left[-{2a\over c}\vb a
 +{1\over2}\left({2a\over c}\right)^2\dot{\vb a}
 -{1\over6}\left({2a\over c}\right)^3\ddot{\vb a}+\cdots\right]\\
 &=-\delta m\,\vb a
 +{q^2\over6\pi\varepsilon_0c^3}\dot{\vb a}
 -{q^2a\over9\pi\varepsilon_0c^4}\ddot{\vb a}+\cdots.
\end{align*}
频率响应、频域力、时域延迟和局域渐近是同一个球壳自作用的四种写法。只有最后一行出现三阶导数，因此 runaway 问题属于局域截断，而不是精确延迟核的必然性质。

外力为零时频域方程 $-i\omega[m_0+M_{\mathrm{em}}^{\mathrm{shell}}(\omega)]\widetilde{\vb v}=0$ 给出无外力模态：
\begin{equation*}
 \omega\left[m_0+M_{\mathrm{em}}^{\mathrm{shell}}(\omega)\right]=0,
\end{equation*}
除平移零模（$\omega=0$）外，稳定性要求响应函数在上半频率平面无零点。在上半平面取 $\omega=i\xi$（$\xi>0$），此时 $e^{i\zeta_a}=e^{-2\xi a/c}$，$m_0+M_{\mathrm{em}}^{\mathrm{shell}}(i\xi)=m_0+\delta m\frac{1-e^{-2\xi a/c}}{2\xi a/c}$ 是 $\xi$ 的单调减函数（$\frac{1-e^{-x}}{x}$ 在 $x>0$ 递减），从 $m$（$\xi\to0$）降到 $m_0$（$\xi\to\infty$）；除非 $m_0\geq0$，否则它必然穿越零点，给出 $e^{\xi t}$ 增长的不稳定模态。因此无零点的条件可写成
\begin{equation*}
 m_0\geq0
 \qquad\Longleftrightarrow\qquad
 a\geq c\tau_{\mathrm{rr}},
\end{equation*}
其中 $m=m_0+\delta m$。若在保持物理质量 $m$ 固定时把 $a$ 压到 $c\tau_{\mathrm{rr}}$ 以下，就必须让裸质量变成负值，精确响应的稳定性随之丧失。点粒子三阶方程中的异常解由这一极限结构产生，而不是由“辐射必然导致反因果”产生。

光滑电荷给出同一母式的另一种形象特例。取
\begin{equation*}
 \rho_0(r)={q\over\pi^{3/2}a^3}e^{-r^2/a^2},
 \qquad
 f(k)=e^{-a^2k^2/4},
 \qquad
 x={\omega a\over\sqrt2c}.
\end{equation*}
将它代入形状因子公式：常数部分 $\int_0^\infty e^{-a^2k^2/2}\d k=\frac1a\sqrt{\frac\pi2}$，频率依赖部分用
\begin{equation*}
 \mathcal P\int_0^\infty{e^{-\alpha k^2}\over k^2-b^2}\d k
 =-{\pi\over2b}e^{-\alpha b^2}\operatorname{erfi}(b\sqrt\alpha),
\end{equation*}
同一个 $M_{\mathrm{em}}(\omega)$ 化为闭式
\begin{equation*}
 M_{\mathrm{em}}^{\mathrm G}(\omega)
 ={q^2\over3\pi^2\varepsilon_0c^2}
 \left[
 {\sqrt\pi\over\sqrt2a}
 -{\pi\omega\over2c}e^{-x^2}\operatorname{erfi}(x)
 +i{\pi\omega\over2c}e^{-x^2}
 \right].
\end{equation*}
其零频值仍是 $4U_{\mathrm{es}}/(3c^2)$，而虚部被 $e^{-\omega^2a^2/(2c^2)}$ 平滑截断。球壳和 Gaussian 分布具有不同高频核，却共享同一个低频辐射系数，这正是一般角积分公式的内容。

到这里，$4U_{\mathrm{es}}/(3c^2)$ 已经由自力响应得到。现在从场动量重新计算同一个零频惯性，便可看清传统 $4/3$ 问题的来源。将静电构型以小速度 $\vb v$ 平移，最低阶磁场为
\begin{equation*}
 \vb B={1\over c^2}\vb v\times\vb E.
\end{equation*}
其中 $(\vb E\times\vb B)_i=\frac{1}{c^2}\epsilon_{ijk}\epsilon_{klm}E_jv_lE_m=\frac{1}{c^2}(E^2\delta_{ij}-E_iE_j)v_j$（用到 $\epsilon_{ijk}\epsilon_{klm}=\delta_{il}\delta_{jm}-\delta_{im}\delta_{jl}$），电磁动量为
\begin{align*}
 P_{\mathrm{em},i}
 &=\varepsilon_0\int(\vb E\times\vb B)_i\d^3x\\
 &={\varepsilon_0\over c^2}
 \int\left(E^2\delta_{ij}-E_iE_j\right)\d^3x\,v_j.
\end{align*}
静电场的 Fourier 分量为（$\widetilde\Phi_0=\widetilde\rho_0/(\varepsilon_0k^2)$，$\vb E_0=-\nabla\Phi_0$）
\begin{equation*}
 \widetilde E_i(\vb k)
 =-{ik_i\over\varepsilon_0k^2}\widetilde\rho_0(\vb k).
\end{equation*}
将它代入电磁动量并用 Parseval 恒等式（$\int E_iE_j\d^3x=\int\frac{k_ik_j|\widetilde\rho_0|^2}{\varepsilon_0^2k^4}\frac{\d^3k}{(2\pi)^3}$，$\int E^2\d^3x=\int\frac{|\widetilde\rho_0|^2}{\varepsilon_0^2k^2}\frac{\d^3k}{(2\pi)^3}$），方括号中的惯性张量化成
\begin{align*}
 {\varepsilon_0\over c^2}
 \int\left(E^2\delta_{ij}-E_iE_j\right)\d^3x
 &={1\over\varepsilon_0c^2}
 \int{\d^3k\over(2\pi)^3}
 {|\widetilde\rho_0|^2\over k^2}
 \left(\delta_{ij}-{k_i k_j\over k^2}\right)\\
 &=\mathcal M_{ij}^{\mathrm{em}}(0).
\end{align*}
因此场动量和前面由自力得到的零频附加惯性是同一个量：
\begin{equation*}
 P_{\mathrm{em},i}
 =\mathcal M_{ij}^{\mathrm{em}}(0)v_j.
\end{equation*}
球对称时，上式立即给出
\begin{equation*}
 \vb P_{\mathrm{em}}
 ={4U_{\mathrm{es}}\over3c^2}\vb v.
\end{equation*}
这不是第二个无关的 $4/3$ 结论，而是同一 $\mathcal M^{\mathrm{em}}(0)$ 分别在自力和场动量中的出现，这是 $4/3$ 的第二次出现。

然而纯电磁场不是孤立守恒系统。对任意静止构型，在小速度 boost 后，同一实验室同时面上的总动量为
\begin{equation*}
 P^i={v_j\over c^2}
 \left[
 \delta_{ij}\int T^{00}\d^3x
 +\int T^{ij}\d^3x
 \right]+O(v^3/c^3).
\end{equation*}
其来源：小速度 boost 下 $T^{i0}=T'^{i0}+\frac{v_j}{c}T'^{ij}+\frac{v^i}{c}T'^{00}+O(v^2/c^2)$，静止构型 $T'^{i0}=0$，由 $P^i=\frac1c\int T^{i0}\d^3x$ 即得该式。它表明，能量密度给出 $E/c^2$，积分应力还会给出额外惯性。纯电磁球对称场满足
\begin{equation*}
 \int T_{\mathrm{em}}^{ij}\d^3x
 ={U_{\mathrm{es}}\over3}\delta_{ij},
\end{equation*}
所以总系数成为 $4U_{\mathrm{es}}/(3c^2)$，这是 $4/3$ 的第三次出现。稳定带电体还必须含有抵抗 Coulomb 排斥的非电磁应力 $T_{\mathrm P}^{\mu\nu}$。对局域、静态且总能动量守恒的系统，
\begin{equation*}
 \partial_\mu T_{\mathrm{tot}}^{\mu\nu}=0,
 \qquad
 T_{\mathrm{tot}}^{\mu\nu}=T_{\mathrm{em}}^{\mu\nu}+T_{\mathrm P}^{\mu\nu},
\end{equation*}
由
\begin{equation*}
 \partial_j(x^iT_{\mathrm{tot}}^{jk})
 =T_{\mathrm{tot}}^{ik}+x^i\partial_jT_{\mathrm{tot}}^{jk}
\end{equation*}
在全空间积分。静态平衡 $\partial_jT^{jk}=0$ 使右端第二项为零；左端是散度的全空间积分，化为无穷远面积分，对有限源消失，得到 Laue 条件
\begin{equation*}
 \int T_{\mathrm{tot}}^{ij}\d^3x=0.
\end{equation*}
把 Laue 条件代回 boost 公式，总动量才化为
\begin{equation*}
 \vb P_{\mathrm{tot}}
 ={E_{\mathrm{tot}}\over c^2}\vb v.
\end{equation*}
因此 $U_{\mathrm{es}}/c^2$ 与 $4U_{\mathrm{es}}/(3c^2)$ 并非互相矛盾的两个“质量”。后者是非守恒电磁子系统在实验室同时面上的动力学惯性，前者是能量项；只有把维持稳定的物质应力一并加入，上式才给出协变的总质量。

薄球壳提供这一应力闭合的熟悉图像。由 Gauss 定律 $E(a^+)=q/(4\pi\varepsilon_0a^2)$，壳外电场对球面产生向外 Maxwell 压力
\begin{equation*}
 p_{\mathrm{em}}
 ={\varepsilon_0E^2(a^+)\over2}
 ={q^2\over32\pi^2\varepsilon_0a^4}.
\end{equation*}
若以表面张力 $\mathcal T$ 维持平衡，Laplace 关系 $2\mathcal T/a=p_{\mathrm{em}}$ 给出
\begin{equation*}
 \mathcal T={q^2\over64\pi^2\varepsilon_0a^3}.
\end{equation*}
该张力的积分应力抵消电磁场的积分应力，使 Laue 条件成立；至于稳定物质本身含有多少静能，则仍由具体微观模型决定。至此有限尺寸正则化的图像闭合：自力的 $4/3$、场动量的 $4/3$ 与 boost 公式中的 $4/3$ 是同一个 $\mathcal M^{\mathrm{em}}(0)$ 的不同出场方式，而有效质量 $\mathsf M(\omega)$ 把附加惯性、频率依赖惯性与延迟自力统一在同一响应函数中。下一节把这个 $\mathsf M(\omega)$ 用于辐射振子，看它如何同时决定附加质量、频移、线宽与散射截面。

\subsection{辐射反作用的应用}

前面建立了辐射反作用的两条主线：在无反作用轨道上进行频谱分析，以及由 Dirac 方程降阶得到的缓慢轨道反作用。本节把两条主线合并起来，应用到两个具体体系：辐射振子，以及有限质量二体库仑运动。后者先沿无反作用轨道求完整频谱，再把频谱积分给出的能量、角动量通量绝热地反馈到轨道常数。

\subsubsection{辐射振子、自然线宽与散射截面}

用辐射振子说明同一个 $\mathsf M(\omega)$ 如何同时决定附加质量、频移、线宽和散射。设恢复力为 $-K\vb x$，在尚未取点极限时，由前面自作用总有效质量的频域运动方程直接得到
\begin{equation*}
	\left[K-\omega^2\mathsf M(\omega)\right]\widetilde{\vb x}
	=q\widetilde{\vb E},
\end{equation*}
故偶极矩 $\widetilde{\vb p}=q\widetilde{\vb x}=\alpha(\omega)\widetilde{\vb E}$ 的精确线性极化率为
\begin{equation*}
	\alpha(\omega)
	={q^2\over K-\omega^2\mathsf M(\omega)}.
\end{equation*}
定义物理质量 $m=\mathsf M(0)$ 和 $\omega_0^2=K/m$，并写
\begin{equation*}
	\Delta\mathsf M(\omega)=\mathsf M(\omega)-m.
\end{equation*}
则极化率中同一个分母改写为
\begin{equation*}
	K-\omega^2\mathsf M(\omega)
	=m(\omega_0^2-\omega^2)-\omega^2\Delta\mathsf M(\omega).
\end{equation*}
弱自作用时，极点 $\omega_{\mathrm p}=\omega_0+\delta\omega-i\Gamma/2$ 由分母改写式一阶给出
\begin{equation*}
	\delta\omega
	=-{\omega_0\over2m}\Re\Delta\mathsf M(\omega_0),
	\qquad
	\Gamma
	={\omega_0\over m}\Im\mathsf M(\omega_0).
\end{equation*}
所以频移来自同一响应的色散实部，线宽来自同一响应的耗散虚部。点粒子长波极限下
\begin{equation*}
	\Delta\mathsf M(\omega)=im\tau_{\mathrm{rr}}\omega,
\end{equation*}
故
\begin{equation*}
	\delta\omega=0,
	\qquad
	\Gamma=\tau_{\mathrm{rr}}\omega_0^2,
\end{equation*}
而极化率化为
\begin{equation*}
	\alpha(\omega)
	={q^2/m\over\omega_0^2-\omega^2-i\tau_{\mathrm{rr}}\omega^3}.
\end{equation*}
若还有非辐射宽度 $\Gamma'$，只需把分母再减去 $i\Gamma'\omega$。在共振附近可用 $\tau_{\mathrm{rr}}\omega^3\simeq\Gamma\omega$，总宽度为 $\Gamma_{\mathrm t}=\Gamma+\Gamma'$。

同一个极化率决定散射、消光和吸收：
\begin{equation*}
	\sigma_{\mathrm{sca}}(\omega)
	={\omega^4\over6\pi\varepsilon_0^2c^4}|\alpha(\omega)|^2,
	\qquad
	\sigma_{\mathrm{ext}}(\omega)
	={\omega\over\varepsilon_0c}\Im\alpha(\omega),
	\qquad
	\sigma_{\mathrm{abs}}=\sigma_{\mathrm{ext}}-\sigma_{\mathrm{sca}}.
\end{equation*}
当阻尼完全来自辐射时，将点极限极化率代入上式，可直接验证
\begin{equation*}
	\sigma_{\mathrm{ext}}=\sigma_{\mathrm{sca}},
	\qquad
	\sigma_{\mathrm{abs}}=0.
\end{equation*}
自由粒子极限 $\omega_0=0$ 且 $\omega\tau_{\mathrm{rr}}\ll1$ 时，同一个散射截面进一步化为 Thomson 值
\begin{equation*}
	\sigma_{\mathrm T}
	={q^4\over6\pi\varepsilon_0^2m^2c^4}
	={8\pi\over3}
	\left({q^2\over4\pi\varepsilon_0mc^2}\right)^2.
\end{equation*}
至此，附加惯性、辐射功率、线宽、频移和散射截面都由同一个因果响应 $\mathsf M(\omega)$ 贯通起来。

后面的二体库仑问题同时应用两条已经建立的主线：无反作用轨道上的频谱分析，以及由 Dirac 方程降阶得到的缓慢轨道反作用。质心系中的电偶极矩为
\begin{equation*}
	\vb p=q_D\vb r,
	\qquad
	q_D=\mu\left({q_1\over M_1}-{q_2\over M_2}\right).
\end{equation*}
在长波、弱反作用和降阶近似下，降阶后的自作用运动方程对相对坐标给出
\begin{equation*}
	\mu\ddot{\vb r}
	=\vb F(\vb r)
	+{q_D^2\over6\pi\varepsilon_0c^3}\dddot{\vb r}.
\end{equation*}
若 $q_D=0$，偶极频谱和偶极反作用同时消失，领先项转为电四极阶。因而库仑例题中先在无反作用轨道上求完整频谱，再把频谱积分得到的能量、角动量通量绝热地反馈到轨道常数，正是前面世界管守恒关系在弱辐射条件下的具体实现。

本节中三个常被写成"电磁质量"的系数现在可以按其推导位置区分。$U_{\mathrm{es}}/c^2$ 是静止系电磁能量除以 $c^2$；$4U_{\mathrm{es}}/(3c^2)$ 是球对称电磁子系统在实验室同时面上的零频平移惯性；$q^2/(8\pi\varepsilon_0c^2\epsilon)$ 是 Dirac 正交世界管中的奇异束缚动量系数。它们分别来自不同的、已经在上文明确写出的物理量，不能在省略源模型、积分超曲面和极限次序以后互相替换。

\subsubsection{例：有限质量二体库仑运动的辐射}

本例不再重新建立频谱理论，而是逐项实现前面的一般程序。首先在质心系中把两个可反冲点电荷的多极矩化为相对坐标的函数；随后求解统一包含吸引与排斥的库仑圆锥曲线。正相对能量的双曲轨道按照非周期过程的 Fourier 变换约定产生连续谱，负相对能量的椭圆轨道按照 Fourier 级数约定产生离散线谱；固定轨道的结果分别代入一般电偶极、电四极频域频谱公式，而散射系综则使用能量加权的一般谱辐射截面公式。这样，轨道动力学、特殊函数计算和辐射学归一化在逻辑上彼此分离，又能在最终结果中相互校验。

\textbf{质心系约化与统一轨道。}

设两粒子的质量和有符号电荷分别为 $(M_1,q_1)$、$(M_2,q_2)$。取
\begin{equation*}
	M=M_1+M_2,
	\qquad
	\mu=\frac{M_1M_2}{M},
	\qquad
	\vb R=\frac{M_1\vb r_1+M_2\vb r_2}{M},
	\qquad
	\vb r=\vb r_1-\vb r_2.
\end{equation*}
在质心系 $\vb R=0$ 中
\begin{equation*}
	\vb r_1=\frac{M_2}{M}\vb r,
	\qquad
	\vb r_2=-\frac{M_1}{M}\vb r.
\end{equation*}
因此电偶极矩约化为
\begin{equation*}
	\vb p=q_D\vb r,
	\qquad
	q_D=\frac{q_1M_2-q_2M_1}{M}
	=\mu\left(\frac{q_1}{M_1}-\frac{q_2}{M_2}\right),
\end{equation*}
无迹电四极矩约化为
\begin{equation*}
	D_{ij}=q_Q\left(3r_ir_j-r^2\delta_{ij}\right),
	\qquad
	q_Q=\frac{q_1M_2^2+q_2M_1^2}{M^2}
	=\mu^2\left(\frac{q_1}{M_1^2}+\frac{q_2}{M_2^2}\right).
\end{equation*}
轨道磁偶极矩则为
\begin{equation*}
	\vb m_{\mathrm{orb}}
	=\frac{1}{2}\sum_{a=1}^2q_a\vb r_a\times\dot{\vb r}_a
	=\frac{q_Q}{2\mu}\vb L,
	\qquad
	\vb L=\mu\vb r\times\dot{\vb r}.
\end{equation*}
对任意中心力零阶轨道，$\vb L$ 守恒，故 $\ddot{\vb m}_{\mathrm{orb}}=0$。于是本问题的正频率轨道辐射没有 $M1$ 项；保留至相对阶 $v^2/c^2$ 时，只需计算 $E1$ 与 $E2$。

库仑相互作用写成
\begin{equation*}
	V(r)=\sigma\frac{\mathcal C}{r},
	\qquad
	\mathcal C=\frac{|q_1q_2|}{4\pi\varepsilon_0},
	\qquad
	\sigma=\operatorname{sgn}(q_1q_2).
\end{equation*}
其中 $\sigma=+1$ 表示排斥，$\sigma=-1$ 表示吸引。相对运动拉格朗日量为
\begin{equation*}
	L_{\mathrm{rel}}=\frac{1}{2}\mu\dot{\vb r}^{\,2}-\sigma\frac{\mathcal C}{r},
\end{equation*}
从而
\begin{equation*}
	\mu\ddot{\vb r}=\sigma\mathcal C\frac{\vb r}{r^3}.
\end{equation*}
有限靶质量只通过 $\mu,q_D,q_Q$ 进入结果，并不破坏相对运动的中心力可积性。

在完成源的约化以后，考虑正相对能量的非束缚运动。设无穷远相对速度为 $v$，碰撞参数为 $b$，则
\begin{equation*}
	\mathcal E=\frac{1}{2}\mu v^2,
	\qquad
	L=\mu vb.
\end{equation*}
令 $w(\theta)=1/r$。由 Binet 方程
\begin{equation*}
	\frac{\d^2w}{\d\theta^2}+w
	=-\frac{\mu}{L^2w^2}\left(\sigma\mathcal Cw^2\right)
	=-\sigma\frac{\mu\mathcal C}{L^2},
\end{equation*}
并取近心点为 $\theta=0$，可得
\begin{equation*}
	r(\theta)=\frac{p}{\varepsilon\cos\theta-\sigma},
	\qquad
	p=\frac{L^2}{\mu\mathcal C},
	\qquad
	\varepsilon^2=1+\frac{2\mathcal EL^2}{\mu\mathcal C^2}.
\end{equation*}
定义无量纲参数
\begin{equation*}
	b_0=\frac{\mathcal C}{\mu v^2},
	\qquad
	\omega_0=\frac{v}{b_0},
	\qquad
	s=\frac{b}{b_0},
	\qquad
	\varepsilon=\sqrt{1+s^2},
	\qquad
	\nu=\frac{\omega}{\omega_0}.
\end{equation*}
此时 $p=b_0s^2$。为进行傅里叶积分，引入双曲参数 $u$：
\begin{equation*}
	\begin{aligned}
	x(u)&=b_0\left(\varepsilon+\sigma\cosh u\right),\\
	y(u)&=b_0s\sinh u,\\
	r(u)&=b_0\left(\varepsilon\cosh u+\sigma\right).
	\end{aligned}
\end{equation*}
确有 $x^2+y^2=r^2$，且由 $p=\varepsilon x-\sigma r$ 可恢复上面的圆锥曲线轨道。面积速度守恒给出
\begin{equation*}
	x\dot y-y\dot x=\frac{L}{\mu}=vb.
\end{equation*}
另一方面
\begin{equation*}
	x\frac{\d y}{\d u}-y\frac{\d x}{\d u}
	=b_0^2s\left(\varepsilon\cosh u+\sigma\right),
\end{equation*}
故
\begin{equation*}
	\frac{\d t}{\d u}
	=\frac{1}{\omega_0}\left(\varepsilon\cosh u+\sigma\right).
\end{equation*}
取近心点 $u=0$ 时 $t=0$，得到统一的时间参数式
\begin{equation*}
	t(u)=\frac{1}{\omega_0}\left(\varepsilon\sinh u+\sigma u\right).
\end{equation*}
空间参数式与时间参数式合并即为完整的双曲轨道参数化。最近距离为
\begin{equation*}
	r_{\min}=b_0(\varepsilon+\sigma),
\end{equation*}
散射角的绝对值满足
\begin{equation*}
	\tan\frac{\Theta}{2}=\frac{1}{s},
	\qquad
	\sin\frac{\Theta}{2}=\frac{1}{\varepsilon}.
\end{equation*}

\textbf{母积分与虚阶 Bessel 函数。}

为计算双曲轨道上的连续谱，需要先建立一个能够生成全部轨道积分的母公式。

所有连续谱积分均包含相位
\begin{equation*}
	\omega t(u)=\nu\left(\varepsilon\sinh u+\sigma u\right).
\end{equation*}
为统一处理 $\sinh u$、$\cosh u$ 的有限次幂，定义整数 $n$ 的母积分
\begin{equation*}
	\mathcal J_n^{(\sigma)}(\nu,\varepsilon)
	=\int_{-\infty}^{\infty}
	e^{nu}e^{i\nu(\varepsilon\sinh u+\sigma u)}\d u.
\end{equation*}
严格地说，可先在被积函数中加入收敛因子，再在积分完成后令收敛参数趋于零。将积分路径平移为 $u=w+i\pi/2$，利用
\begin{equation*}
	\sinh\left(w+\frac{i\pi}{2}\right)=i\cosh w
\end{equation*}
以及第二类修正 Bessel 函数的积分表示
\begin{equation*}
	K_\lambda(z)=\frac{1}{2}\int_{-\infty}^{\infty}
	e^{-z\cosh w+\lambda w}\d w,
	\qquad \Re z>0,
\end{equation*}
得到
\begin{equation*}
	\mathcal J_n^{(\sigma)}
	=2e^{-\sigma\pi\nu/2}e^{in\pi/2}
	K_{n+i\sigma\nu}(\nu\varepsilon).
\end{equation*}
特别地
\begin{equation*}
	\mathcal J_0^{(\sigma)}
	=2e^{-\sigma\pi\nu/2}K_{i\nu}(\nu\varepsilon),
\end{equation*}
因为 $K_{-i\nu}=K_{i\nu}$。以下简记
\begin{equation*}
	K=K_{i\nu}(\nu\varepsilon),
	\qquad
	K'=\left.\frac{\partial K_{i\nu}(x)}{\partial x}\right|_{x=\nu\varepsilon}.
\end{equation*}
整数平移阶数可由递推式
\begin{equation*}
	K_{\lambda+1}(x)=\frac{\lambda}{x}K_\lambda(x)-K_\lambda'(x),
	\qquad
	K_{\lambda-1}(x)=-\frac{\lambda}{x}K_\lambda(x)-K_\lambda'(x)
\end{equation*}
逐级约化，而二阶导数可由虚阶 Bessel 方程
\begin{equation*}
	x^2K_{i\nu}''(x)+xK_{i\nu}'(x)-(x^2-\nu^2)K_{i\nu}(x)=0
\end{equation*}
消去。因此有限次多项式乘以双曲轨道相位的积分最终都可写成 $K_{i\nu}$ 与 $K_{i\nu}'$ 的组合。

\textbf{电偶极与电四极连续谱。}

先计算电偶极项。由 $\vb p=q_D\vb r$，只需求相对加速度的傅里叶变换。从轨道参数式得到
\begin{equation*}
	\dot x=\sigma v\frac{\sinh u}{\varepsilon\cosh u+\sigma},
	\qquad
	\dot y=vs\frac{\cosh u}{\varepsilon\cosh u+\sigma}.
\end{equation*}
又有
\begin{equation*}
	\frac{\d}{\d u}e^{i\omega t(u)}
	=i\nu\left(\varepsilon\cosh u+\sigma\right)e^{i\omega t(u)}.
\end{equation*}
于是，在加入收敛因子后分部积分，
\begin{equation*}
	\begin{aligned}
	\widetilde a_x
	&=\int_{-\infty}^{\infty}\frac{\d\dot x}{\d u}e^{i\omega t}\d u
	=-i\sigma v\nu\int_{-\infty}^{\infty}\sinh u\,e^{i\omega t}\d u,\\
	\widetilde a_y
	&=\int_{-\infty}^{\infty}\frac{\d\dot y}{\d u}e^{i\omega t}\d u
	=-ivs\nu\int_{-\infty}^{\infty}\cosh u\,e^{i\omega t}\d u.
	\end{aligned}
\end{equation*}
第一项由
\begin{equation*}
	\frac{\partial\mathcal J_0^{(\sigma)}}{\partial\varepsilon}
	=i\nu\int_{-\infty}^{\infty}\sinh u\,e^{i\omega t}\d u
\end{equation*}
计算。第二项利用全导数积分
\begin{equation*}
	0=\int_{-\infty}^{\infty}\frac{\d}{\d u}e^{i\omega t}\d u
	=i\nu\int_{-\infty}^{\infty}
	\left(\varepsilon\cosh u+\sigma\right)e^{i\omega t}\d u
\end{equation*}
化为 $\mathcal J_0^{(\sigma)}$。代入母积分结果后得到
\begin{equation*}
	\widetilde a_x=-2\sigma v\nu e^{-\sigma\pi\nu/2}K',
	\qquad
	\widetilde a_y=2i\sigma v\nu\frac{s}{\varepsilon}e^{-\sigma\pi\nu/2}K.
\end{equation*}
因此
\begin{equation*}
	\widetilde{\ddot p}_i\widetilde{\ddot p}_i^*
	=4q_D^2v^2\nu^2e^{-\sigma\pi\nu}
	\left(K'^2+\frac{s^2}{\varepsilon^2}K^2\right),
\end{equation*}
代入一般电偶极频域频谱公式，得到任意质量比、任意碰撞参数的电偶极连续谱闭式
\begin{equation*}
	\frac{\d E_{E1}^{(\sigma)}(b)}{\d\omega}
	=\frac{2q_D^2v^2}{3\pi^2\varepsilon_0c^3}
	\nu^2e^{-\sigma\pi\nu}
	\left[
	K'^2+\frac{s^2}{\varepsilon^2}K^2
	\right].
\end{equation*}
吸引与排斥的差别全部包含在 $e^{-\sigma\pi\nu}$ 中；其余函数只依赖于 $\nu$ 与 $s=b/b_0$。

电四极项需要相对坐标二次型的傅里叶变换。由有效四极电荷，只需计算三个独立的二次轨道矩
\begin{equation*}
	R_{xx}=\int_{-\infty}^{\infty}x^2e^{i\omega t}\d t,
	\qquad
	R_{yy}=\int_{-\infty}^{\infty}y^2e^{i\omega t}\d t,
	\qquad
	R_{xy}=\int_{-\infty}^{\infty}xye^{i\omega t}\d t.
\end{equation*}
由双曲空间参数式及 $\d t/\d u=(\varepsilon\cosh u+\sigma)/\omega_0$，先写成
\begin{equation*}
	\begin{aligned}
	R_{xx}&=\frac{b_0^3}{v}\int_{-\infty}^{\infty}
	(\varepsilon+\sigma\cosh u)^2(\varepsilon\cosh u+\sigma)e^{i\omega t}\d u,\\
	R_{yy}&=\frac{b_0^3s^2}{v}\int_{-\infty}^{\infty}
	\sinh^2u(\varepsilon\cosh u+\sigma)e^{i\omega t}\d u,\\
	R_{xy}&=\frac{b_0^3s}{v}\int_{-\infty}^{\infty}
	\sinh u(\varepsilon+\sigma\cosh u)(\varepsilon\cosh u+\sigma)e^{i\omega t}\d u.
	\end{aligned}
\end{equation*}
为把每一项直接归入母积分，定义
\begin{equation*}
	\mathcal C_m=\frac{\mathcal J_m^{(\sigma)}+\mathcal J_{-m}^{(\sigma)}}{2},
	\qquad
	\mathcal S_m=\frac{\mathcal J_m^{(\sigma)}-\mathcal J_{-m}^{(\sigma)}}{2}.
\end{equation*}
其中 $\mathcal C_m$、$\mathcal S_m$ 分别是 $\cosh(mu)$、$\sinh(mu)$ 与轨道相位乘积的积分。利用
\begin{equation*}
	\cosh^2u=\frac{1+\cosh2u}{2},
	\quad
	\cosh^3u=\frac{3\cosh u+\cosh3u}{4},
	\quad
	\sinh u\cosh u=\frac{\sinh2u}{2},
\end{equation*}
\begin{equation*}
	\sinh u(\cosh^2u+1)=\frac{\sinh3u+5\sinh u}{4},
\end{equation*}
二次轨道矩积分逐项展开为
\begin{equation*}
	\begin{aligned}
	R_{xx}=\frac{b_0^3}{v}\bigg[&
	\frac{\varepsilon}{4}\mathcal C_3
	+\varepsilon\left(\varepsilon^2+\frac{11}{4}\right)\mathcal C_1
	+\frac{\sigma(2\varepsilon^2+1)}{2}\mathcal C_2
	+\frac{\sigma(4\varepsilon^2+1)}{2}\mathcal C_0\bigg],\\
	R_{yy}=\frac{b_0^3s^2}{v}\bigg[&
	\frac{\varepsilon}{4}\mathcal C_3
	-\frac{\varepsilon}{4}\mathcal C_1
	+\frac{\sigma}{2}\mathcal C_2
	-\frac{\sigma}{2}\mathcal C_0\bigg],\\
	R_{xy}=\frac{b_0^3s}{v}\bigg[&
	\frac{\varepsilon^2+1}{2}\mathcal S_2
	+\frac{\sigma\varepsilon}{4}
	\left(\mathcal S_3+5\mathcal S_1\right)\bigg].
	\end{aligned}
\end{equation*}
母积分结果将 $\mathcal C_m,\mathcal S_m$ 写成 $K_{i\sigma\nu\pm m}(\nu\varepsilon)$；再反复使用递推式，所有整数平移阶数消去。定义
\begin{equation*}
	\mathcal A_\sigma=\frac{4b_0^3}{v\nu^2}e^{-\sigma\pi\nu/2},
\end{equation*}
可得
\begin{equation*}
	\begin{aligned}
	R_{xx}&=\sigma\mathcal A_\sigma
	\left(\frac{\nu s^2}{\varepsilon}K'+\frac{1}{\varepsilon^2}K\right),\\
	R_{yy}&=-\sigma\mathcal A_\sigma
	\left(\frac{\nu s^2}{\varepsilon}K'-\frac{s^2}{\varepsilon^2}K\right),\\
	R_{xy}&=i\sigma\mathcal A_\sigma s
	\left(\frac{1}{\varepsilon}K'-\frac{\nu s^2}{\varepsilon^2}K\right).
	\end{aligned}
\end{equation*}
轨道位于 $xy$ 平面，故
\begin{equation*}
	\begin{aligned}
	\widetilde D_{xx}&=q_Q(2R_{xx}-R_{yy}),
	&\widetilde D_{yy}&=q_Q(2R_{yy}-R_{xx}),\\
	\widetilde D_{zz}&=-q_Q(R_{xx}+R_{yy}),
	&\widetilde D_{xy}&=\widetilde D_{yx}=3q_QR_{xy}.
	\end{aligned}
\end{equation*}
直接收缩得到
\begin{equation*}
	\widetilde D_{ij}\widetilde D_{ij}^*
	=6q_Q^2\left[
	|R_{xx}|^2+|R_{yy}|^2+3|R_{xy}|^2
	-\operatorname{Re}(R_{xx}R_{yy}^*)
	\right].
\end{equation*}
将二次轨道矩分量代入并整理，可写成
\begin{equation*}
	\widetilde D_{ij}\widetilde D_{ij}^*
	=\frac{96q_Q^2b_0^6}{v^2\nu^4}e^{-\sigma\pi\nu}\mathcal F_Q(\nu,s),
\end{equation*}
其中
\begin{equation*}
	\begin{aligned}
	\mathcal F_Q(\nu,s)=\frac{1}{\varepsilon^4}\bigg[&
	3\varepsilon^2s^2(1+\nu^2s^2)K'^2
	+3\varepsilon\nu s^2(1-3s^2)KK'\\
	&+\left(1-s^2+s^4+3\nu^2s^6\right)K^2
	\bigg].
	\end{aligned}
\end{equation*}
利用 $\omega=\nu v/b_0$，把收缩结果代入一般电四极频域频谱公式，得到电四极连续谱闭式
\begin{equation*}
	\frac{\d E_{E2}^{(\sigma)}(b)}{\d\omega}
	=\frac{2q_Q^2v^4}{15\pi^2\varepsilon_0c^5}
	\nu^2e^{-\sigma\pi\nu}\mathcal F_Q(\nu,s).
\end{equation*}
至此，有限质量二体库仑散射在固定碰撞参数下的 $E1$、$E2$ 连续谱均已化为虚阶第二类修正 Bessel 函数的闭式。
\textbf{总辐射能量与角动量。}

若不分辨频率，则同一结果可由 Parseval 恒等式写成
\begin{equation*}
	\Delta E_X^{(\sigma)}(b)
	=\int_0^\infty\frac{\d E_X^{(\sigma)}(b)}{\d\omega}\d\omega,
	\qquad X=E1,E2.
\end{equation*}
直接在时域积分可以更简洁地求出这个频率积分，并同时得到辐射角动量。令
\begin{equation*}
	f(\theta)=\varepsilon\cos\theta-\sigma,
	\qquad
	\theta_\infty=\arccos\frac{\sigma}{\varepsilon},
	\qquad
	s=\sqrt{\varepsilon^2-1},
\end{equation*}
则双曲轨道可写成
\begin{equation*}
	r=\frac{p}{f(\theta)},
	\qquad
	\frac{\d t}{\d\theta}
	=\sqrt{\frac{\mu p^3}{\mathcal C}}\frac{1}{f(\theta)^2},
	\qquad
	-\theta_\infty<\theta<\theta_\infty.
\end{equation*}
由 $\vb p=q_D\vb r$ 和统一库仑相对运动方程，
\begin{equation*}
	\ddot{\vb p}
	=\sigma\frac{q_D\mathcal C}{\mu}\frac{\vb r}{r^3},
	\qquad
	\dot{\vb p}\times\ddot{\vb p}
	=-\sigma\frac{q_D^2\mathcal C}{\mu^2r^3}\vb L.
\end{equation*}
代入前文的瞬时 $E1$ 能量与角动量通量，只需计算
\begin{equation*}
	\begin{aligned}
	\int_{-\theta_\infty}^{\theta_\infty}f(\theta)^2\d\theta
	&=(\varepsilon^2+2)\theta_\infty-3\sigma s,\\
	\int_{-\theta_\infty}^{\theta_\infty}f(\theta)\d\theta
	&=2\left(s-\sigma\theta_\infty\right).
	\end{aligned}
\end{equation*}
于是一次双曲散射的电偶极总辐射能量和辐射角动量分别为
\begin{equation*}
	\Delta E_{E1}^{(\sigma)}
	=\frac{q_D^2\mathcal C^{3/2}}
	{6\pi\varepsilon_0c^3\mu^{3/2}p^{5/2}}
	\left[(\varepsilon^2+2)\theta_\infty-3\sigma s\right],
\end{equation*}
\begin{equation*}
	\Delta\vb L_{E1,\mathrm{rad}}^{(\sigma)}
	=\frac{q_D^2\mathcal C}
	{3\pi\varepsilon_0c^3\mu p}
	\left(\theta_\infty-\sigma s\right)\uv L.
\end{equation*}
对于吸引散射，$\theta_\infty=\pi-\arccos(1/\varepsilon)$；对于排斥散射，$\theta_\infty=\arccos(1/\varepsilon)$。因此吸引与排斥两类双曲散射的偶极损失，都由上面两式给出。

四极总损失同样可从固定碰撞参数频谱积分得到，也可在时域直接完成。由
\begin{equation*}
	\begin{aligned}
	\dddot D_{ij}\dddot D_{ij}
	&=\frac{24q_Q^2\mathcal C^2}{\mu^2r^6}
	\left[12r^2\dot{\vb r}^{\,2}
	-11(\vb r\cdot\dot{\vb r})^2\right],\\
	\ddot{\vb D}\crossdot\dddot{\vb D}
	&=\frac{36\sigma\mathcal C q_Q^2}{\mu}
	(\vb r\times\dot{\vb r})
	\left[
	\frac{2\sigma\mathcal C}{\mu r^4}
	+\frac{3(\vb r\cdot\dot{\vb r})^2}{r^5}
	-\frac{2\dot{\vb r}^{\,2}}{r^3}
	\right],
	\end{aligned}
\end{equation*}
以及
\begin{equation*}
	\dot r=\frac{L\varepsilon\sin\theta}{\mu p},
	\qquad
	r\dot\theta=\frac{Lf(\theta)}{\mu p},
	\qquad
	L=\sqrt{\mu\mathcal C p},
\end{equation*}
四极能量积分归结为
\begin{equation*}
	\int_{-\theta_\infty}^{\theta_\infty}
	\left[12f^4+\varepsilon^2\sin^2\theta\,f^2\right]\d\theta
	=24\left[
	\left(1+\frac{73}{24}\varepsilon^2+\frac{37}{96}\varepsilon^4\right)\theta_\infty
	-\sigma s\frac{602+673\varepsilon^2}{288}
	\right],
\end{equation*}
而角动量积分为
\begin{equation*}
	\int_{-\theta_\infty}^{\theta_\infty}
	f\left[\varepsilon^2\sin^2\theta+2\sigma f-2f^2\right]\d\theta
	=\sigma(8+7\varepsilon^2)\theta_\infty
	-(13+2\varepsilon^2)s.
\end{equation*}
因此
\begin{equation*}
	\Delta E_{E2}^{(\sigma)}
	=\frac{4q_Q^2\mathcal C^{5/2}}
	{5\pi\varepsilon_0c^5\mu^{5/2}p^{7/2}}
	\left[
	\left(1+\frac{73}{24}\varepsilon^2+\frac{37}{96}\varepsilon^4\right)\theta_\infty
	-\sigma s\frac{602+673\varepsilon^2}{288}
	\right],
\end{equation*}
\begin{equation*}
	\Delta\vb L_{E2,\mathrm{rad}}^{(\sigma)}
	=\frac{q_Q^2\mathcal C^2}
	{10\pi\varepsilon_0c^5\mu^2p^2}
	\left[
	(8+7\varepsilon^2)\theta_\infty
	-\sigma(13+2\varepsilon^2)s
	\right]\uv L.
\end{equation*}
这些公式在 $q_D=0$ 时给出领先阶辐射；特别地，同号且荷质比相等的双曲散射正是上面两式取 $\sigma=+1$ 的情形。

\textbf{谱辐射截面。}

固定碰撞参数的谱给出单条经典轨道的辐射。按照能量加权的一般谱辐射截面定义，轴对称库仑散射的谱辐射截面为
\begin{equation*}
	\frac{\d\chi_X^{(\sigma)}}{\d\omega}
	=2\pi\int_0^\infty b\d b\,
	\frac{\d E_X^{(\sigma)}(b)}{\d\omega},
	\qquad X=E1,E2.
\end{equation*}
由于
\begin{equation*}
	b\d b=b_0^2\varepsilon\d\varepsilon,
	\qquad
	x=\nu\varepsilon,
\end{equation*}
偶极截面成为
\begin{equation*}
	\frac{\d\chi_{E1}^{(\sigma)}}{\d\omega}
	=\frac{4q_D^2v^2b_0^2}{3\pi\varepsilon_0c^3}e^{-\sigma\pi\nu}
	\int_\nu^\infty x\left[
	K_{i\nu}'(x)^2+
	\left(1-\frac{\nu^2}{x^2}\right)K_{i\nu}(x)^2
	\right]\d x.
\end{equation*}
由虚阶 Bessel 方程，
\begin{equation*}
	\frac{\d}{\d x}\left[xK_{i\nu}(x)K_{i\nu}'(x)\right]
	=x\left[
	K_{i\nu}'(x)^2+
	\left(1-\frac{\nu^2}{x^2}\right)K_{i\nu}(x)^2
	\right].
\end{equation*}
$K_{i\nu}(x)$ 在 $x\to\infty$ 指数衰减，故
\begin{equation*}
	\frac{\d\chi_{E1}^{(\sigma)}}{\d\omega}
	=\frac{4q_D^2v^2b_0^2}{3\pi\varepsilon_0c^3}
	e^{-\sigma\pi\nu}
	\left[-\nu K_{i\nu}(\nu)K_{i\nu}'(\nu)\right].
\end{equation*}
四极截面同样可以完全积出。把 $\mathcal F_Q$ 的定义中
\begin{equation*}
	\varepsilon=\frac{x}{\nu},
	\qquad
	s^2=\frac{x^2-\nu^2}{\nu^2}
\end{equation*}
代入，记 $K=K_{i\nu}(x)$，则可由虚阶 Bessel 方程逐项验证
\begin{equation*}
	x\mathcal F_Q
	=\frac{\d}{\d x}
	\left[A(x)K^2+B(x)KK'+C(x)K'^2\right],
\end{equation*}
其中
\begin{equation*}
	\begin{aligned}
	A(x)&=-\frac{(x^2-\nu^2)(2x^2-3\nu^2)}{2x^2},\\
	B(x)&=3\frac{(x^2-\nu^2)^2+x^2}{x},\\
	C(x)&=\frac{3\nu^2-13x^2}{2}.
	\end{aligned}
\end{equation*}
在 $x=\nu$ 处
\begin{equation*}
	A(\nu)=0,
	\qquad
	B(\nu)=3\nu,
	\qquad
	C(\nu)=-5\nu^2,
\end{equation*}
而上边界仍为零，故
\begin{equation*}
	\int_\nu^\infty x\mathcal F_Q\d x
	=5\nu^2K_{i\nu}'(\nu)^2
	-3\nu K_{i\nu}(\nu)K_{i\nu}'(\nu).
\end{equation*}
最终得到四极谱辐射截面的闭式
\begin{equation*}
	\frac{\d\chi_{E2}^{(\sigma)}}{\d\omega}
	=\frac{4q_Q^2v^4b_0^2}{15\pi\varepsilon_0c^5}
	e^{-\sigma\pi\nu}
	\left[
	5\nu^2K_{i\nu}'(\nu)^2
	-3\nu K_{i\nu}(\nu)K_{i\nu}'(\nu)
	\right].
\end{equation*}
若需要光子数截面，只需除以单光子能量：
\begin{equation*}
	\frac{\d\sigma_{\gamma,X}}{\d\omega}
	=\frac{1}{\hbar\omega}\frac{\d\chi_X}{\d\omega}.
\end{equation*}
\textbf{软频与硬频渐近。}

精确闭式建立后，其软频和硬频结构都可直接由虚阶 Bessel 函数的渐近式读出。

固定 $s$ 而令 $\nu\to0$ 时，
\begin{equation*}
	K_{i\nu}(\nu\varepsilon)
	=\ln\frac{2}{\nu\varepsilon}-\gamma_{\mathrm E}+o(1),
	\qquad
	K_{i\nu}'(\nu\varepsilon)
	=-\frac{1}{\nu\varepsilon}+o(\nu^{-1}).
\end{equation*}
代入固定碰撞参数的闭式，得到
\begin{equation*}
	\frac{\d E_{E1}}{\d\omega}
	\longrightarrow
	\frac{2q_D^2v^2}{3\pi^2\varepsilon_0c^3}
	\frac{1}{1+s^2}
	=\frac{q_D^2}{6\pi^2\varepsilon_0c^3}|\Delta\vb v|^2,
\end{equation*}
以及
\begin{equation*}
	\frac{\d E_{E2}}{\d\omega}
	\longrightarrow
	\frac{2q_Q^2v^4}{5\pi^2\varepsilon_0c^5}
	\frac{s^2}{(1+s^2)^2}
	=\frac{q_Q^2v^4}{10\pi^2\varepsilon_0c^5}\sin^2\Theta.
\end{equation*}
偶极软谱只取决于初末速度差。四极软谱在严格正碰时为零，这是因为初末速度虽反向，但张量 $v_iv_j$ 不变。

令
\begin{equation*}
	L_\nu=\ln\frac{2}{\nu}-\gamma_{\mathrm E},
\end{equation*}
由截面闭式可得
\begin{equation*}
	\frac{\d\chi_{E1}^{(\sigma)}}{\d\omega}
	=\frac{4q_D^2v^2b_0^2}{3\pi\varepsilon_0c^3}
	\left[L_\nu+o(1)\right],
\end{equation*}
\begin{equation*}
	\frac{\d\chi_{E2}^{(\sigma)}}{\d\omega}
	=\frac{4q_Q^2v^4b_0^2}{15\pi\varepsilon_0c^5}
	\left[5+3L_\nu+o(1)\right].
\end{equation*}
软频最低阶与 $\sigma$ 无关；对数来自越来越大的碰撞参数，实际体系中的屏蔽长度或有限观测时间将给出红外截断。

固定 $s>0$ 而令 $\nu\to\infty$，定义
\begin{equation*}
	\Delta(s)=s-\arctan s>0.
\end{equation*}
虚阶 Bessel 函数的陡降区渐近为
\begin{equation*}
	K_{i\nu}(\nu\varepsilon)
	\sim\sqrt{\frac{\pi}{2\nu s}}
	\exp\left[-\frac{\pi\nu}{2}-\nu\Delta(s)\right],
	\qquad
	K_{i\nu}'(\nu\varepsilon)
	\sim-\frac{s}{\varepsilon}K_{i\nu}(\nu\varepsilon).
\end{equation*}
故
\begin{equation*}
	\frac{\d E_{E1}^{(\sigma)}}{\d\omega}
	\sim
	\frac{2q_D^2v^2}{3\pi\varepsilon_0c^3}
	\frac{\nu s}{\varepsilon^2}
	\exp\left[-2\nu\Delta(s)-(1+\sigma)\pi\nu\right],
\end{equation*}
\begin{equation*}
	\frac{\d E_{E2}^{(\sigma)}}{\d\omega}
	\sim
	\frac{2q_Q^2v^4}{5\pi\varepsilon_0c^5}
	\frac{\nu^3s^5}{\varepsilon^4}
	\exp\left[-2\nu\Delta(s)-(1+\sigma)\pi\nu\right].
\end{equation*}
对吸引势，$s\ll1$ 时 $\Delta(s)=s^3/3+O(s^5)$，故固定碰撞参数谱的特征截止满足 $\nu s^3\sim1$。上述固定 $s$ 渐近对 $s\to0$ 不一致；谱截面的高频极限必须使用转折点 $x=\nu$ 附近的一致渐近式
\begin{equation*}
	e^{\pi\nu/2}K_{i\nu}(\nu)
	\sim\frac{\sqrt3\,\Gamma(1/3)}{6^{2/3}\nu^{1/3}},
	\qquad
	-e^{\pi\nu/2}K_{i\nu}'(\nu)
	\sim\frac{\sqrt3\,\Gamma(2/3)}{6^{1/3}\nu^{2/3}}.
\end{equation*}
代入截面闭式，得到
\begin{equation*}
	\frac{\d\chi_{E1}^{(\sigma)}}{\d\omega}
	\sim
	\frac{4q_D^2v^2b_0^2}{3\sqrt3\,\varepsilon_0c^3}
	e^{-(1+\sigma)\pi\nu},
\end{equation*}
以及
\begin{equation*}
	\frac{\d\chi_{E2}^{(\sigma)}}{\d\omega}
	\sim
	\frac{8\sqrt\pi}{3^{7/6}\Gamma(1/6)}
	\frac{q_Q^2v^4b_0^2}{\varepsilon_0c^5}
	\nu^{2/3}e^{-(1+\sigma)\pi\nu}.
\end{equation*}
因此排斥势具有额外的 $e^{-2\pi\nu}$ 截止；理想点吸引势的偶极截面趋于常数，而四极截面按 $\nu^{2/3}$ 增长。后两者来自越来越小的碰撞参数，只是经典点势模型的中间渐近，最终会被有限尺寸、相对论、量子反冲或辐射反作用截断。

\textbf{相对论固定中心库仑散射的软谱。}

前面的双曲轨道闭式建立在非相对论二体动力学上。若只要求软频极限，则可以放宽速度条件而不必求完整的相对论 Fourier 积分，因为一般相对论软谱公式只需要初末速度。为给出一个完全解析的例子，考虑质量为 $m$、电荷为 $q$ 的试探粒子在固定点电荷产生的外部库仑势中运动：
\begin{equation*}
	V_\sigma(r)=\sigma\frac{\kappa}{r},
	\qquad
	\kappa=\frac{|qQ|}{4\pi\varepsilon_0},
	\qquad
	\sigma=\operatorname{sgn}(qQ).
\end{equation*}
这里固定中心不辐射，全部软辐射来自试探粒子的速度改变。应当强调，这一模型是外部静电场中的相对论单粒子问题；若两个有限质量电荷都相对论性，完整 Maxwell--Lorentz 动力学不能简单化为一个瞬时的约化质量中心势。

试探粒子的守恒能量和角动量为
\begin{equation*}
	\mathcal E
	=\sqrt{m^2c^4+c^2\left(p_r^2+\frac{L^2}{r^2}\right)}
	+\sigma\frac{\kappa}{r},
	\qquad
	L=p_\infty b,
\end{equation*}
其中
\begin{equation*}
	\mathcal E=\gamma_\infty mc^2,
	\qquad
	p_\infty=\gamma_\infty mv_\infty,
	\qquad
	\beta_\infty=\frac{v_\infty}{c}.
\end{equation*}
令 $u=1/r$。由 Hamilton 方程可得 $p_r=-L\d u/\d\varphi$，于是径向能量方程化为
\begin{equation*}
	\left(\frac{\d u}{\d\varphi}\right)^2
	=\frac{p_\infty^2}{L^2}
	-\frac{2\sigma\mathcal E\kappa}{L^2c^2}u
	-\left(1-\frac{\kappa^2}{L^2c^2}\right)u^2.
\end{equation*}
对 $\varphi$ 求导，得到线性轨道方程
\begin{equation*}
	\frac{\d^2u}{\d\varphi^2}+\lambda^2u
	=-\frac{\sigma\mathcal E\kappa}{L^2c^2},
	\qquad
	\lambda=\sqrt{1-a^2},
	\qquad
	a=\frac{\kappa}{Lc}.
\end{equation*}
先考虑 $a<1$ 的散射支。把最近点取在 $\varphi=0$，解为
\begin{equation*}
	u(\varphi)=A\cos(\lambda\varphi)-\sigma u_c,
	\qquad
	u_c=\frac{\mathcal E\kappa}{L^2c^2-\kappa^2},
\end{equation*}
其中由径向能量方程的常数项得到
\begin{equation*}
	A^2=u_c^2+\frac{p_\infty^2}{L^2\lambda^2}.
\end{equation*}
令
\begin{equation*}
	\delta
	=\arctan\frac{a}{\beta_\infty\lambda}.
\end{equation*}
利用 $\mathcal E/(p_\infty c)=1/\beta_\infty$，可验证
\begin{equation*}
	\frac{u_c}{A}
	=\frac{a/(\beta_\infty\lambda)}
	{\sqrt{1+a^2/(\beta_\infty^2\lambda^2)}}
	=\sin\delta.
\end{equation*}
轨道的两条渐近线满足 $u(\pm\varphi_\infty)=0$，因而
\begin{equation*}
	\lambda\varphi_\infty=\frac{\pi}{2}-\sigma\delta.
\end{equation*}
从入射无穷远到出射无穷远所扫过的极角为
\begin{equation*}
	\Delta\varphi_\sigma
	=2\varphi_\infty
	=\frac{\pi-2\sigma\delta}{\lambda},
\end{equation*}
相对于无相互作用直线的有符号偏转角为
\begin{equation*}
	\chi_\sigma
	=\pi-\Delta\varphi_\sigma
	=\pi-\frac{\pi-2\sigma\arctan[a/(\beta_\infty\lambda)]}{\lambda}.
\end{equation*}
这就是相对论固定中心库仑散射的 Darwin 角。当轨道发生绕转时，$\chi_\sigma$ 可以超出通常的主值区间；软谱只通过初末速度的内积依赖它，因此应保留 $\cos\chi_\sigma$，而不必人为把角度折回 $[0,\pi]$。吸引势在 $a\geq1$ 时失去离心势垒并出现落向中心的轨道，上面的偏转角公式所描述的是 $a<1$ 的普通散射支。

初末速率相同，均为 $\beta_\infty c$，故两条渐近外线之间的相对 Lorentz 因子为
\begin{equation*}
	\Gamma_\sigma(b)
	=\gamma_\infty^2
	\left(1-\beta_\infty^2\cos\chi_\sigma(b)\right)
	=1+2\gamma_\infty^2\beta_\infty^2
	\sin^2\frac{\chi_\sigma(b)}{2}.
\end{equation*}
于是无需求解相对论轨道的时间参数或完整 Lienard--Wiechert Fourier 积分，便可由单电荷相对论软谱公式立即得到固定碰撞参数的解析软谱
\begin{equation*}
	\left.\frac{\d E_{\mathrm{soft}}^{(\sigma)}(b)}{\d\omega}\right|_{\omega\to0}
	=\frac{q^2}{2\pi^2\varepsilon_0c}
	\left[
	\frac{\Gamma_\sigma(b)\operatorname{arccosh}\Gamma_\sigma(b)}
	{\sqrt{\Gamma_\sigma(b)^2-1}}-1
	\right].
\end{equation*}
相应的角分布同样是闭式：
\begin{equation*}
	\left.\frac{\d^2E_{\mathrm{soft}}^{(\sigma)}}{\d\omega\d\Omega}\right|_{\omega\to0}
	=\frac{q^2}{16\pi^3\varepsilon_0c}
	\left|
	\vb n\times\left\{\vb n\times
	\left[
	\frac{\boldsymbol\beta_f}{1-\vb n\cdot\boldsymbol\beta_f}
	-\frac{\boldsymbol\beta_i}{1-\vb n\cdot\boldsymbol\beta_i}
	\right]\right\}
	\right|^2,
\end{equation*}
其中 $|\boldsymbol\beta_i|=|\boldsymbol\beta_f|=\beta_\infty$ 且 $\boldsymbol\beta_i\cdot\boldsymbol\beta_f=\beta_\infty^2\cos\chi_\sigma$。

几个极限可用来检验这一闭式。首先，当 $a\ll1$ 时，
\begin{equation*}
	\lambda=1+O(a^2),
	\qquad
	\chi_\sigma
	=\frac{2\sigma\kappa}{\gamma_\infty m v_\infty^2b}+O(\kappa^2).
\end{equation*}
若再取 $\beta_\infty\ll1$，软谱闭式退化为
\begin{equation*}
	\left.\frac{\d E_{\mathrm{soft}}}{\d\omega}\right|_{\omega\to0}
	=\frac{2q^2v_\infty^2}{3\pi^2\varepsilon_0c^3}
	\sin^2\frac{\chi_\sigma}{2}
	+O(c^{-5}),
\end{equation*}
并由非相对论 Rutherford 关系 $\sin^2(\chi/2)=1/(1+b^2/b_0^2)$ 恢复固定碰撞参数偶极软谱的固定中心极限。

其次，在大碰撞参数端，把小偏转角展开与小相对速度展开结合给出
\begin{equation*}
	\left.\frac{\d E_{\mathrm{soft}}(b)}{\d\omega}\right|_{b\to\infty}
	\simeq
	\frac{2q^2\kappa^2}{3\pi^2\varepsilon_0m^2v_\infty^2c^3}\frac{1}{b^2}.
\end{equation*}
因此未屏蔽库仑场的谱辐射截面仍有 Coulomb 对数：
\begin{equation*}
	\left.\frac{\d\chi_{\mathrm{soft}}}{\d\omega}\right|_{\omega\to0}
	\simeq
	\frac{4q^2\kappa^2}{3\pi\varepsilon_0m^2v_\infty^2c^3}
	\ln\frac{b_{\max}}{b_{\min}}.
\end{equation*}
相对论因子在这个大 $b$ 的领先对数系数中抵消；速度和强场修正进入有限项以及小碰撞参数区域。最后，当 $\gamma_\infty\sin(|\chi_\sigma|/2)\gg1$ 时，软谱闭式变为
\begin{equation*}
	\left.\frac{\d E_{\mathrm{soft}}}{\d\omega}\right|_{\omega\to0}
	\simeq
	\frac{q^2}{2\pi^2\varepsilon_0c}
	\left[
	2\ln\left(2\gamma_\infty\left|\sin\frac{\chi_\sigma}{2}\right|\right)-1
	\right],
\end{equation*}
显示了相对论软轫致辐射的对数增强和前向束射。

软谱闭式的精确性具有明确含义：它对 $\beta_\infty$ 和偏转角不作展开，对 $a<1$ 的固定中心相对论库仑散射轨道使用精确的 Darwin 角，但只给出 $\omega\to0$ 的领先项。其软条件是辐射周期远长于有效散射时间，即 $\omega\tau_{\mathrm{sc}}\ll1$；它不要求 $\beta_\infty\ll1$，也不要求电偶极长波条件 $\omega a_{\mathrm{eff}}/c\ll1$。若固定中心近似被放宽而靶粒子也发生反冲，一般的相对论软谱积分公式仍然成立，只需把两粒子的入射、出射四条外线全部代入；困难仅转移到相对论二体散射角和渐近四动量的动力学求解上。
\textbf{偶极抵消的特殊情形。}

有限质量二体问题还有一个重要的特殊情形，即电偶极矩由于荷质比关系而抵消。

当两粒子的荷质比相同，
\begin{equation*}
	\frac{q_1}{M_1}=\frac{q_2}{M_2},
\end{equation*}
有效偶极电荷公式给出 $q_D=0$，整个电偶极谱消失，而固定碰撞参数的电四极谱与其截面闭式成为领先结果。若 $q_D\neq0$，由软频结果得到
\begin{equation*}
	\frac{\d E_{E2}/\d\omega}{\d E_{E1}/\d\omega}
	\longrightarrow
	\frac{3}{5}\left(\frac{q_Q}{q_D}\right)^2
	\frac{v^2}{c^2}\frac{s^2}{1+s^2}.
\end{equation*}
因此在没有电荷质量组合抵消时，四极通常比偶极低相对阶 $v^2/c^2$；当 $|q_D|\lesssim(v/c)|q_Q|$ 时，二者才可能在同一量级竞争。相反，若 $M_1=M_2$ 且 $q_1=-q_2$，则 $q_Q=0$ 而 $q_D=q_1$，最低阶仍为纯电偶极辐射。

\textbf{椭圆束缚运动的谱线与绝热演化。}

负相对能量只可能出现在吸引势 $\sigma=-1$，此时轨道为周期性椭圆。按照周期运动的 Fourier 级数约定与一般周期谱线功率公式，连续 Fourier 变换应改为离散 Fourier 级数，辐射由整数谐波线组成。设椭圆半长轴为 $a$，离心率为 $0\leq e<1$，平均角频率为
\begin{equation*}
	\Omega=\sqrt{\frac{\mathcal C}{\mu a^3}}.
\end{equation*}
以偏近点角 $u$ 参数化
\begin{equation*}
	x=a(\cos u-e),
	\qquad
	y=a\sqrt{1-e^2}\sin u,
	\qquad
	\Omega t=u-e\sin u.
\end{equation*}
按照周期 Fourier 级数约定，需要计算
\begin{equation*}
	\mathcal I_{N,m}(e)
	=\frac{1}{2\pi}\int_0^{2\pi}
	e^{imu}(1-e\cos u)e^{iN(u-e\sin u)}\d u.
\end{equation*}
利用 Jacobi--Anger 展开
\begin{equation*}
	e^{-iNe\sin u}=\sum_{k=-\infty}^{\infty}J_k(-Ne)e^{iku}
\end{equation*}
并逐项积分，得到
\begin{equation*}
	\mathcal I_{N,m}(e)
	=J_{N+m}(Ne)-\frac{e}{2}
	\left[J_{N+m-1}(Ne)+J_{N+m+1}(Ne)\right].
\end{equation*}
对 $N\geq1$ 有 $\mathcal I_{N,0}=0$。定义
\begin{equation*}
	C_{N,m}=\frac{\mathcal I_{N,m}+\mathcal I_{N,-m}}{2},
	\qquad
	S_{N,m}=\frac{\mathcal I_{N,m}-\mathcal I_{N,-m}}{2i}.
\end{equation*}
坐标的 Fourier 系数为
\begin{equation*}
	r_{x,N}=\frac{a}{N}J_N'(Ne),
	\qquad
	r_{y,N}=i\frac{a\sqrt{1-e^2}}{Ne}J_N(Ne).
\end{equation*}
代入一般周期谱线功率公式，得到第 $N$ 条电偶极谱线
\begin{equation*}
	P_{E1,N}
	=\frac{q_D^2a^2\Omega^4}{3\pi\varepsilon_0c^3}N^2
	\left[J_N'(Ne)^2+\frac{1-e^2}{e^2}J_N(Ne)^2\right],
	\qquad N\geq1.
\end{equation*}
四极所需的二次轨道矩系数为
\begin{equation*}
	\begin{aligned}
	R_{xx,N}&=a^2\left(\frac{1}{2}C_{N,2}-2eC_{N,1}\right),\\
	R_{yy,N}&=-\frac{a^2(1-e^2)}{2}C_{N,2},\\
	R_{xy,N}&=a^2\sqrt{1-e^2}
	\left(\frac{1}{2}S_{N,2}-eS_{N,1}\right).
	\end{aligned}
\end{equation*}
相应的 $D_{ij,N}$ 仍由双曲情形中四极分量与二次轨道矩的关系给出，因此
\begin{equation*}
	P_{E2,N}
	=\frac{q_Q^2(N\Omega)^6}{60\pi\varepsilon_0c^5}
	\left[
	|R_{xx,N}|^2+|R_{yy,N}|^2+3|R_{xy,N}|^2
	-\operatorname{Re}(R_{xx,N}R_{yy,N}^*)
	\right].
\end{equation*}
总平均谱写成
\begin{equation*}
	\frac{\d\overline P}{\d\omega}
	=\sum_{N=1}^{\infty}
	\left(P_{E1,N}+P_{E2,N}+\cdots\right)
	\delta(\omega-N\Omega).
\end{equation*}
对所有谐波求和必须恢复时域中的周期平均功率。以电偶极项为例，令
\begin{equation*}
	p=a(1-e^2),
	\qquad
	r=\frac{p}{1+e\cos\theta},
	\qquad
	\frac{\d t}{\d\theta}
	=\sqrt{\frac{\mu p^3}{\mathcal C}}
	\frac{1}{(1+e\cos\theta)^2},
	\qquad
	T=\frac{2\pi}{\Omega},
\end{equation*}
则由瞬时多极通量或等价地由谱线 Parseval 求和得到
\begin{equation*}
	-\frac{\d E_{\mathrm{orb}}}{\d t}
	=\sum_{N=1}^{\infty}P_{E1,N}
	=\frac{q_D^2\mathcal C^2}
	{6\pi\varepsilon_0c^3\mu^2p^4}
	\left(1+\frac{e^2}{2}\right)(1-e^2)^{3/2},
\end{equation*}
以及
\begin{equation*}
	-\frac{\d\vb L}{\d t}
	=\frac{q_D^2\mathcal C^{3/2}}
	{6\pi\varepsilon_0c^3\mu^{3/2}p^{5/2}}
	(1-e^2)^{3/2}\uv L.
\end{equation*}
这里 $E_{\mathrm{orb}}=-\mathcal C(1-e^2)/(2p)$ 且 $L=\sqrt{\mu\mathcal C p}$。由上面两式消去 $E_{\mathrm{orb}}$ 与 $L$，得到缓慢辐射反作用下的轨道参数方程
\begin{equation*}
	\frac{\dot p}{p}
	=-\frac{q_D^2\mathcal C}{3\pi\varepsilon_0c^3\mu^2p^3}(1-e^2)^{3/2},
	\qquad
	\frac{\dot e}{e}
	=-\frac{q_D^2\mathcal C}{4\pi\varepsilon_0c^3\mu^2p^3}(1-e^2)^{3/2}.
\end{equation*}
因此
\begin{equation*}
	\frac{e}{e_i}=\left(\frac{p}{p_i}\right)^{3/4}.
\end{equation*}
若记
\begin{equation*}
	\tau=\frac{\pi\varepsilon_0c^3\mu^2p_i^3}{\mathcal Cq_D^2},
	\qquad
	F(e)=\sqrt{1-e^2}+\frac{1}{\sqrt{1-e^2}},
\end{equation*}
则 $e_i\neq0$ 时的积分解可以写成简洁的隐式形式
\begin{equation*}
	t=\frac{4\tau}{e_i^4}\left[F(e_i)-F(e)\right],
	\qquad
	p=p_i\left(\frac{e}{e_i}\right)^{4/3}.
\end{equation*}
在这一绝热模型中，形式上的并合时刻为
\begin{equation*}
	t_{\mathrm c}
	=\frac{4\tau}{e_i^4}\left[F(e_i)-2\right].
\end{equation*}
圆轨道极限 $e_i\to0$ 给出
\begin{equation*}
	e(t)=0,
	\qquad
	p(t)=p_i\left(1-\frac{t}{\tau}\right)^{1/3},
	\qquad
	t_{\mathrm c}=\tau.
\end{equation*}
接近形式并合时，局域速度、有限尺寸和辐射反作用均不再是小修正，因此隐式解只描述缓慢演化的非相对论阶段。

圆轨道 $e=0$ 时，偶极矩只含基频，四极矩只含二次谐波：
\begin{equation*}
	P_{E1,1}=\frac{q_D^2a^2\Omega^4}{6\pi\varepsilon_0c^3},
	\qquad
	P_{E2,2}=\frac{2q_Q^2a^4\Omega^6}{5\pi\varepsilon_0c^5},
\end{equation*}
其余谱线均为零。若再有 $q_D=0$，最低非零辐射线便位于 $2\Omega$。

\textbf{适用范围。}

最后说明上述闭式的适用范围。这些解析式对质量比和库仑偏转角没有作展开，但仍属于非相对论长波多极理论。首先必须要求整个有效轨道上的局域相对速度远小于 $c$。排斥轨道上速度的最大值是无穷远速度，故 $v/c\ll1$ 基本足够；吸引轨道的近心速度为
\begin{equation*}
	v_{\mathrm p}=v\sqrt{\frac{\varepsilon+1}{\varepsilon-1}},
\end{equation*}
因而还需 $v_{\mathrm p}/c\ll1$。其次，多极展开要求对产生给定频率的有效空间尺度有 $\omega r_{\mathrm{eff}}/c\ll1$。排斥硬频区相应为 $\beta\nu\ll1$；吸引硬频截面的主导碰撞满足 $s\sim\nu^{-1/3}$，相应条件为 $\beta\nu^{1/3}\ll1$。

此外还应满足辐射反作用所造成的能量损失远小于入射相对动能，经典轨道角动量 $L/\hbar=\mu vb/\hbar$ 足够大，光子反冲满足 $\hbar\omega\ll\mu v^2/2$，并且主要贡献的最近距离大于粒子或核的有限尺寸、主要贡献的最大距离小于屏蔽尺度。在这些条件共同成立的区域，固定碰撞参数的 $E1$、$E2$ 连续谱闭式与相应的截面闭式给出非束缚二体库仑运动的连续谱；$E1$、$E2$ 谱线功率公式给出吸引束缚轨道的离散谱。

\begin{notes}
作为类比，广义相对论中最低阶引力波辐射由质量四极矩贡献。考虑做二体椭圆运动的质量为 \(M_1\) 和 \(M_2\) 的质点，记
\begin{equation*}
	\alpha = G M_1 M_2, \quad \mu = \frac{M_1 M_2}{M_1 + M_2}.
\end{equation*}

设 2 指向 1 的位矢为 \(\vb{r}\)，则二体运动方程、二体角动量、质心系中体系的质量四极矩为
\begin{equation*}
	\mu \frac{\d^2\vb{r}}{\d t^2} = -\frac{\alpha}{r^3}\vb{r}, \quad \vb{L} = \mu \vb{r} \times \frac{\d \vb{r}}{\d t}, \quad \vb{Q} = \mu \left( 3\vb{r}\vb{r} - r^2 \vb{I} \right).
\end{equation*}

进而，质量四极矩的各阶变化率为
\begin{equation*}
	\begin{aligned}
	\frac{\d \vb{Q}}{\d t} &= \mu \left( 3\frac{\d \vb{r}}{\d t}\vb{r} + 3\vb{r}\frac{\d \vb{r}}{\d t} - 2\left( \vb{r} \cdot \frac{\d \vb{r}}{\d t} \right)\vb{I} \right), \\
	\frac{\d^2\vb{Q}}{\d t^2} &= 2\alpha p \left( -\frac{3\vb{r}\vb{r} - r^2 \vb{I}}{r^3 p} + \frac{\mu^2}{L^2} \left( 3\frac{\d \vb{r}}{\d t}\frac{\d \vb{r}}{\d t} - \left| \frac{\d \vb{r}}{\d t} \right|^2 \vb{I} \right) \right), \\
	\frac{\d^3\vb{Q}}{\d t^3} &= \frac{2\alpha}{r^3} \left( 9\frac{\vb{r}\vb{r}}{r^2} \left( \vb{r} \cdot \frac{\d \vb{r}}{\d t} \right) - 6\left( \frac{\d \vb{r}}{\d t}\vb{r} + \vb{r}\frac{\d \vb{r}}{\d t} \right) + \left( \vb{r} \cdot \frac{\d \vb{r}}{\d t} \right)\vb{I} \right).\\
	\frac{\d^3\vb{Q}}{\d t^3} : \frac{\d^3\vb{Q}}{\d t^3} &= \frac{24\alpha^2}{r^6} \left( -11\left( \vb{r} \cdot \frac{\d \vb{r}}{\d t} \right)^2 + 12 r^2 \left| \frac{\d \vb{r}}{\d t} \right|^2 \right).\\
	\frac{\d^2\vb{Q}}{\d t^2} \crossdot \left( \frac{\d^3\vb{Q}}{\d t^3} \right) &= \frac{36\alpha^2 p \vb{L}}{\mu r^5} \left( \frac{2r}{p} - \frac{3\mu^2}{L^2} \left( \vb{r} \cdot \frac{\d \vb{r}}{\d t} \right)^2 + \frac{2\mu^2}{L^2} r^2 \left| \frac{\d \vb{r}}{\d t} \right|^2 \right).
	\end{aligned}
\end{equation*}

由于观测时更关心周期性椭圆运动的平均效应，因此考虑周期平均后的辐射损失能量、角动量。运用开普勒运动的结论
\begin{equation*}
	r = \frac{p}{1 + e\cos\theta}, \quad \frac{\d \vb{r}}{\d t} = \frac{L}{\mu r} \left( \frac{e\sin\theta}{1 + e\cos\theta}\uv{r} + \uv{\theta} \right), \quad L = \sqrt{\mu\alpha p}, \quad T = \frac{2\pi}{(1 - e^2)^{3/2}} \sqrt{\frac{p^3}{\alpha\mu}}.
\end{equation*}

并代入
\begin{equation*}
	\frac{\d t}{\d \theta} = \sqrt{\frac{\mu p^3}{\alpha}} \frac{1}{(1 + e\cos\theta)^2}.
\end{equation*}

经过繁琐的积分得
\begin{equation*}
	\begin{aligned}
	-\frac{\d E}{\d t} &= \frac{G}{45 c^5 T} \int_{0}^{2\pi} \left( \frac{\d^3\vb{Q}}{\d t^3} : \frac{\d^3\vb{Q}}{\d t^3} \right) \frac{\d t}{\d \theta} \d \theta
	= \frac{32 G}{5 c^5} \frac{\alpha^3}{\mu p^5} (1 - e^2)^{3/2} \left( 1 + \frac{73}{24}e^2 + \frac{37}{96}e^4 \right), \\
	-\frac{\d \vb{L}}{\d t} &= \frac{2 G}{45 c^5 T} \int_{0}^{2\pi} \left( \frac{\d^2\vb{Q}}{\d t^2} \crossdot \left( \frac{\d^3\vb{Q}}{\d t^3} \right) \right) \frac{\d t}{\d \theta} \d \theta
	= \frac{32 G}{5 c^5} \frac{\alpha^{5/2} \uv{L}}{\mu^{1/2} p^{7/2}} (1 - e^2)^{3/2} \left( 1 + \frac{7}{8}e^2 \right).
	\end{aligned}
\end{equation*}

用两质点质量及轨道半长轴 \(a = p(1-e^2)^{-1}\) 表示，即为
\begin{equation*}
	\begin{aligned}
	-\frac{\d E}{\d t} &= \frac{32 G^4 M_1^2 M_2^2 (M_1 + M_2)}{5 c^5 a^5} (1 - e^2)^{-7/2} \left( 1 + \frac{73}{24}e^2 + \frac{37}{96}e^4 \right), \\
	-\frac{\d \vb{L}}{\d t} &= \frac{32 G^{7/2} M_1^2 M_2^2 (M_1 + M_2)^{1/2}}{5 c^5 a^{7/2}} (1 - e^2)^{-2} \left( 1 + \frac{7}{8}e^2 \right).
	\end{aligned}
\end{equation*}
\end{notes}
